\chapter{Rigidity over a base field \(k\): morphisms from a genus-\(0\) curve to a group scheme are constant}
\label{chap:RigidityKbar}
% archon:covers AlgebraicJacobian/RigidityKbar.lean AlgebraicJacobian/Cotangent/ChartAlgebra.lean

% STRATEGY NOTE (iter-157): The project's COMMITTED genus-0 rigidity route is now route
% (c) --- the characteristic-free abelian-variety rigidity-lemma stack developed in its own
% chapter \cref{chap:AbelianVarietyRigidity} (Lean file
% AlgebraicJacobian/AbelianVarietyRigidity.lean, upstream of Jacobian.lean). The headline
% there, thm:rigidity_genus0_curve_to_AV (char-free, [IsAlgClosed k-bar] only), is what
% genusZeroWitness consumes. THIS chapter now documents the over-k-bar declaration
% rigidity_over_kbar (which carries [CharZero k-bar]) together with the FALLBACK route (a):
% the df=0 / cotangent-triviality + Serre-duality decomposition and the closed, axiom-clean
% chart-algebra envelope (which supplies only the converse "df=0 implies constant"). Route
% (a) is retained in tree, off the critical path; the df=0 framing below is the fallback,
% NOT the committed route.

This chapter introduces the project's named declaration for the M2.a sub-step of the genus-\(0\) Albanese-witness argument: rigidity of morphisms from a smooth proper geometrically irreducible curve of genus \(0\) to a smooth proper geometrically irreducible group scheme, \textbf{over an algebraically closed base field \(\bar k\) of characteristic zero} (the alg-closed pivot of \S "Iter-152 alg-closed pivot" below; descent of the constancy conclusion to a general base field \(k\) is a separate downstream step, hosted in the genus-\(0\) witness chapter). The classical statement (Mumford, \emph{Abelian Varieties}, Chapter~II~\S4) is

\begin{quote}
  \emph{Let \(A\) be an abelian variety of positive dimension over an algebraically closed field \(\bar k\). Then every morphism \(f \colon \mathbb P^1_{\bar k} \to A\) is constant.}
\end{quote}

This is the keystone classical input for the C.2.d sub-step of Jacobian.tex \S C.2: the genus-\(0\) universal-property step of \cref{thm:nonempty_jacobianWitness} that concludes the universal-pointed-morphism factorisation by ruling out positive-dimensional images of \(f \colon C \to A\) when \(C\) has genus \(0\) and a \(k\)-rational marked point.

\paragraph{Disposition (iter-155): \cref{thm:rigidity_over_kbar} is a gated NAMED GAP.} The named declaration \texttt{AlgebraicGeometry.rigidity\_over\_kbar} is landed with a \texttt{sorry} body and is held as a \textbf{named gap} in the same family as \cref{thm:nonempty_jacobianWitness} (Jacobian.tex) and the deferred Serre-duality / Riemann--Roch stack --- honestly disclosed, not on a prover lane. The substantive content (the C.2.d sub-step below, ruling out a \(1\)-dimensional image of \(f\)) reduces to the vanishing of the cotangent pull-back \(\mathrm df \colon f^{*}\Omega_{A/\bar k} \to \Omega_{C/\bar k}\), and a decisive iter-155 \texttt{mathlib-analogist} consult (\texttt{analogies/df-zero-production-iter155.md}) established that \textbf{producing \(\mathrm df = 0\) is irreducibly a global-sections fact}, gated on two genuinely-absent inputs:

\begin{itemize}
  \item \textbf{(gap i) Global cotangent-triviality of \(A\):} \(\Omega_{A/\bar k} \cong \mathcal O_A^{\oplus g}\), generated by translation-invariant differentials. The presheaf globalisation that would build this (\texttt{mulRight\_globalises\_cotangent}) was \textbf{excised from the Lean tree at iter-145}; \texttt{GrpObj.cotangentSpaceAtIdentity} \(+\) \texttt{shearMulRight} survive only as a \emph{design template}. Estimated \(\sim 800\)--\(1500\) LOC. Mathlib's only ``group cotangent is trivial'' idiom is the smooth-manifold \texttt{GroupLieAlgebra} stack, which does not port to schemes.
  \item \textbf{(gap ii) Serre duality \(H^0(C, \Omega_{C/\bar k}) = 0\)} for the genus-\(0\) curve. Given the global trivialisation of gap (i), \(\mathrm df \in H^0\bigl(C, \mathcal{H}om(f^{*}\Omega_{A/\bar k}, \Omega_{C/\bar k})\bigr) \cong \bigl(H^0(C, \Omega_{C/\bar k})\bigr)^{\oplus g}\), so \(\mathrm df = 0 \iff H^0(C, \Omega_{C/\bar k}) = 0\). Mathlib snapshot \texttt{b80f227} has no dualizing sheaf and no Riemann--Roch; the project's \texttt{Cohomology/*} stack computes \(H^1(C, \mathcal O_C)\) (genus) but has \textbf{no \(H^0\)-vanishing lemma and no genus\(\leftrightarrow\Omega\) bridge}. The genus-\(0\) hypothesis \(H^1(C, \mathcal O_C) = 0\) does \emph{not} yield \(H^0(C, \Omega_{C/\bar k}) = 0\) without the duality. Estimated \(\sim 3000\)--\(8000\) LOC.
\end{itemize}

\noindent The chart-algebra envelope of \cref{sec:RigidityKbar_shared_pile} (\cref{lem:KaehlerDifferential_mem_range_algebraMap_of_D_eq_zero} (KDM), \cref{lem:constants_integral_over_base_field}, and the delegate \cref{lem:chart_algebra_df_zero_factors_through_constant_on_chart}), which \emph{is} closed and axiom-clean, supplies only the \textbf{converse} half of the bridge: ``if \(\mathrm df = 0\) then \(f\) factors through \(\Spec \bar k\) (is constant)''. It \emph{cannot} produce \(\mathrm df = 0\): on each affine chart \(V \subseteq C\), \(\Omega_{C/\bar k}(V)\) is free of rank \(1\) (nonzero), so a chart-by-chart computation never detects the vanishing --- the vanishing materialises only on global sections. Earlier framings in this chapter (the piece-(iv) ``DEFERRED, not on the C.2.d critical path'' disposition and the ``No-Serre-duality layer'') asserted that the chart-algebra route \emph{avoids} Serre duality; \textbf{that is false}, and is corrected throughout \cref{sec:RigidityKbar_shared_pile}: \(H^0(C, \Omega_{C/\bar k}) = 0\) is \emph{on} the critical path for producing \(\mathrm df = 0\), not avoided.

The project does \textbf{not} open a ``produce \(\mathrm df = 0\)'' prover lane until one of two candidate closure routes is built; the decision between them is a leading open strategic question and the chapter commits to neither:
\begin{itemize}
  \item \textbf{Route (a):} build gaps (i) \(+\) (ii) directly.
  \item \textbf{Route (b):} the dual-abelian-variety keystone --- ``morphisms \(\mathbb P^1_{\bar k} \to A\) are constant because \(\Pic^0(\mathbb P^1_{\bar k}) = 0\)''. Earlier rejected as ``needs \(\Pic^0\)'', but since the positive-genus arm of \cref{thm:nonempty_jacobianWitness} (Route A) needs the \(\Pic^0\) engine \emph{regardless}, route (b) may amortise the shared machinery and be globally cheaper than a standalone Serre-duality build.
\end{itemize}

\paragraph{Iter-152 alg-closed pivot.} The project committed an architectural pivot: rigidity is proved \textbf{only over an algebraically closed base field \(\bar k\) of characteristic zero}, and the constancy conclusion descends to a general base field \(k\) downstream. The reason is that the chart-algebra core lemma \cref{lem:KaehlerDifferential_mem_range_algebraMap_of_D_eq_zero} (KDM) is mathematically \emph{false} as a bare \(B\)-only algebra lemma: a prover lane exhibited two counterexamples that satisfy every hypothesis yet have \(\ker D \supsetneq \operatorname{range}(\mathtt{algebraMap})\) --- namely \(B = k \times k\) (finite \'etale over \(k\)) and \(B = \mathbb Q(\sqrt 2)\) over \(\mathbb Q\). The missing content is precisely ``\(k\) algebraically closed in \(B\)''. Adding the \emph{joint} hypotheses \texttt{[IsAlgClosed kbar]} \(+\) \texttt{[IsDomain B]} makes it true: \texttt{[IsDomain B]} kills the \(B = k \times k\) counterexample and \texttt{[IsAlgClosed kbar]} kills the \(\mathbb Q(\sqrt 2)/\mathbb Q\) counterexample. (Neither hypothesis suffices alone; the iter-151 note that ``\texttt{[IsDomain B]} alone fails'' considered the hypotheses only one-at-a-time.) Over an algebraically closed base the \cref{lem:constants_integral_over_base_field} constants lemma also collapses to a one-liner via \texttt{IsAlgClosed.algebraMap\_bijective\_of\_isIntegral}, and the entire (S3.\textsc{sep}/\textsc{pi}) base-change-to-\(\bar k\) chain is descoped from the M2.a critical path.

\textbf{The Lean declaration's signature now carries \texttt{[IsAlgClosed kbar]} \(+\) \texttt{[CharZero kbar]}} (matching the theorem's existing docstring, which already described \(\bar k\) as algebraically closed --- the hypothesis was simply absent from the formal statement). Prior strategy iterations had instead routed through base-change to \(\bar k\) followed by Galois descent of morphism equality (M2.c, \(\sim 4\)--\(8\) iter / \(300\)--\(500\) LOC); the iter-127 framing then claimed all three pieces of the shared cotangent-vanishing pile build over an \emph{arbitrary} base field, dropping descent entirely. The alg-closed pivot supersedes that framing for the rigidity theorem itself: rigidity is established over \(\bar k\), and the descent of the constancy conclusion to a general \(k\) is a short downstream consequence of \texttt{AlgebraicGeometry.Flat.epi\_of\_flat\_of\_surjective} (faithfully-flat covers are epimorphisms of schemes), with the detailed descent living in the genus-\(0\) witness chapter. The group-scheme cotangent-triviality pieces (i)/(iii) below remain genuinely intrinsic to the group scheme and are stated over a general base where that is sound; only the curve-side \(\mathrm df = 0 \Rightarrow\) constancy chain (the KDM/constants core) requires the alg-closed base.

\section{Statement}

% NOTE: classical rigidity for morphisms from \(\mathbb P^1\) (or genus-0
% pointed proper curves) to abelian / group varieties; standard
% references are Mumford, *Abelian Varieties*, Ch.~II \S4 (rational
% curves on abelian varieties; paywalled, not bundled) and Hartshorne,
% *Algebraic Geometry*, Ex.~IV.4 (not bundled in references/). The
% cotangent-route closure used here ties into Stacks Tag 04QM (smooth
% over fields ⇒ geometrically reduced; references/stacks-varieties.tex,
% lemma-smooth-geometrically-normal) for the chart-level differential
% machinery.
\begin{theorem}
  \label{thm:rigidity_over_kbar}
  \lean{AlgebraicGeometry.rigidity_over_kbar}
  \leanok
  \uses{thm:GrpObj_eq_of_eqOnOpen, rem:rigidity_over_kbar_dim_dichotomy,
    lem:GrpObj_omega_free, lem:GrpObj_omega_rank_eq_dim,
    lem:GrpObj_cotangentSpace, lem:GrpObj_lieAlgebra_finrank,
    lem:Scheme_Over_ext_of_diff_zero,
    lem:chart_algebra_df_zero_factors_through_constant_on_chart,
    lem:KaehlerDifferential_mem_range_algebraMap_of_D_eq_zero,
    lem:constants_integral_over_base_field,
    thm:Scheme_AffineCoverMVSquare_HModule_prime_sequence_exact}
  % SOURCE: Mumford, Abelian Varieties, Ch. II §4 (verbatim text not yet
  % retrieved — paywalled, no open copy in references/). No SOURCE QUOTE
  % is supplied: the source is not bundled and recalled text must not be
  % substituted (citation discipline).
  \textit{Source: Mumford, Abelian Varieties, II §4.}
  Let \(\bar k\) be an algebraically closed field of characteristic zero (Lean: \texttt{[IsAlgClosed kbar]} \(+\) \texttt{[CharZero kbar]}; in the use case, the algebraic closure of a base field \(k\), with descent of the conclusion to \(k\) handled downstream in the genus-\(0\) witness chapter). These two hypotheses match the theorem's existing docstring and are required by the curve-side \(\mathrm df = 0 \Rightarrow\) constancy chain (\cref{lem:KaehlerDifferential_mem_range_algebraMap_of_D_eq_zero}, \cref{lem:constants_integral_over_base_field}); see \S "Iter-152 alg-closed pivot". Let \(C : \mathrm{Over}\,(\Spec \bar k)\) be a smooth proper geometrically irreducible curve over \(\bar k\) with \(\genus(C) = 0\) (i.e.\ \(\dim_{\bar k} H^1(C, \mathcal O_C) = 0\)). Let \(A : \mathrm{Over}\,(\Spec \bar k)\) be a smooth proper geometrically irreducible group scheme over \(\bar k\), with identity element \(\eta_A \in A(\bar k)\). Let
  \[
    f \colon C \longrightarrow A
  \]
  be a morphism of \(\bar k\)-schemes over \(\Spec \bar k\), and let \(p \in C(\bar k)\) be a \(\bar k\)-rational point such that \(f(p) = \eta_A\). Then \(f\) equals the constant morphism at \(\eta_A\):
  \[
    f \;=\; \bigl(C \to \Spec \bar k \xrightarrow{\,\eta_A\,} A\bigr).
  \]

  \emph{Encoding note (iter-126).} The Lean signature uses the abstract characterization \(\mathtt{SmoothOfRelativeDimension}\ 1\ C.\mathtt{hom} + \mathtt{IsProper}\ C.\mathtt{hom} + \mathtt{GeometricallyIrreducible}\ C.\mathtt{hom} + \genus(C) = 0\) instead of a literal \(\mathbb P^1_{\bar k}\) scheme. The two encodings are equivalent over \(\bar k\) once the curve has a \(\bar k\)-rational point (Brauer--Severi triviality + Riemann--Roch on \(C \cong \mathbb P^1_{\bar k}\)); the abstract form is used because Mathlib snapshot \texttt{b80f227} does not ship a packaged \(\mathtt{ProjectiveSpace}\,n\,S\) as a \(\mathtt{Scheme.Over}\,S\), and the C.2.b–C.2.d proof decomposition below works for any smooth proper geometrically integral genus-0 curve over \(\bar k\) without explicit \(\mathbb P^1\) API.
\end{theorem}

% SOURCE QUOTE PROOF: Mumford, Abelian Varieties, Ch. II §4 (verbatim text
% not retrieved — paywalled, no open copy in references/). No verbatim
% SOURCE QUOTE PROOF is supplied: the keystone classical input (C.2.d,
% "proper rational curves on an abelian variety are constant") is due to
% Mumford, but the source is not bundled and recalled text must not be
% substituted (citation discipline). The proof body below is a
% project-internal assembly of the chapter's own decomposition
% (C.2.b–C.2.e and the shared cotangent-vanishing pile, pieces (i)–(iv)),
% not a transcription of an external proof.
\begin{proof}
  \uses{thm:GrpObj_eq_of_eqOnOpen, rem:rigidity_over_kbar_dim_dichotomy,
    lem:GrpObj_omega_free, lem:GrpObj_omega_rank_eq_dim,
    lem:GrpObj_cotangentSpace, lem:GrpObj_lieAlgebra_finrank,
    lem:Scheme_Over_ext_of_diff_zero,
    lem:chart_algebra_df_zero_factors_through_constant_on_chart,
    lem:KaehlerDifferential_mem_range_algebraMap_of_D_eq_zero,
    lem:constants_integral_over_base_field,
    thm:Scheme_AffineCoverMVSquare_HModule_prime_sequence_exact}
  Write \(g := \dim_{\bar k} A\). The proof assembles the decomposition of
  \cref{sec:RigidityKbar_proof_decomposition} (steps C.2.b–C.2.e) on top of
  the shared cotangent-vanishing pile of \cref{sec:RigidityKbar_shared_pile}
  (pieces (i)–(iv)). The keystone step C.2.d is gated on a route choice (route
  (a) cotangent-bundle vs.\ route (b) dual-abelian-variety); we present the
  argument conditionally on whichever sub-lemmas the chapter names, so the
  closing recipe is a finite assembly of in-chapter (and cross-chapter)
  declarations.

  \emph{Reduction to the rigidity lemma (C.2.b).} Let
  \(c \colon C \to A\) denote the constant morphism at \(\eta_A\), i.e.\ the
  composite \(C \to \Spec \bar k \xrightarrow{\eta_A} A\). To prove \(f = c\)
  it suffices, by the project's rigidity lemma
  \cref{thm:GrpObj_eq_of_eqOnOpen} (\texttt{AlgebraicGeometry.Scheme.Over.ext\_of\_eqOnOpen}),
  to exhibit a non-empty open \(U \subseteq C\) on which
  \(U.\iota \circ f.\mathrm{left} = U.\iota \circ c.\mathrm{left}\) as scheme
  morphisms. The hypotheses of that lemma are met by the standing data:
  \(C\) is smooth over \(\bar k\), hence reduced, and geometrically
  irreducible; \(A\) is proper over \(\bar k\), hence separated. Thus once the
  set-level agreement is in hand on a non-empty open, the scheme-morphism
  equality \(f = c\) follows globally (this is the C.2.e promotion, integrated
  into the application of \cref{thm:GrpObj_eq_of_eqOnOpen}).

  \emph{Image-dimension dichotomy (C.2.c).} If \(g = 0\) then
  \(A = \Spec \bar k\) and the conclusion is vacuous
  (\cref{rem:rigidity_over_kbar_dim_dichotomy}); so assume \(g \geq 1\). The
  morphism \(f\) goes from the proper, geometrically integral, \(1\)-dimensional
  \(\bar k\)-scheme \(C\) to the smooth proper \(g\)-dimensional \(\bar k\)-scheme
  \(A\); its set-theoretic image \(f(C) \subseteq A\) is an irreducible closed
  subset (continuous image of an irreducible space under a proper, hence closed,
  map). The image therefore has dimension \(0\) or \(1\):
  \begin{itemize}
    \item If \(\dim f(C) = 0\), the image is a single point, which the pointed
      hypothesis \(f(p) = \eta_A\) pins to \(\{\eta_A\}\). Then \(f\) is constant
      on underlying topological points, and the set-level equality
      \(f = c\) on \(C\) is promoted to scheme-morphism equality via
      \cref{thm:GrpObj_eq_of_eqOnOpen} as in C.2.b/C.2.e.
    \item If \(\dim f(C) = 1\), the image is an irreducible rational curve in
      \(A\). This case is excluded by the keystone C.2.d below.
  \end{itemize}

  \emph{Keystone: no rational curves on a positive-dimensional abelian variety
  (C.2.d).} It remains to rule out a \(1\)-dimensional image. By Mumford's
  classical fact (\textit{Source: Mumford, Abelian Varieties, II §4}), a proper
  rational curve on an abelian variety is constant; we reduce this to the
  vanishing of the cotangent pull-back
  \(\mathrm df \colon f^{*}\Omega_{A/\bar k} \to \Omega_{C/\bar k}\). Producing
  \(\mathrm df = 0\) is a \emph{global-sections} fact, assembled as follows.
  \begin{itemize}
    \item \emph{Group-cotangent triviality (piece (i)).} The relative cotangent
      \(\Omega_{A/\bar k}\) of the group scheme \(A\) is free of rank \(g\): by
      \cref{lem:GrpObj_omega_free} it is the pull-back along \(A \to \Spec \bar k\)
      of the cotangent space at the identity \(\mathfrak g^{\vee} = \eta_A^{*}\Omega_{A/\bar k}\)
      (\cref{lem:GrpObj_cotangentSpace}), a finitely generated free
      \(\bar k\)-module, and by \cref{lem:GrpObj_omega_rank_eq_dim} (resting on
      the rank computation \cref{lem:GrpObj_lieAlgebra_finrank}) its rank equals
      \(g\). Pulling back along \(f\) gives a global trivialisation
      \(f^{*}\Omega_{A/\bar k} \cong \mathcal O_C^{\oplus g}\), under which
      \(\mathrm df\) is a \(g\)-tuple of global sections of \(\Omega_{C/\bar k}\),
      i.e.\ \(\mathrm df \in \bigl(H^0(C, \Omega_{C/\bar k})\bigr)^{\oplus g}\).
    \item \emph{Genus-\(0\) cohomological vanishing (piece (iv)).} Consequently
      \(\mathrm df = 0 \iff H^0(C, \Omega_{C/\bar k}) = 0\). The right-hand
      vanishing is the genus-\(0\) instance of the curve-side cotangent
      cohomology; it is established by a two-chart \v Cech / Mayer--Vietoris
      computation directly on \(\Omega_{C/\bar k}^{\oplus g}\) via
      \cref{thm:Scheme_AffineCoverMVSquare_HModule_prime_sequence_exact} (the
      same idiom the project uses for \(H^1(C, \mathcal O_C) = 0\) on a genus-\(0\)
      curve), rather than by a packaged Serre-duality theorem. \emph{This step
      is gated on the route choice:} route (a) builds pieces (i)\(+\)(ii)
      directly to obtain the trivialisation and the vanishing; route (b) obtains
      the same conclusion through the dual abelian variety via
      \(\Pic^0(\mathbb P^1_{\bar k}) = 0\). The argument below is stated
      conditionally on whichever route supplies \(\mathrm df = 0\); neither route
      is committed here.
    \item \emph{Converse: \(\mathrm df = 0 \Rightarrow f\) constant (piece (ii)).}
      Given \(\mathrm df = 0\), the chart-algebra envelope furnishes constancy.
      On each affine chart \(V = \Spec R \subseteq C\) lying over a chart
      \(W = \Spec B \subseteq A\), the per-chart helper
      \cref{lem:chart_algebra_df_zero_factors_through_constant_on_chart} shows
      that the chart map \(f^{\sharp} \colon B \to R\) factors as
      \(B \to \bar k \to R\): it delegates the ring-side kernel extraction to
      the K\"ahler core \cref{lem:KaehlerDifferential_mem_range_algebraMap_of_D_eq_zero}
      (KDM) and the integral-closure of the constants to
      \cref{lem:constants_integral_over_base_field}. The hypotheses
      \texttt{[CharZero kbar]} and \texttt{[IsAlgClosed kbar]} carried by the
      theorem are exactly what these two helpers require: characteristic zero
      makes the Frobenius-iteration layer (piece (iii)) vacuous, so the
      derivation-kernel argument applies directly, and algebraic closedness
      collapses the constants lemma to the bijectivity of
      \(\mathtt{algebraMap}\,\bar k\,\bar k(b)\). Assembling the chart
      factorisations by sheaf compatibility and globalising via
      \cref{lem:Scheme_Over_ext_of_diff_zero}, the morphism \(f\) factors
      through \(\Spec \bar k\); hence \(f\) is constant, contradicting
      \(\dim f(C) = 1\). This rules out the \(1\)-dimensional image.
  \end{itemize}

  \emph{Conclusion (C.2.b/C.2.e).} The dichotomy C.2.c thus collapses to the
  single point case: \(f(C) = \{\eta_A\}\). Taking \(U := C\), the resulting
  set-level agreement \(f = c\) is promoted to scheme-morphism equality by
  \cref{thm:GrpObj_eq_of_eqOnOpen} (reduced source \(C\), separated target
  \(A\)). Therefore
  \[
    f \;=\; \bigl(C \to \Spec \bar k \xrightarrow{\,\eta_A\,} A\bigr),
  \]
  as claimed.
\end{proof}

\begin{remark}
  \label{rem:rigidity_over_kbar_dim_dichotomy}
  
  The conclusion is vacuously true when \(A\) is zero-dimensional: in that case \(A = \Spec \bar k\) and every morphism \(\mathbb P^1_{\bar k} \to A\) factors through \(\Spec \bar k\) trivially. The substantive content of the theorem is the positive-dimensional case, which decomposes into the steps of \cref{sec:RigidityKbar_proof_decomposition}.
\end{remark}

\section{Proof decomposition}
\label{sec:RigidityKbar_proof_decomposition}

The proof of \cref{thm:rigidity_over_kbar} decomposes as follows. We reuse the sub-step labelling of Jacobian.tex \S C.2 so cross-references between the two chapters remain unambiguous. \textbf{This decomposition is the eventual closing recipe for a gated named gap} (see the iter-155 disposition in the chapter introduction): the substantive step C.2.d reduces to producing \(\mathrm df = 0\), which is gated on gaps (i)\(+\)(ii) [route (a)] or on the dual-AV keystone [route (b)]. The steps below stay inert --- the Lean body is a \texttt{sorry} --- until a route is chosen and built; none names a present, assembled Mathlib infrastructure for the image-dimension and \(\mathrm df = 0\) content.

\begin{itemize}
  \item[\textbf{(C.2.b)}] \emph{Reduction to the project's rigidity lemma.} Let \(c \colon \mathbb P^1_{\bar k} \to A\) denote the constant morphism at \(\eta_A\). The conclusion \(f = c\) follows from \cref{thm:GrpObj_eq_of_eqOnOpen} (i.e.\ \texttt{AlgebraicGeometry.Scheme.Over.ext\_of\_eqOnOpen}, after the iter-125 refactor) applied to \(g_1 := f\), \(g_2 := c\), \(X := \mathbb P^1_{\bar k}\), \(Y := A\), once we produce a non-empty open \(U \subseteq \mathbb P^1_{\bar k}\) on which \(U.\iota \circ f.\mathrm{left} = U.\iota \circ c.\mathrm{left}\) as scheme morphisms. We will take \(U := \mathbb P^1_{\bar k}\). The refactored rigidity lemma's hypotheses are trivially satisfied: \(\mathbb P^1_{\bar k}\) is smooth (hence reduced) and geometrically irreducible over \(\bar k\), and \(A\) is proper (hence separated) over \(\bar k\).

  \item[\textbf{(C.2.c)}] \emph{Image dimension argument.} The morphism \(f\) goes from a proper geometrically integral \(\bar k\)-scheme of dimension \(1\) to a smooth proper \(\bar k\)-scheme of dimension \(g := \dim_{\bar k} A\). The set-theoretic image \(f(\mathbb P^1_{\bar k}) \subseteq A\) is an irreducible closed subset of \(A\) (continuous image of an irreducible space under a proper map). \emph{Gating note (iter-155): this dichotomy names no present, assembled Mathlib infrastructure} --- the ``image of a proper map is closed and irreducible'' fact and a usable notion of scheme/image dimension are not packaged for this setting in snapshot \texttt{b80f227}; C.2.c is part of the gated assembly, to be detailed once a route (a)/(b) is chosen. Dichotomy:
  \begin{itemize}
    \item If \(g = 0\): \(A = \Spec \bar k\) and the conclusion is vacuous (\cref{rem:rigidity_over_kbar_dim_dichotomy}).
    \item If \(g \geq 1\) and the image is a single point: the image must be \(\{\eta_A\}\) (the pointed hypothesis \(f(p) = \eta_A\) pins it). Then \(f\) is constant on underlying topological points, and the conclusion follows by promoting set-level equality to scheme-morphism equality (sub-step C.2.e of Jacobian.tex, applied via \cref{thm:GrpObj_eq_of_eqOnOpen}).
    \item If \(g \geq 1\) and the image is one-dimensional: the image is an irreducible rational curve in \(A\). This case is ruled out by the keystone classical input (C.2.d below).
  \end{itemize}

  \item[\textbf{(C.2.d)}] \emph{Keystone: proper rational curves on an abelian variety are constant --- the gated content.} The remaining case --- the image of \(f\) has dimension exactly \(1\) --- is ruled out by Mumford's classical fact stated in the chapter introduction, and is the load-bearing reduction of the whole theorem. It reduces to the vanishing \(\mathrm df = 0\) of the cotangent pull-back \(f^{*}\Omega_{A/\bar k} \to \Omega_{C/\bar k}\). Per the iter-155 disposition (chapter introduction) and \texttt{analogies/df-zero-production-iter155.md}, producing \(\mathrm df = 0\) is a \textbf{global-sections} fact: with the global trivialisation \(f^{*}\Omega_{A/\bar k} \cong \mathcal O_C^{\oplus g}\) (gap i), \(\mathrm df = 0 \iff H^0(C, \Omega_{C/\bar k}) = 0\) (gap ii, Serre duality on the genus-\(0\) curve). The two standard proofs sketched in Jacobian.tex \S C.2.d furnish the two gated closure routes that the project commits to neither of yet: \textbf{route (a)} = the cotangent-bundle route, building gaps (i)\(+\)(ii) directly (triviality of \(\Omega_{A/\bar k}\) + \(H^0(C,\Omega_{C/\bar k}) = 0\)); \textbf{route (b)} = the dual-abelian-variety route via \(\Pic^0(\mathbb P^1_{\bar k}) = 0\), which may amortise the \(\Pic^0\) engine that the positive-genus arm of \cref{thm:nonempty_jacobianWitness} needs regardless. The chart-algebra envelope of \cref{sec:RigidityKbar_shared_pile} supplies only the converse ``\(\mathrm df = 0 \Rightarrow\) constant''; it does \emph{not} produce \(\mathrm df = 0\), and (contrary to earlier framing) does not avoid the Serre-duality input.

  \item[\textbf{(C.2.e)}] \emph{Promotion of set-level equality to scheme-morphism equality.} Standard reduced-source / separated-target argument, integrated into the application of \cref{thm:GrpObj_eq_of_eqOnOpen} in C.2.b.
\end{itemize}

\section{The shared cotangent-vanishing Mathlib pile}
\label{sec:RigidityKbar_shared_pile}

The C.2.d keystone above (\emph{morphisms from \(\mathbb P^1_{\bar k}\) to a positive-dimensional abelian variety over \(\bar k\) are constant}) and the genus-\(0\) identification alternative of \cref{thm:nonempty_jacobianWitness} (M2.d-alt of STRATEGY.md, which sidesteps Riemann--Roch by running the cotangent-vanishing argument directly) depend on the same Mathlib pile. STRATEGY.md \S M2.d-alt names this the \emph{shared cotangent-vanishing pile}. The iter-126 mathlib-analogist consult (\texttt{analogies/cotangent-vanishing-pile.md}) decomposed the pile into the four pieces below, revised the per-piece LOC budgets upward where the directive had under-scoped, and made the characteristic-\(p\) option pick:

\begin{itemize}
  \item[(i)] \emph{Group-scheme cotangent triviality.} The relative cotangent presheaf \(\Omega_{G/k}\) of a smooth proper geometrically irreducible group scheme \(G\) over a base field \(k\) (this piece is genuinely intrinsic to the group scheme and needs no algebraic-closure hypothesis, even under the iter-152 alg-closed pivot of \S "Iter-152 alg-closed pivot"; only the curve-side \(\mathrm df = 0 \Rightarrow\) constancy chain requires \(\bar k\) algebraically closed) is free of rank \(\dim G\). Per the iter-126 analogist's naming-idiom alignment (Mathlib snapshot \texttt{b80f227} has \emph{no} \texttt{AbelianVariety} file; the \texttt{GrpObj} namespace per Yang+Merten 2026 is the canonical home), target names are \texttt{AlgebraicGeometry.GrpObj.omega\_free} and \texttt{AlgebraicGeometry.GrpObj.omega\_rank\_eq\_dim}. A prover-ready three-step sub-lemma decomposition (i.a)\,\(\to\)\,(i.b)\,\(\to\)\,(i.c) for the iter-128+ build lane is given in \cref{subsec:RigidityKbar_piece_i_decomposition}. Estimated build: \textbf{800--1500 LOC} (iter-126 analogist's correction over an earlier 200--400 LOC framing; the triviality is the chain Lie-algebra-of-\texttt{GrpObj} + shear-iso globalisation + relative-cotangent-presheaf trivialisation on top of the project's existing \texttt{relativeDifferentialsPresheaf}, inheriting the presheaf-vs-sheaf bridge cost from \texttt{analogies/cotangent-presheaf-design.md}).

  \item[(ii)] \emph{Scheme-level \(\mathrm d f = 0\) implies \(f\) factors through \(\Spec k\).} The companion algebraic fact: if a morphism \(f \colon X \to Y\) between smooth proper \(k\)-schemes has \(\mathrm d f \colon f^* \Omega_{Y/k} \to \Omega_{X/k}\) the zero map (e.g.\ because \(\Omega_{Y/k}\) is trivial and \(H^0(X, \Omega_{X/k}) = 0\)), then \(f\) factors through the constant morphism. In characteristic \(0\) this is the standard differential-vanishing-implies-constancy argument. In characteristic \(p > 0\) it requires the Frobenius iteration of piece (iii). Per the analogist's recommendation: the ring-level half aligns with Mathlib's \texttt{Differential.ContainConstants} typeclass (\texttt{Mathlib.RingTheory.Derivation.DifferentialRing}, lines 62--70; one Mathlib consumer at \texttt{FieldTheory.Differential.Liouville}); the scheme-level half is a NEEDS\_MATHLIB\_GAP\_FILL gap and the canonical name is \texttt{AlgebraicGeometry.Scheme.Over.ext\_of\_diff\_zero} (following the iter-125-refactored \texttt{ext\_of\_eqOnOpen} idiom). Estimated build: \textbf{250--500 LOC}.

  \item[(iii)] \emph{Characteristic-\(p\) handling --- \textbf{committed Option A: Frobenius iteration}.} Per the analogist's decisive pick (\texttt{analogies/cotangent-vanishing-pile.md} \S "Option A wins"): the three options on the table are
  \begin{itemize}
    \item Frobenius iteration (replace \(f\) with \(f \circ F^n\) for the \(n\)-th iterate of Frobenius; descent of the conclusion uses smoothness of \(A\));
    \item Mumford's argument (Mumford, \emph{Abelian Varieties}, Chapter~II~\S4: abelian varieties contain no rational curves in characteristic \(p\));
    \item Lifting to characteristic \(0\) via Witt vectors.
  \end{itemize}
  Options B (Mumford) and C (Witt) each carry standalone multi-thousand-LOC dependencies (Mumford requires the very Pic\(^0\) / cotangent machinery the project is trying to streamline; Witt-vector lifting requires a scheme-level Witt functor + smooth deformation + descent that Mathlib entirely lacks). Option A consumes Mathlib's first-class ring-level \texttt{frobenius} / \texttt{iterateFrobenius} from \texttt{Mathlib.Algebra.CharP.Frobenius} (lines 30--110+) and aligns with the project's existing scheme-level Frobenius design context (\texttt{AlgebraicJacobian/Rigidity.lean:48--55}); the scheme-level Frobenius lift remains a NEEDS\_MATHLIB\_GAP\_FILL gap, but Option A's net cost is dominated by adapting the ring-level idiom upward to schemes. Estimated build: \textbf{300--600 LOC}.

  \item[(iv)] \emph{Serre duality on a smooth proper curve --- \textbf{DEFERRED named gap, ON the C.2.d critical path (iter-155 correction)}.} The honest-closure verdict (\texttt{analogies/serre-duality.md}, iter-110) is \textbf{3000--8000 LOC}. \textbf{Correction (iter-155).} Earlier text here claimed piece (iv) was used ``by the M2.d-alt genus-\(0\) identification alternative, NOT by C.2.d itself'' and that ``C.2.d / M2.a body closure does NOT depend on piece (iv)''. \emph{That is false.} The iter-155 \texttt{mathlib-analogist} consult (\texttt{analogies/df-zero-production-iter155.md}) established that the C.2.d keystone reduces to producing \(\mathrm df = 0\), which is a global-sections fact: with the gap-(i) trivialisation \(f^{*}\Omega_{A/\bar k} \cong \mathcal O_C^{\oplus g}\), one has \(\mathrm df \in \bigl(H^0(C, \Omega_{C/\bar k})\bigr)^{\oplus g}\), so \(\mathrm df = 0 \iff H^0(C, \Omega_{C/\bar k}) = 0\) --- which is exactly the genus-\(0\) instance of this Serre-duality piece (iv). Hence piece (iv) (or an equivalent that yields \(H^0(C, \Omega_{C/\bar k}) = 0\)) is \textbf{on} the C.2.d critical path under the cotangent-bundle closure route (a); only the dual-AV route (b) can in principle bypass it, by amortising the \(\Pic^0\) engine. The chart-by-chart KDM stack does \emph{not} avoid this input (see the correction in \cref{subsec:RigidityKbar_piece_i_decomposition}). Piece (iv) remains a deferred named gap; it is not opened on a prover lane until a closure route is committed.
\end{itemize}

\subsection{Piece (i): sub-lemma decomposition for iter-128+ build}
\label{subsec:RigidityKbar_piece_i_decomposition}

We decompose piece (i) above into three sub-steps (i.a)\,\(\to\)\,(i.b)\,\(\to\)\,(i.c) that are well-suited to staged Lean formalisation. Sub-step (i.a) defines the Lie algebra (as a \(k\)-vector space) of a group scheme at the identity and pins its rank; sub-step (i.b) globalises that fibre across \(G\) via a functorial shear isomorphism; sub-step (i.c) concludes the freeness and rank of the relative cotangent \(\Omega_{G/k}\).

All three sub-steps of this group-scheme cotangent-triviality piece are stated over an arbitrary base field \(k\) (the cotangent triviality of \(G\) is intrinsic to the group scheme and does not need the alg-closed hypothesis; see \S "Iter-152 alg-closed pivot" in the chapter introduction, where the alg-closed base is required only for the curve-side \(\mathrm df = 0 \Rightarrow\) constancy chain, not here); the prose carefully avoids any reliance on \(k\)-rational or \(\bar k\)-rational point lifting, working functorially throughout. Sub-step (i.a) is the natural first target for the iter-128+ prover lane: it has the smallest signature surface (a finite-rank free \(k\)-module attached to the identity section of \(G\)), and no scheme-level globalisation.

\paragraph{Iter-144 chart-algebra pivot --- disposition (per \texttt{mathlib-analogist-chart-algebra-iter144} + STRATEGY.md \S "Iter-144 chart-algebra pivot" block).}
Effective iter-145+, the M2.body-pile commits to the \textbf{chart-algebra route}. Under this commitment:
\begin{itemize}
  \item Piece (i.b) Step 2 d\_app sub-sorry inside \cref{lem:GrpObj_omega_basechange_proj}'s proof, the IsIso theorem \cref{lem:GrpObj_basechange_along_proj_two_inv_app_isIso} (Lean target \texttt{AlgebraicGeometry.GrpObj.basechange\_along\_proj\_two\_inv\_app\_isIso}), and piece (i.b) Main \cref{lem:GrpObj_mulRight_globalises} (Lean target \texttt{AlgebraicGeometry.GrpObj.mulRight\_globalises\_cotangent}) are \textbf{DESCOPED from critical-path}. \emph{Status correction (iter-155):} \texttt{mulRight\_globalises\_cotangent} was subsequently \textbf{EXCISED from the Lean tree at iter-145}; \texttt{GrpObj.cotangentSpaceAtIdentity} \(+\) \texttt{shearMulRight} survive only as a \emph{design template}, \textbf{not} a present declaration. This excised gap-(i) globalisation (\(\sim 800\)--\(1500\) LOC) is exactly the gap (i) that the keystone \(\mathrm df = 0\) (chapter introduction) is gated on; the blueprint blocks below for these declarations are therefore an off-path record of the bundled-route design, not a live in-tree obligation.
  \item Piece (i.c.1/i.c.2/i.c.3) (\cref{lem:GrpObj_omega_free} + \cref{lem:GrpObj_omega_rank_eq_dim} assembly) is \textbf{DESCOPED}. Chart-algebra routes the freeness and rank facts through chart-level \texttt{Algebra.IsStandardSmooth.free\_kaehlerDifferential} per-chart, not through a global sheaf-of-modules construction.
  \item Piece (iii) scheme-level absolute Frobenius PHANTOM (\textasciitilde\(800\)--\(1500\) LOC, \texttt{Scheme.absoluteFrobenius} ABSENT in Mathlib snapshot \texttt{b80f227}) is \textbf{DESCOPED}; chart-algebra uses ring-level \texttt{RingHom.iterateFrobenius\_comm} from \texttt{Mathlib.Algebra.CharP.Frobenius} per-chart, giving Frobenius iteration as a direct \texttt{RingHom}-commutation with no scheme-level lift.
  \item Piece (ii) is \textbf{INFLATED} to absorb the chart-algebra (\(\alpha\)) \texttt{Algebra.IsPushout}-from-affine-product helper and (\(\beta\)) per-chart translation-invariance Kähler-derivation argument upstream; see the \S "Iter-144 chart-algebra envelope for piece (ii)" paragraph below for the updated envelope (\textasciitilde\(600\)--\(1050\) LOC total).
\end{itemize}

Piece (i.a) --- \cref{lem:GrpObj_cotangentSpace} (\texttt{cotangentSpaceAtIdentity}) and \cref{lem:GrpObj_lieAlgebra_finrank} (\texttt{cotangentSpaceAtIdentity\_finrank\_eq}), closed iter-128 \(\to\) iter-132 --- remains in-tree as standalone infrastructure. Chart-algebra does \emph{not} consume it on the rigidity critical path, but the cotangent-space-at-the-identity definition and its rank are independently useful downstream (Lie-algebra dual \(\mathfrak g\) on the tangent side, etc.).

\paragraph{Iter-144 chart-algebra envelope for piece (ii) (per \texttt{mathlib-analogist-chart-algebra-iter144}).}
Piece (ii) (scheme-level \(\mathrm df = 0 \Rightarrow f\) factors through \(\Spec k\), the C.2.d / M2.a body consumer of piece (i)'s cotangent-trivial output) is the load-bearing body under the iter-144 chart-algebra pivot. Total envelope \textasciitilde\(600\)--\(1050\) LOC, decomposed as:
\begin{itemize}
  \item Chart-algebra (\(\alpha\)) \texttt{Algebra.IsPushout}-from-affine-product helper: \textasciitilde\(80\)--\(150\) LOC. Constructs the \texttt{Algebra.IsPushout} instance on the chart-level ring extensions \(\Gamma(\Spec k, U) \to \Gamma(W, \mathcal O_G)\) for affine charts \(V \to W\) in the product \(G \otimes G\). Built from \texttt{pullbackSpecIso} and the \texttt{\_hom\_fst} / \texttt{\_inv\_fst} simp companions in \texttt{Mathlib.AlgebraicGeometry.Pullbacks}, chained with \texttt{CommRingCat.isPushout\_iff\_isPushout} and \texttt{isPullback\_SpecMap\_of\_isPushout} (cf.\ \cref{lem:GrpObj_omega_basechange_proj} Route (a) item 1).
  \item Chart-algebra (\(\beta\)) per-chart translation-invariance Kähler-derivation argument: \textasciitilde\(150\)--\(300\) LOC. The load-bearing core of M2.a's \(\mathrm df = 0\) derivation chain (per STRATEGY.md \S "Soundness rules: M2.a \(\mathrm df = 0\) derivation chain --- articulated (added iter-144)"). The three-layer chain is:
    \begin{enumerate}
      \item \textbf{Chart-local layer.} On each affine chart \(W = \Spec B \subseteq A.\mathrm{left}\) with affine \(V \subseteq C.\mathrm{left}\) such that \(f\) restricts to a ring map \(f^{\sharp} \colon B \to \Gamma(C, V)\), the differential \(\mathrm df = 0\) at the chart level means the composite \(\Omega_{B/k} \to \Omega_{\Gamma(C, V)/k}\) obtained from the universal Kähler derivation \(\mathrm d \colon B \to \Omega_{B/k}\) post-composed with \(f^{\sharp}\) vanishes. Combined with piece (i)'s chart-level free-of-rank-\(n\) presentation of \(\Omega_{B/k}\) (\texttt{Algebra.IsStandardSmooth.free\_kaehlerDifferential}), the kernel-extraction step gives \(f^{\sharp}(b) \in \Gamma(C, V)^{p^{e}}\) for every \(b \in B\), where \(p = \operatorname{char} k\) and \(e \geq 0\) is a chart-uniform Frobenius depth.
      \item \textbf{Char-\(p\) Frobenius-iteration layer.} The chart-local conclusion \(f^{\sharp}(b) \in \Gamma(C, V)^{p^{e}}\) is promoted to "for every \(b \in B\), there exists \(e\) with \(f^{\sharp}(b) = \mathrm{iterateFrobenius}_{p, e}(c_{b})\) for some \(c_{b} \in \Gamma(C, V)\)" via \texttt{RingHom.iterateFrobenius\_comm} (\texttt{Mathlib.Algebra.CharP.Frobenius}, verified iter-144). The uniform-\(e\)-on-\(W\) glue extracts a globally-bounded Frobenius depth on \(C\) from the smoothness of \(A\) at the image of \(f\) (using \(A\)'s tangent-bundle triviality at each point in the image of \(f\) to control the depth uniformly across \(W\)). In characteristic \(0\) this layer is vacuous: \(\mathrm df = 0 \Rightarrow f^{\sharp}(b) \in k \subseteq \Gamma(C, V)\) directly via the standard derivation-kernel argument (no Frobenius).
      \item \textbf{Constancy-from-\(\mathrm df = 0\) layer (the converse only).} \emph{Iter-155 correction.} This layer was previously titled the ``No-Serre-duality layer'' and claimed the chart-algebra (\(\beta\)) chain ``replaces'' the \(H^0(C, \Omega_{C/k}) = 0\) route, ``eliminating piece (iv) as a critical-path dependency''. \textbf{That claim is false} (per \texttt{analogies/df-zero-production-iter155.md}) and is corrected here. What the chain actually delivers is the \emph{converse} implication: \textbf{given} the chart-uniform input \(f^{\sharp}(b) \in \Gamma(C, V)^{p^{e}}\) (which is the chart-level shadow of the hypothesis \(\mathrm df = 0\)), the conclusion ``\(f\) factors through \(\Spec k\)'' follows from the integrally-closed-constants helper (next item) and the scheme-level lift via \texttt{Scheme.Over.ext\_of\_eqOnOpen} (final item). It does \emph{not} produce the hypothesis \(\mathrm df = 0\), and it does not avoid the \(H^0(C, \Omega_{C/k}) = 0\) input: that input is exactly what is needed to \emph{establish} \(\mathrm df = 0\) in the first place. Concretely, \(\mathrm df = 0\) is a global-sections fact --- with the gap-(i) trivialisation \(f^{*}\Omega_{A/k} \cong \mathcal O_C^{\oplus g}\), \(\mathrm df \in \bigl(H^0(C, \Omega_{C/k})\bigr)^{\oplus g}\), so \(\mathrm df = 0 \iff H^0(C, \Omega_{C/k}) = 0\) --- and the chart-by-chart machinery cannot detect it because \(\Omega_{C/k}(V)\) is free of rank \(1\) (nonzero) on every affine chart \(V\). Hence piece (iv) (or an equivalent yielding \(H^0(C, \Omega_{C/k}) = 0\)) is \textbf{on} the critical path for the keystone, not eliminated by this chain.
    \end{enumerate}
  \item Algebra-level core \texttt{KaehlerDifferential.mem\_range\_algebraMap\_of\_D\_eq\_zero}: \textasciitilde\(200\)--\(350\) LOC. Mathlib-shaped helper for the ring-level half of \(\mathrm df = 0 \Rightarrow f^{\sharp}\) factors through \(k\), packaged as the statement "for \(D \colon R \to M\) the universal Kähler derivation with \(D \circ f^{\sharp} = 0\), the image of \(f^{\sharp}\) lies in the integral closure of the image of the base ring". This is the chart-uniform core consumed by layer (1) of the \(\beta\)-chain.
  \item Integrally-closed-constants helper: \textasciitilde\(50\)--\(100\) LOC. The wrapper around "constants in \(C\) are integral over \(k\) for \(C\) smooth proper geometrically irreducible", used to promote the chart-uniform Frobenius-depth conclusion to a global \(\Spec k\)-factorisation.
  \item Scheme-level lift via \texttt{Scheme.Over.ext\_of\_eqOnOpen}: \textasciitilde\(100\)--\(150\) LOC. Closes the argument by exhibiting a non-empty open \(U \subseteq C\) on which \(f|_{U}\) equals the constant morphism at the image of \(f\)'s chart-uniform output, then applies the iter-125-refactored \cref{thm:GrpObj_eq_of_eqOnOpen} (\texttt{AlgebraicGeometry.Scheme.Over.ext\_of\_eqOnOpen}) to globalise.
\end{itemize}

The chart-algebra route replaces the bundled scheme-level absolute-Frobenius construction (piece (iii) PHANTOM) with chart-uniform ring-level idioms already in Mathlib snapshot \texttt{b80f227}. \textbf{It does not, however, replace the Serre-duality cohomology vanishing (piece (iv))} --- the iter-155 correction above and in \cref{subsec:RigidityKbar_piece_i_decomposition} retracts that earlier claim. The chart-algebra (\(\beta\)) chain is the \emph{converse} ``\(\mathrm df = 0 \Rightarrow\) constant''; producing \(\mathrm df = 0\) remains a global-sections fact gated on \(H^0(C, \Omega_{C/k}) = 0\) (piece (iv)) plus the gap-(i) trivialisation. The Lean target name for the scheme-level converse half is \texttt{AlgebraicGeometry.Scheme.Over.ext\_of\_diff\_zero} (per the iter-126 analogist's idiom alignment); see \cref{lem:Scheme_Over_ext_of_diff_zero} for the present thin-wrapper Lean decl versus the gated substantive \(\mathrm df = \mathrm dg\) refinement.

\paragraph{Iter-128 / iter-129 prover lane re-scoping (per \texttt{strategy-critic-iter128} + iter-129 rename).} Sub-step (i.a) was originally a 2-target pair and is now (post-iter-129) a 3-target trio: \cref{lem:GrpObj_cotangentSpace} (the \emph{definition} of \(\mathfrak g^\vee\), plus the finitely-generated-free claim — iter-128 lane, body landed), \cref{lem:GrpObj_cotangent_bridge} (the \emph{bridge} to the local-ring cotangent space \(\mathfrak m / \mathfrak m^{2}\) — iter-130+ build target), and \cref{lem:GrpObj_lieAlgebra_finrank} (the \emph{rank-equals-relative-dimension} claim — iter-130+ build target). The iter-128 prover lane attacks ONLY the first: it scaffolds the Lean declaration \texttt{AlgebraicGeometry.GrpObj.cotangentSpaceAtIdentity} (renamed iter-129 from the placeholder \texttt{lieAlgebra}; signature relaxed to a free \texttt{\{n : \mathbb N\} [SmoothOfRelativeDimension n G.hom]} binder so downstream consumers of arbitrary relative dimension can instantiate it) and closes its body. The bridge lemma \cref{lem:GrpObj_cotangent_bridge} is vestigial on the live path under Replacement (B) (see \S Q4 collapse footer at the rank-lemma proof's end) and remains deferred indefinitely; the rank lemma \cref{lem:GrpObj_lieAlgebra_finrank} is the iter-132 prover-lane target (closure path: Step 1 chart-side Kähler rank from \cref{thm:smooth_locally_free_omega}'s existential + Step 2 base-change to \(k\) via \texttt{Module.finrank\_baseChange}, both detailed in that lemma's proof body). Rationale: (i) bundling definition + bridge + rank into one prover-iter task is borderline-too-ambitious per \texttt{progress-critic-iter128}; (ii) the rank-lemma RHS Lean encoding needs blueprint pinning iter-129 before being prover-ready (it must read \texttt{Module.finrank k (cotangentSpaceAtIdentity G) = n} parameterised by an explicit \texttt{[SmoothOfRelativeDimension n G.hom]} instance; the bare \(\dim G\) in the prose does not type-check directly because Mathlib has no \(\dim G\) function on \texttt{Over (Spec k)}, per \texttt{strategy-critic-iter128} CHALLENGE on (i.a) signature ambiguity).

Throughout this subsection, \(k\) is an arbitrary field and \(G : \Over\,(\Spec k)\) is a smooth proper geometrically irreducible group scheme over \(k\), with identity section \(\eta_G : \mathbf{1} \to G\).

\begin{lemma}\leanok
[(i.a) Cotangent space at the identity of a group scheme]
  \label{lem:GrpObj_cotangentSpace}
  \lean{AlgebraicGeometry.GrpObj.cotangentSpaceAtIdentity}
  \uses{thm:smooth_locally_free_omega}
  % Lean signature stub (pinned post-iter-129 fixup; relaxed from the
  % iter-128 hardcoded rel-dim-1 form so the same declaration serves
  % consumers of arbitrary relative dimension, e.g. \cref{thm:rigidity_over_kbar}
  % applied to an abelian variety of dimension \(g\)):
  %   noncomputable def cotangentSpaceAtIdentity (G : Over (Spec (.of k)))
  %       [CategoryTheory.GrpObj G] {n : ℕ} [SmoothOfRelativeDimension n G.hom]
  %       [IsProper G.hom] [GeometricallyIrreducible G.hom] : ModuleCat k
  The pullback along the identity section \(\eta_G\) of the relative cotangent presheaf,
  \[
    \mathfrak g^{\vee} \;:=\; \eta_G^{*}\,\Omega_{G/k},
  \]
  is a finitely generated free \(k\)-module. We call \(\mathfrak g^{\vee}\) the \emph{cotangent space of \(G\) at the identity}. The \emph{Lie algebra of \(G\)} as a \(k\)-vector space is recovered as the \(k\)-linear dual \(\mathfrak g := \mathfrak g^{\vee\vee} = \mathrm{Hom}_{k}(\mathfrak g^{\vee}, k)\) (the bracket structure is not used in the rigidity argument and is not packaged here).

  \emph{Lean encoding note (iter-129 dualisation pinning).} The Lean declaration \texttt{AlgebraicGeometry.GrpObj.cotangentSpaceAtIdentity} returns a \texttt{ModuleCat k} which is the \emph{un-dualised} cotangent \(\mathfrak g^{\vee}\) — the direct output of the pullback-along-section construction described in the proof, with no internal \texttt{Module.Dual} hop. Downstream consumers that need the tangent \(\mathfrak g\) take \texttt{Module.Dual k (cotangentSpaceAtIdentity G)} externally. This convention pins what was left open in the iter-128 encoding note: the rigidity argument is symmetric under dualisation (only the finite-rank-free claim is consumed), so the un-dualised form is preferred as it matches the body construction without an extra contravariant hop. The declaration embeds into Mathlib's \texttt{ModuleCat} infrastructure and supplies the \texttt{Module.Free k} / \texttt{Module.Finite k} typeclasses needed for the rank lemma \cref{lem:GrpObj_lieAlgebra_finrank} via the bridge lemma \cref{lem:GrpObj_cotangent_bridge}.
\end{lemma}

\begin{proof}

\leanok
  The identity section \(\eta_G\) is a canonical \(k\)-rational point of \(G\) supplied by the \(\GrpObj\) structure. The iter-131 Lean construction (Replacement (B), pure-term \texttt{noncomputable def}) realises \(\mathfrak g^{\vee} = \eta_G^{*}\,\Omega_{G/k}\) as a chart-base-changed Kähler module: by the forward Jacobian criterion \cref{thm:smooth_locally_free_omega} (already in the tree as \texttt{AlgebraicGeometry.Scheme.smooth\_locally\_free\_omega} in \texttt{AlgebraicJacobian/Differentials.lean}), there exists an affine chart \(V \subseteq G\) around the image of the identity section \(\eta_G(\mathrm{pt})\), with affine base chart \(U \subseteq \Spec k\) and the structural ring map \(\Gamma(\Spec k, U) \to \Gamma(G, V)\) exhibiting \texttt{Algebra.IsStandardSmoothOfRelativeDimension} \(n\) (and consequently \(\Omega_{\Gamma(G,V) \,/\, \Gamma(\Spec k,\,U)}\) is a free \(\Gamma(G, V)\)-module of rank \(n\) by \texttt{Algebra.IsStandardSmoothOfRelativeDimension.rank\_kaehlerDifferential} in \texttt{Mathlib.RingTheory.Smooth.StandardSmoothCotangent}). The cotangent space at the identity is then defined as the base change to \(k\) of this algebraic Kähler module along the ring map \(\psi_V \colon \Gamma(G, V) \to k\) obtained by restricting the identity section \(\eta_G\) to \(V\) and composing with the canonical isomorphism \(\Gamma(\Spec k, \cdot) \cong k\):
  \[
    \mathfrak g^{\vee} \;\;:=\;\; k \,\otimes_{\Gamma(G,\,V)}\, \Omega_{\Gamma(G,\,V) \,/\, \Gamma(\Spec k,\,U)}.
  \]
  This is a finitely generated free \(k\)-module of rank \(n\) by base-change of a free module along \(\psi_V\) (the rank is preserved; see \cref{lem:GrpObj_lieAlgebra_finrank}).

  \emph{Caveat on canonicity.} The construction depends on a chart choice \(V\) extracted via a \texttt{Classical.choose}-chain from the existential statement of \texttt{smooth\_locally\_free\_omega}; it is therefore not canonically attached to \(G\) in the strictest sense. (The iter-131 fix-up refactored what was previously a top-level \texttt{Classical.choice}-wrapper into a pure-term \texttt{let}-chain of \texttt{Classical.choose} / \texttt{Classical.choose\_spec} accessors, exposing the outer head symbol \texttt{(ModuleCat.extendScalars\,\(\psi_V\).hom).obj\,(ModuleCat.of\,\(\Gamma\)(G, V)\,\(\Omega\)[\(\Gamma\)(G, V)\,/\,\(\Gamma\)(Spec k, U)])} to downstream consumers; see \S "Iter-131 \texttt{Classical.choose}-chain body shape" below for the Lean encoding details, and the companion \cref{lem:GrpObj_cotangentSpace_extendScalars_witness} for the structural-shape rewrite handle.) For the live consumer (rigidity over \(k\), via \cref{thm:rigidity_over_kbar}), the existence of a rank-\(n\) free \(k\)-module is the only structural content needed, so canonicity is not load-bearing. A canonical (chart-independent) presentation via the stalk-side cotangent \(\mathfrak m_{\eta_G}/\mathfrak m_{\eta_G}^{2}\) is available as a future bridge (\cref{lem:GrpObj_cotangent_bridge} below, currently \texttt{\textbackslash notready}) at additional cost; the project defers that bridge until a non-rigidity consumer materialises that requires canonicity.
\end{proof}

\begin{lemma}\leanok
[Structural-shape rewrite handle for the cotangent space at the identity]
  \label{lem:GrpObj_cotangentSpace_extendScalars_witness}
  \lean{AlgebraicGeometry.GrpObj.cotangentSpaceAtIdentity_eq_extendScalars}
  \uses{lem:GrpObj_cotangentSpace, thm:smooth_locally_free_omega}
  % Lean signature stub (iter-131 strong-acceptance companion lemma, closed
  % in-tree as `AlgebraicGeometry.GrpObj.cotangentSpaceAtIdentity_eq_extendScalars`;
  % promoted to a first-class blueprint citizen iter-133 per the iter-132/iter-133
  % blueprint-reviewer MED-B finding so it appears in the dependency graph
  % alongside the rank-lemma route patterns; statement matches the Lean
  % declaration verbatim):
  %   theorem cotangentSpaceAtIdentity_eq_extendScalars
  %       (G : Over (Spec (.of k))) [CategoryTheory.GrpObj G]
  %       {n : ℕ} [SmoothOfRelativeDimension n G.hom]
  %       [IsProper G.hom] [GeometricallyIrreducible G.hom] :
  %       ∃ (U : (Spec (.of k)).Opens) (V : G.left.Opens) (e : V ≤ G.hom ⁻¹ᵁ U)
  %           (htop : (⊤ : (Spec (.of k)).Opens) ≤
  %             (CategoryTheory.CommaMorphism.left η[G]) ⁻¹ᵁ V),
  %         letI : Algebra ↥Γ(Spec (.of k), U) ↥Γ(G.left, V) :=
  %           (Scheme.Hom.appLE G.hom U V e).hom.toAlgebra
  %         (cotangentSpaceAtIdentity (n := n) G : ModuleCat k) =
  %           (ModuleCat.extendScalars
  %             ((CategoryTheory.CommaMorphism.left η[G]).appLE V ⊤ htop ≫
  %               (Scheme.ΓSpecIso (.of k)).hom).hom).obj
  %             (ModuleCat.of Γ(G.left, V) Ω[Γ(G.left, V) ⁄ Γ(Spec (.of k), U)])
  
  Under the hypotheses of \cref{lem:GrpObj_cotangentSpace}, there exist affine opens \(U \subseteq \Spec k\), \(V \subseteq G.\mathrm{left}\), an inclusion \(e \colon V \leq G.\mathrm{hom}^{-1}\,U\), and a top-inclusion \(h_{\mathrm{top}} \colon \top \leq \eta_{G}^{-1}\,V\) such that, under the \(\Gamma(\Spec k, U)\)-algebra structure on \(\Gamma(G, V)\) from \texttt{appLE}, the cotangent space at the identity equals the explicit chart-base-changed Kähler module:
  \[
    \mathfrak g^{\vee} \;=\; \texttt{cotangentSpaceAtIdentity}\,G \;=\; (\texttt{ModuleCat.extendScalars}\,\psi_{V}.\mathrm{hom}).\mathrm{obj}\bigl(\texttt{ModuleCat.of}\,\Gamma(G, V)\,\Omega[\Gamma(G, V)/\Gamma(\Spec k, U)]\bigr),
  \]
  where \(\psi_{V} \colon \Gamma(G, V) \to k\) is the ring map obtained by restricting the identity section \(\eta_{G}\) to \(V\) and composing with the canonical iso \(\Gamma(\Spec k, \cdot) \cong k\).
\end{lemma}

\begin{proof}
  
  \leanok
  Reproduce the body's \texttt{Classical.choose}-chain on \cref{thm:smooth_locally_free_omega}'s existential (the same chart triple \((U, V, e)\) and the same membership witness \(h_{xV} \colon x_{0} \in V\)). The top-inclusion \(h_{\mathrm{top}} \colon \top \leq \eta_{G}^{-1}\,V\) follows from \(h_{xV}\) via \texttt{Subsingleton.elim} on points of \(\Spec k\) (a field's \(\Spec\) has a unique point). The displayed equation closes by \texttt{rfl} once the right-hand side's \texttt{Classical.choose}-chain is unpacked on the same existential as the body's, since both extract the same chart witnesses. The Lean proof is closed in-tree as \texttt{AlgebraicGeometry.GrpObj.cotangentSpaceAtIdentity\_eq\_extendScalars}.
\end{proof}

\begin{lemma}[(i.a-bridge) Bridge to the local-ring cotangent space at the identity]
  \label{lem:GrpObj_cotangent_bridge}
  \lean{AlgebraicGeometry.TODO.GrpObj_cotangent_bridge}
  \uses{lem:GrpObj_cotangentSpace, lem:kaehler_localization_subsingleton}
  % Iter-146 disposition: \lean{...cotangentSpaceAtIdentity_iso_localRingCotangent}
  % stripped — no Lean declaration of that name exists in tree (was a
  % Replacement (A) bridge proposed in
  % \texttt{analogies/lieAlgebra-rank-bridge.md} iter-129 but never built);
  % broken hint would mislead any downstream tool following the dependency
  % graph. The block remains \notready as before.
  % Lean signature stub (iter-130+ build target; bridges the iter-131
  % chart-base-changed Kähler body of \cref{lem:GrpObj_cotangentSpace} to
  % the local-ring cotangent space supplied by Mathlib). Currently
  % \notready under Replacement (B): the live rank-lemma closure path
  % (\cref{lem:GrpObj_lieAlgebra_finrank} Steps 1+2) bypasses this
  % bridge entirely, so it is vestigial on the live path and is
  % deferred until a non-rigidity consumer materialises that requires
  % the canonical (chart-independent) stalk-side presentation
  % \mathfrak m / \mathfrak m^{2} — see the rank-lemma proof's footer
  % paragraph for the iter-131 strategy-critic Q4 collapse verdict.
  %   noncomputable def cotangentSpaceAtIdentity_iso_localRingCotangent
  %       (G : Over (Spec (.of k)))
  %       [CategoryTheory.GrpObj G] {n : ℕ} [SmoothOfRelativeDimension n G.hom]
  %       [IsProper G.hom] [GeometricallyIrreducible G.hom] :
  %       cotangentSpaceAtIdentity G ≅
  %         ModuleCat.of k (Ideal.IsLocalRing.CotangentSpace
  %           (R := Scheme.stalk G.left (η_G.left.base ⟨_, trivial⟩)))
  
  \notready
  Let \(G\) be as in \cref{lem:GrpObj_cotangentSpace}, and let \(\mathcal O_{G, \eta_G}\) denote the local ring of \(G\) at the identity (the stalk along \(\eta_G : \Spec k \to G\)), with maximal ideal \(\mathfrak m_{\eta_G}\) and residue map \(\pi \colon \mathcal O_{G, \eta_G} \to k\). Then there is a canonical \(k\)-linear isomorphism
  \[
    \mathfrak g^{\vee} \;=\; \eta_G^{*}\,\Omega_{G/k}
    \;\xrightarrow{\;\sim\;}\;
    \mathfrak m_{\eta_G} \big/ \mathfrak m_{\eta_G}^{2}.
  \]
  Concretely, the right-hand side is \texttt{Ideal.IsLocalRing.CotangentSpace} applied to the local ring \(\mathcal O_{G, \eta_G}\) (as a module over its residue field, identified with \(k\) via the \(\GrpObj\)-supplied identity section); the left-hand side is the iter-130+ chart-base-changed Kähler Lean body of \cref{lem:GrpObj_cotangentSpace}.
\end{lemma}

\begin{proof}
  
  The bridge is the composition of two standard identifications (both at the level of \(k\)-modules):

  \emph{Step 1: from the chart-base-changed Kähler module to the differential module of the local stalk-residue map.} The iter-131 body of \cref{lem:GrpObj_cotangentSpace} is the chart-base-change
  \[
    k \;\otimes_{\Gamma(G, V)}\; \Omega_{\Gamma(G, V)\,/\,\Gamma(\Spec k,\,U)}
  \]
  along the ring map \(\psi_V \colon \Gamma(G, V) \to k\) obtained from the identity section restricted to \(V\). To bridge this chart-base-changed module to the stalk-side cotangent, factor \(\psi_V\) through the localisation at the identity stalk (and a parallel base-side localisation \(\Gamma(\Spec k, U) \to k\)):
  \[
    \Gamma(G, V) \;\xrightarrow{\;\text{loc}\;}\; \mathcal O_{G, \eta_G} \;\xrightarrow{\;\pi\;}\; k,
    \qquad
    \Gamma(\Spec k, U) \;\xrightarrow{\;\text{loc}\;}\; k.
  \]
  Localisation kills relative differentials (\texttt{Algebra.FormallyUnramified.subsingleton\_kaehlerDifferential}, used already as \cref{lem:kaehler_localization_subsingleton} of Differentials.tex), so extension of scalars along the localisation step on each side preserves \(\Omega\) as the Kähler differential module \(\Omega_{\mathcal O_{G, \eta_G}/k}\). The full extension along \(\psi_V\) — factored through the local stalk — therefore identifies \(k \otimes_{\Gamma(G, V)} \Omega_{\Gamma(G, V)/\Gamma(\Spec k, U)}\) with \(k \otimes_{\mathcal O_{G, \eta_G}} \Omega_{\mathcal O_{G, \eta_G}/k}\), independent of the chart \(V\).

  \emph{Step 2: from the Kähler module of a local \(k\)-algebra with residue field \(k\) to \(\mathfrak m / \mathfrak m^{2}\).} For a local ring \((R, \mathfrak m, k)\) that is a \(k\)-algebra with residue map \(\pi \colon R \to k\) a \(k\)-algebra retraction (here supplied by the identity section \(\eta_G\) producing \(R = \mathcal O_{G, \eta_G}\) over \(k\) with residue field canonically \(k\)), the standard identification
  \[
    k \,\otimes_{R}\, \Omega_{R/k} \;\;\cong\;\; \mathfrak m \big/ \mathfrak m^{2}
  \]
  holds via the universal derivation \(\mathrm{d} \colon R \to \Omega_{R/k}\) restricted to \(\mathfrak m\) (the kernel of \(\pi\)). This is Stacks Tag~02G1 (or Hartshorne~II.8, Proposition~II.8.7 in the Noetherian case) and is in Mathlib at the ring level via \texttt{Algebra.Extension.CotangentSpace} infrastructure (\texttt{Mathlib.RingTheory.Extension.Cotangent.Basic}); the present blueprint statement chooses the local-ring framing \texttt{Ideal.IsLocalRing.CotangentSpace} (\texttt{Mathlib.RingTheory.Ideal.Cotangent}, the abbreviation \texttt{(maximalIdeal R).Cotangent}) as the cleanest container for the rank step in \cref{lem:GrpObj_lieAlgebra_finrank}.

  Composing Steps 1 and 2 yields the displayed isomorphism. The bridge identifies the chart-base-changed Kähler module of Replacement (B) (depending on a chosen chart \(V\) extracted via the iter-131 \texttt{Classical.choose}-chain on \cref{thm:smooth_locally_free_omega}) with the canonical stalk-side cotangent \(\mathfrak m_{\eta_G} / \mathfrak m_{\eta_G}^{2}\). The two are abstractly isomorphic as \(k\)-modules — both have rank \(n\) and arise as cotangent spaces at the identity — but the isomorphism is non-trivial to construct: it is a localisation-along-the-identity-stalk step (Step 1 above) composed with the standard \(k \otimes_{R} \Omega_{R/k} \cong \mathfrak m / \mathfrak m^{2}\) identification (Step 2 above). The cost of this bridge is estimated at 300--600 LOC and is deferred until a non-rigidity consumer requires canonicity (e.g.\ a future Lie-algebra-bracket consumer of \(\mathfrak g\)). Regularity enters only at the rank-pinning step (\cref{lem:GrpObj_lieAlgebra_finrank}).
\end{proof}

\begin{lemma}\leanok
[(i.a) Rank of the cotangent space at the identity]
  \label{lem:GrpObj_lieAlgebra_finrank}
  \lean{AlgebraicGeometry.GrpObj.cotangentSpaceAtIdentity_finrank_eq}
  \uses{lem:GrpObj_cotangentSpace, thm:smooth_locally_free_omega}
  % Lean signature stub (iter-130+ rank-lemma build target; depends on the
  % bridge lemma \cref{lem:GrpObj_cotangent_bridge} plus the regular-local-ring
  % + standard-smooth rank-pinning Mathlib pieces named in the proof):
  %   theorem cotangentSpaceAtIdentity_finrank_eq (G : Over (Spec (.of k)))
  %       [CategoryTheory.GrpObj G] {n : ℕ} [SmoothOfRelativeDimension n G.hom]
  %       [IsProper G.hom] [GeometricallyIrreducible G.hom] :
  %       Module.finrank k (cotangentSpaceAtIdentity G) = n
  
  In the setting of \cref{lem:GrpObj_cotangentSpace}, with \([\mathtt{SmoothOfRelativeDimension}\,n\,G.\mathtt{hom}]\) supplying the relative dimension parameter,
  \[
    \finrank_{k}\,\mathfrak g^{\vee} \;=\; n,
  \]
  and consequently (since rank is preserved under \(k\)-linear duality on finitely-generated free modules) \(\finrank_{k}\,\mathfrak g = n\) as well.

  \emph{Lean encoding note (iter-128 per \texttt{strategy-critic-iter128}; iter-129 rename).} The right-hand side is the relative-dimension parameter \(n\) from the \texttt{[SmoothOfRelativeDimension n G.hom]} instance, NOT a bare \(\dim G\) (which has no canonical Lean function on \texttt{Over (Spec k)}). The mathematical content is unchanged because for a smooth proper geometrically irreducible group scheme, the relative dimension equals the dimension of the cotangent space at the identity equals what informal prose calls \(\dim G\). The Lean target name is \texttt{cotangentSpaceAtIdentity\_finrank\_eq} (iter-129 rename, replacing the iter-128 placeholder \texttt{lieAlgebra\_finrank\_eq\_dim}, aligned with the renamed (i.a) declaration \cref{lem:GrpObj_cotangentSpace}).
\end{lemma}

\begin{proof}
  
  \leanok
  We close the rank computation under Replacement (B) (the iter-131 chart-base-change body of \cref{lem:GrpObj_cotangentSpace}) using the existential output of \cref{thm:smooth_locally_free_omega} directly. The proof is in four steps: Steps 1+2 form the live closure path; Step 3 documents an alternative canonical route through \(\mathfrak m / \mathfrak m^{2}\) that is currently deferred under Replacement (B); Step 4 carries the conclusion across \(k\)-linear duality to the tangent side.

  \emph{Step 1: chart-side Kähler rank.} Let \(V \subseteq G.\mathrm{left}\) and \(U \subseteq \Spec k\) be the affine charts (around the image \(x_0 = \eta_G(\mathrm{pt}) \in G.\mathrm{left}\)) extracted by \cref{thm:smooth_locally_free_omega} (project lemma \texttt{AlgebraicGeometry.Scheme.smooth\_locally\_free\_omega} in \texttt{AlgebraicJacobian/Differentials.lean}); the same charts are pinned by the body of \cref{lem:GrpObj_cotangentSpace} via the iter-131 \texttt{Classical.choose}-chain and re-extractable through the companion lemma \texttt{cotangentSpaceAtIdentity\_eq\_extendScalars} (see \S "Iter-131 \texttt{Classical.choose}-chain body shape" below for the recommended downstream rewrite pattern). The existential's payload, on the \(\Gamma(\Spec k, U)\)-algebra \(\Gamma(G, V)\) (with the algebra structure given by the appLE ring map of \(G.\mathrm{hom}\)), supplies in one shot both
  \[
    \texttt{hfree} \colon \;\texttt{Module.Free}\;\Gamma(G, V)\;\Omega_{\Gamma(G, V)\,/\,\Gamma(\Spec k,\,U)}
    \qquad\text{and}\qquad
    \texttt{hrank} \colon \;\texttt{Module.rank}\;\Gamma(G, V)\;\Omega_{\Gamma(G, V)\,/\,\Gamma(\Spec k,\,U)} \;=\; n,
  \]
  internally consuming the Mathlib lemmas \texttt{Algebra.IsStandardSmooth.free\_kaehlerDifferential} and \texttt{Algebra.IsStandardSmoothOfRelativeDimension.rank\_kaehlerDifferential} (\texttt{Mathlib.RingTheory.Smooth.StandardSmoothCotangent}). Converting the rank-as-cardinal equation \texttt{hrank} to a \(\finrank\) equation via \texttt{Module.finrank\_eq\_of\_rank\_eq} (\texttt{Mathlib.LinearAlgebra.Dimension.Finrank}; for a commutative semiring \(R\) with \texttt{StrongRankCondition}, \(\texttt{Module.rank}\,R\,M = \uparrow n \Rightarrow \texttt{Module.finrank}\,R\,M = n\)) gives
  \[
    \finrank_{\Gamma(G, V)}\,\Omega_{\Gamma(G, V)\,/\,\Gamma(\Spec k,\,U)} \;=\; n.
  \]

  \emph{Step 2: base-change to \(k\) preserves \(\finrank\).} By construction, the body of \cref{lem:GrpObj_cotangentSpace} is the base change \(k \otimes_{\Gamma(G, V)} \Omega_{\Gamma(G, V)/\Gamma(\Spec k, U)}\) along the ring map \(\psi_V \colon \Gamma(G, V) \to k\) obtained from the identity section restricted to \(V\) (composed with \(\Gamma(\Spec k, \cdot) \cong k\)). The Mathlib base-change lemma
  \[
    \texttt{Module.finrank\_baseChange} \colon\;\;
    \texttt{Module.finrank}\;R\;(\texttt{TensorProduct}\;S\;R\;M')
      \;=\; \texttt{Module.finrank}\;S\;M'
  \]
  (\texttt{Mathlib.LinearAlgebra.Dimension.Constructions}; requires \(S\) a commutative semiring with \texttt{StrongRankCondition}, \(M'\) a \texttt{Module.Free}\;\(S\)-module, and \(R\) an \(S\)-algebra with \texttt{StrongRankCondition}), applied with \(S := \Gamma(G, V)\), \(R := k\), and \(M' := \Omega_{\Gamma(G, V)/\Gamma(\Spec k, U)}\) — using \texttt{hfree} from Step 1 to discharge the \texttt{Module.Free} hypothesis — yields
  \[
    \finrank_{k}\,\mathfrak g^{\vee} \;=\; \finrank_{k}\bigl(k \otimes_{\Gamma(G, V)} \Omega_{\Gamma(G, V)/\Gamma(\Spec k, U)}\bigr) \;=\; \finrank_{\Gamma(G, V)}\,\Omega_{\Gamma(G, V)/\Gamma(\Spec k, U)} \;=\; n,
  \]
  which closes the rank lemma on the live (B) path.

  \emph{Companion freeness fact (downstream consumer; not on the rank-lemma critical path).} For downstream consumers that separately need \(k\)-freeness of \(\mathfrak g^{\vee}\) — notably \cref{lem:GrpObj_omega_free} of piece (i.c) — the typeclass \texttt{Algebra.TensorProduct.instFree} (\texttt{Mathlib.RingTheory.TensorProduct.Free}) supplies \(\texttt{Module.Free}\;k\;(k \otimes_{\Gamma(G, V)} \Omega_{\Gamma(G, V)/\Gamma(\Spec k, U)})\) directly from \texttt{hfree}. This is informational here; the rank lemma above closes without invoking it.

  \emph{Step 3 (alternative canonical route, currently deferred).} An alternative closure via the stalk-side cotangent space \(\mathfrak m_{\eta_G}/\mathfrak m_{\eta_G}^{2}\) is available in principle and was the iter-128-era plan. The route is: (i) bridge from the chart-base-changed body to the local-stalk Kähler module via \cref{lem:GrpObj_cotangent_bridge} (currently \texttt{\textbackslash notready}; identifies \(k \otimes_{\Gamma(G, V)} \Omega_{\Gamma(G, V)/\Gamma(\Spec k, U)} \cong \mathfrak m_{\eta_G}/\mathfrak m_{\eta_G}^{2}\), via a localisation-at-the-identity-stalk step composed with the standard \(k \otimes_{R} \Omega_{R/k} \cong \mathfrak m / \mathfrak m^{2}\) identification — Stacks Tag~02G1 / Hartshorne~II.8.7); (ii) invoke regularity of \(\mathcal O_{G, \eta_G}\) (\texttt{IsRegularLocalRing}; from smoothness of \(G \to \Spec k\) at the closed point \(\eta_G\), Stacks Tag~00TR / Hartshorne~II.10.3.B), with Krull dimension of \(\mathcal O_{G, \eta_G}\) equal to the relative dimension \(n\) (Stacks Tag~02G1); (iii) read off the rank from the defining property of regular local rings, \(\finrank_{k}\,\mathfrak m_{\eta_G}/\mathfrak m_{\eta_G}^{2} = \dim \mathcal O_{G, \eta_G} = n\), packaged into \texttt{Ideal.IsLocalRing.CotangentSpace} (\texttt{Mathlib.RingTheory.Ideal.Cotangent}) and \texttt{IsRegularLocalRing} (\texttt{Mathlib.RingTheory.RegularLocalRing.Defs}). This route is preferred only when a downstream consumer requires the canonical (chart-independent) stalk-side presentation; it is \emph{not} on the live closure path under Replacement (B) because the bridge lemma \cref{lem:GrpObj_cotangent_bridge} is \texttt{\textbackslash notready}. The footer paragraph below records the iter-131 strategy-critic verdict on this deferral.

  \emph{Step 4: dualisation conclusion (\(\finrank\,\mathfrak g = n\)).} By the standard rank-equals-rank-of-dual identity for finitely-generated free modules over a field,
  \[
    \finrank_{k}\,\mathfrak g \;=\; \finrank_{k}\,\mathrm{Hom}_{k}(\mathfrak g^{\vee}, k) \;=\; \finrank_{k}\,\mathfrak g^{\vee} \;=\; n,
  \]
  delivering the consequent \(\finrank_{k}\,\mathfrak g = n\) on the tangent side.

  \emph{Mathlib name summary (live closure path under Replacement (B)).} The closure pieces consumed by Steps 1+2 above are: \cref{thm:smooth_locally_free_omega} (project lemma \texttt{AlgebraicGeometry.Scheme.smooth\_locally\_free\_omega}; supplies the existential \(\exists\,U\,V\,e\,(h_{xV} : x_0 \in V)\,(h_U\,h_V : \texttt{IsAffineOpen}\;\cdot)\,(\texttt{hfree})\,(\texttt{hrank})\) in one shot, internally consuming \texttt{Algebra.IsStandardSmoothOfRelativeDimension.rank\_kaehlerDifferential} and \texttt{Algebra.IsStandardSmooth.free\_kaehlerDifferential}), \texttt{Module.finrank\_eq\_of\_rank\_eq} (\texttt{Mathlib.LinearAlgebra.Dimension.Finrank}; rank-as-cardinal to \(\finrank\) converter), and \texttt{Module.finrank\_baseChange} (\texttt{Mathlib.LinearAlgebra.Dimension.Constructions}; base-change preserves \(\finrank\) when the source module is \texttt{Module.Free}). The companion freeness typeclass \texttt{Algebra.TensorProduct.instFree} (\texttt{Mathlib.RingTheory.TensorProduct.Free}) is informational for downstream consumers (e.g.\ \cref{lem:GrpObj_omega_free}) but does not enter the rank-lemma proof. All four names are verified in Mathlib snapshot \texttt{b80f227} by the iter-132 plan agent.

  \emph{Iter-131 strategy-critic Q4 collapse (footer).} The iter-131 strategy-critic Q4 review of \cref{subsec:RigidityKbar_piece_i_decomposition} documented in STRATEGY.md that, under Replacement (B), the trio (definition + bridge + rank) collapses to a duo (definition + rank): the rank lemma closes directly from \cref{thm:smooth_locally_free_omega}'s existential without routing through the bridge \cref{lem:GrpObj_cotangent_bridge}. The bridge therefore becomes vestigial on the live path and is currently \texttt{\textbackslash notready}; it re-enters the live closure only if trigger (a') at iter-133+ piece (i.b) fires — i.e., only if a non-rigidity consumer materialises that requires the canonical (chart-independent) stalk-side cotangent presentation \(\mathfrak m_{\eta_G}/\mathfrak m_{\eta_G}^{2}\) that the bridge supplies.
\end{proof}

\begin{lemma}\leanok
[(i.b-helper) Binary-product shear isomorphism for a group object]
  \label{lem:GrpObj_shearMulRight}
  \lean{AlgebraicGeometry.GrpObj.shearMulRight}
  % Lean signature stub (iter-134 substantive closure; the binary-product
  % upgrade of `CategoryTheory.GrpObj.mulRight`
  % (`Mathlib.CategoryTheory.Monoidal.Grp_.lean:277--281`), which packages
  % only the one-input version `(· · f)` for a fixed global element
  % `f : 𝟙_ ⟶ G`. Here both factors of `G ⊗ G` are free, parametrising
  % the family of right-translations {mulRight a}_{a ∈ G} globally across G.
  % Stated for any `GrpObj G` in any Cartesian monoidal category C, so the
  % declaration is general-purpose and instantiates at `Over (Spec (.of k))`
  % for the rigidity application):
  %   def shearMulRight {C : Type*} [Category C] [CartesianMonoidalCategory C]
  %       (G : C) [GrpObj G] : G ⊗ G ≅ G ⊗ G where
  %     hom := lift (fst G G) μ
  %     inv := lift (fst G G) (lift (fst G G ≫ ι) (snd G G) ≫ μ)
  % together with the `@[simps]`-spawned companion lemmas
  %   shearMulRight_hom_fst : (shearMulRight G).hom ≫ fst G G = fst G G
  %   shearMulRight_hom_snd : (shearMulRight G).hom ≫ snd G G = μ
  
  Let \(G\) be a group object in a Cartesian monoidal category. The \emph{binary-product shear isomorphism} is the self-iso
  \[
    \sigma \colon G \otimes G \;\xrightarrow{\;\sim\;}\; G \otimes G
  \]
  whose forward and inverse maps are built categorically from the multiplication \(\mu\), the inverse \(\iota\), and the binary-product projections of \(G \otimes G\):
  \[
    \sigma_{\mathrm{hom}} \;=\; \texttt{lift\,(fst\,G\,G)\,μ}, \qquad
    \sigma_{\mathrm{inv}} \;=\; \texttt{lift\,(fst\,G\,G)\,(lift\,(fst\,G\,G\,≫\,\(\iota\))\,(snd\,G\,G)\,≫\,μ)},
  \]
  informally \(\sigma\colon (a, b) \mapsto (a, a \cdot b)\) with inverse \((a, b) \mapsto (a, a^{-1} \cdot b)\). The forward map is an isomorphism over the first projection and acts as right-translation on the second coordinate; concretely,
  \[
    \sigma_{\mathrm{hom}} \circ \pr_{1} \;=\; \pr_{1}, \qquad
    \sigma_{\mathrm{hom}} \circ \pr_{2} \;=\; \mu.
  \]
\end{lemma}

\begin{proof}
  
  \leanok
  The compositions \(\sigma_{\mathrm{hom}} \circ \sigma_{\mathrm{inv}} = \mathrm{id}_{G \otimes G} = \sigma_{\mathrm{inv}} \circ \sigma_{\mathrm{hom}}\) are verified componentwise via the binary-product universal property (\texttt{CartesianMonoidalCategory.hom\_ext}). The first-projection components are trivial (both sides collapse to \(\pr_{1}\) by \texttt{lift\_fst}). The second-projection components reduce to identities of the form
  \[
    \texttt{lift\,(fst\,G\,G\,≫\,\(\iota\))\,\(\mu\)\,≫\,\(\mu\)} \;=\; \pr_{2}
    \quad\text{and}\quad
    \texttt{lift\,(fst\,G\,G)\,(lift\,(fst\,G\,G\,≫\,\(\iota\))\,(snd\,G\,G)\,≫\,\(\mu\))\,≫\,\(\mu\)} \;=\; \pr_{2},
  \]
  closed using the Mathlib calculus \texttt{MonObj.lift\_lift\_assoc} (associativity-via-lift reshuffling), \texttt{CategoryTheory.GrpObj.lift\_comp\_inv\_left} / \texttt{lift\_comp\_inv\_right} (left/right-inverse axioms for \(\GrpObj\)), and \texttt{MonObj.lift\_comp\_one\_left} (left-unit axiom). The companion identifications \(\sigma_{\mathrm{hom}} \circ \pr_{1} = \pr_{1}\) and \(\sigma_{\mathrm{hom}} \circ \pr_{2} = \mu\) follow directly from \texttt{lift\_fst} and \texttt{lift\_snd} on \(\sigma_{\mathrm{hom}} = \texttt{lift\,(fst\,G\,G)\,μ}\) and are auto-generated by the \texttt{@[simps]} attribute on \texttt{shearMulRight} into the lemmas \texttt{shearMulRight\_hom\_fst} and \texttt{shearMulRight\_hom\_snd}. The proof shape parallels \texttt{CategoryTheory.GrpObj.mulRight} (\texttt{Mathlib.CategoryTheory.Monoidal.Grp\_.lean:277--281}) and \texttt{CategoryTheory.GrpObj.isPullback} (\texttt{Grp\_.lean:293--323}), but with both factors of \(G \otimes G\) free rather than the one-input version's fixed global element.
\end{proof}

\begin{lemma}[(i.b) Shear-iso globalisation of the cotangent]
  \label{lem:GrpObj_mulRight_globalises}
  \lean{AlgebraicGeometry.TODO.GrpObj_mulRight_globalises}
  \uses{lem:GrpObj_shearMulRight, lem:GrpObj_omega_basechange_proj, lem:GrpObj_omega_restrict_to_identity_section}
  % Iter-146 disposition: \lean{...mulRight_globalises_cotangent}
  % stripped — Lean target EXCISED iter-145 from
  % `AlgebraicJacobian/Cotangent/GrpObj.lean` under the iter-144
  % chart-algebra pivot (the bundled-route piece (i.b) main lemma is
  % DESCOPED from critical-path; the iter-145 prover lane removed the
  % sorry-bodied scaffold rather than keeping it as auditable record).
  % Block kept as traceable iter-133 \(\to\) iter-143 decision history;
  % see also EXCISE disposition paragraph above
  % \cref{lem:GrpObj_omega_basechange_proj}.
  % Lean signature stub (iter-133 blueprint-writer hardening per the iter-133
  % mathlib-analogist verdict at
  % `task_results/mathlib-analogist-mulright-globalises-iter133.md` Decision 4
  % and the persistent design file `analogies/mulright-globalises-cotangent.md`;
  % iter-134+ prover-lane target). The RHS is stated at the
  % `PresheafOfModules`-of-`G.left` level using the abstract pullback
  % presheaf `η_G^* (relativeDifferentialsPresheaf G.hom)`, NOT at the
  % `ModuleCat k`-valued `cotangentSpaceAtIdentity G` level. Decoupling
  % piece (i.b)'s body from chart localisation in this way keeps the lemma
  % in the 210--440 LOC envelope; the chart-localisation bridge from this
  % sheaf-level RHS to `cotangentSpaceAtIdentity G` is a separate
  % ~100--200 LOC chart-localisation identification consumed inside
  % piece (i.c) (`omega_free`), NOT inside this lemma's body:
  %   noncomputable def mulRight_globalises_cotangent
  %       (G : Over (Spec (.of k))) [CategoryTheory.GrpObj G]
  %       {n : ℕ} [SmoothOfRelativeDimension n G.hom]
  %       [IsProper G.hom] [GeometricallyIrreducible G.hom] :
  %       relativeDifferentialsPresheaf G.hom ≅
  %         (PresheafOfModules.pullback (φ_str G)).obj
  %           ((PresheafOfModules.pullback (φ_η G)).obj
  %             (relativeDifferentialsPresheaf G.hom))
  % where:
  %   - `PresheafOfModules.pullback` is the Mathlib functor
  %     (`Mathlib.Algebra.Category.ModuleCat.Presheaf.Pullback`, signature
  %     `(PresheafOfModules S) ⥤ (PresheafOfModules R)` parameterised by a
  %     compatibility morphism `φ : S ⟶ F.op.comp R`).
  %   - `φ_str` packages the compatibility morphism for the structure map
  %     `G.hom : G.left ⟶ Spec (.of k)`; `φ_η` the compatibility morphism for
  %     `η[G].left : Spec (.of k) ⟶ G.left`. Both are obtained via the
  %     `TopCat.Presheaf.pullbackPushforwardAdjunction` of `CommRingCat`-valued
  %     structure presheaves, in the same shape as the `φ'` on
  %     `AlgebraicJacobian/Differentials.lean:52`.
  %   - The home category is `G.left.PresheafOfModules`. The conclusion
  %     `Iso` is inside `Mathlib.Algebra.Category.ModuleCat.Presheaf` (the
  %     presheaf-of-modules-level statement; NOT the sheafified
  %     `SheafOfModules`-level version, which would require an extra
  %     sheafification step not needed for piece (i.b)).
  
  % NOTE iter-135: the Lean target
  % \lean{AlgebraicGeometry.GrpObj.mulRight_globalises_cotangent} ships
  % as an honest sorry-bodied scaffold; the signature is the intended
  % sheaf-level RHS displayed above (the compatibility morphisms for
  % `PresheafOfModules.pullback` are constructed via Mathlib's
  % `Scheme.Hom.toRingCatSheafHom`, per the iter-135 mathlib-analogist
  % verdict in `analogies/phi-compatibility-morphisms.md`). The proof
  % body is `sorry`; `sync_leanok` removes the prior iter-134 `\leanok`
  % from the proof block accordingly. Body closure is iter-136+ work
  % per the Step 1/2/3/Compose outline below.
  % NOTE iter-140: \notready stripped — Lean target
  % `\lean{AlgebraicGeometry.GrpObj.mulRight_globalises_cotangent}`
  % is formalized as an honest sorry-bodied scaffold at
  % `Cotangent/GrpObj.lean:806`. Per CLAUDE.md marker vocab,
  % `\notready` on a formalized statement is stale.
  % NOTE iter-146: the Lean target was subsequently EXCISED iter-145
  % from `AlgebraicJacobian/Cotangent/GrpObj.lean` under the iter-144
  % chart-algebra pivot; the statement is once again blueprint-only.
  \notready
  Let \(\mu \colon G \times_{k} G \to G\) denote the multiplication of \(G\) and let \(\sigma \colon G \times_{k} G \to G \times_{k} G\) be the \emph{shear isomorphism}
  \[
    \sigma \;=\; \langle \pr_{1},\, \mu \rangle, \qquad (a, b) \;\mapsto\; (a,\, a \cdot b),
  \]
  built categorically from \(\mu\) and the binary-product projections in \(\Over\,(\Spec k)\). Then \(\sigma\) is an isomorphism over the first projection, and pullback along \(\sigma\) together with the canonical first-projection-cotangent identification \(\Omega_{(G \times_{k} G)/G} \cong \pr_{2}^{*}\,\Omega_{G/k}\) (\cref{lem:GrpObj_omega_basechange_proj}) yields a natural isomorphism of presheaves of modules on \(G\)
  \[
    \Omega_{G/k} \;\cong\; \pi_{G}^{*}\!\bigl(\eta_{G}^{*}\,\Omega_{G/k}\bigr) \qquad \text{on } G,
  \]
  where \(\pi_{G} \colon G \to \Spec k\) is the structure map (in the standard category-theory notation for the binary product in \(\Over\,(\Spec k)\), this is the "first projection" of the singleton-product-with-base, sometimes written \(\pr_{1}\); we use \(\pi_{G}\) to avoid confusion with the binary-product first projection \(\pr_{1} \colon G \times_{k} G \to G\) appearing above). Informally: the relative cotangent of \(G\) is canonically isomorphic to the constant \(\mathcal O_{G}\)-module with fibre \(\eta_{G}^{*}\,\Omega_{G/k}\), trivialised along the structure map.

  \paragraph{Iter-127 over-k risk register.} The globalisation \textbf{must} be formulated functorially via the shear scheme map \(\sigma\) — which is a morphism in \(\Over\,(\Spec k)\), valid for any base field — and \textbf{not} as a pointwise translation by closed (or \(\bar k\)-rational) points of \(G\). The pointwise-translation argument requires either alg-closure of \(k\) or a Nullstellensatz-style \(\bar k\)-point lifting (as in Mathlib's \texttt{pointEquivClosedPoint}, used in the alg-closed-base branch of \texttt{Mathlib.AlgebraicGeometry.Group.Smooth.lean:50--51}); the functorial shear formulation sidesteps that obstacle entirely. The iter-133 mathlib-analogist verdict (\texttt{analogies/mulright-globalises-cotangent.md} Decision 4) explicitly rules out the pointwise idioms of \texttt{Mathlib.Geometry.Manifold.GroupLieAlgebra.mulInvariantVectorField} (requires \texttt{[Group G]}) and \texttt{Mathlib.Topology.Algebra.Group.Basic.Homeomorph.shearMulRight} (requires \texttt{[Group G]} on the carrier) as the over-k-forbidden alternatives.

  \paragraph{Lean target name vs.\ construction (option (i), iter-133).} The named target \texttt{mulRight\_globalises\_cotangent} retains the \texttt{mulRight} naming convention because the shear iso \(\sigma = \langle \pr_{1}, \mu \rangle\) on \(G \times_{k} G\) is the binary-product upgrade of \texttt{CategoryTheory.GrpObj.mulRight} (\texttt{Mathlib.CategoryTheory.Monoidal.Grp\_.lean:277--281}, which packages a one-input automorphism \(a \mapsto a \cdot f\) for a fixed global element \(f \colon \mathbf{1} \to G\)): the map \(\sigma\) parametrises the family of right-translations \(\{\mathrm{mulRight}\,a\}_{a \in G}\) globally across \(G\) by taking the first coordinate as the parameter, and the second projection \(\pr_{2} \circ \sigma = \mu\) realises the right-translation action on the second coordinate. The "globalises the cotangent" phrasing in the name records that this family-of-translations is what triviallises \(\Omega_{G/k}\) via pullback along \(\sigma\) (cf.\ the inverse-pair proof shape of \texttt{CategoryTheory.GrpObj.isPullback}, \texttt{Grp\_.lean:293--323}, which uses the structurally identical \texttt{lift \dots (lift \dots \,\iota\, \dots) ≫ μ} construction). No standalone \texttt{mulRight}-name function appears in the Lean signature itself; the connection is purely nomenclatural, and the underlying Lean construction is the binary-product shear iso, not a one-input \texttt{mulRight} automorphism.
\end{lemma}

\begin{proof}
  


  \emph{Step 1: Build the shear iso \(\sigma\).} By \cref{lem:GrpObj_shearMulRight}, the binary-product shear \(\sigma_{\mathrm{hom}} = \langle \pr_{1}, \mu \rangle \colon G \times_{k} G \to G \times_{k} G\) in \(\Over\,(\Spec k)\) — informally \((a, b) \mapsto (a, a \cdot b)\) — is a self-iso of \(G \times_{k} G\) with explicit inverse \(\sigma_{\mathrm{inv}} = \langle \pr_{1}, \mu \circ (\iota \circ \pr_{1}, \pr_{2}) \rangle\), and satisfies the companion identifications \(\sigma_{\mathrm{hom}} \circ \pr_{1} = \pr_{1}\) and \(\sigma_{\mathrm{hom}} \circ \pr_{2} = \mu\). The Lean target is \texttt{AlgebraicGeometry.GrpObj.shearMulRight} (closed iter-134, \(\sim 30\) LOC of \texttt{lift\_lift\_assoc} / \texttt{lift\_comp\_inv\_*} / \texttt{lift\_comp\_one\_left} calculus on the \(\GrpObj\) axioms, modelled on \texttt{CategoryTheory.GrpObj.mulRight} and \texttt{CategoryTheory.GrpObj.isPullback} of \texttt{Mathlib.CategoryTheory.Monoidal.Grp\_.lean:277--281,\,293--323}).

  \emph{Step 2: Base-change of differentials.} View \(G \times_{k} G\) as a \(G\)-scheme via the first projection \(\pr_{1}\). The relative cotangent of this \(G\)-scheme, \(\Omega_{(G \times_{k} G)/G}\), is canonically isomorphic to the pullback \(\pr_{2}^{*}\,\Omega_{G/k}\) via \cref{lem:GrpObj_omega_basechange_proj} (the load-bearing helper sub-lemma, NEEDS\_MATHLIB\_GAP\_FILL, \(\sim 150\)--\(300\) LOC):
  \[
    \Omega_{(G \times_{k} G)/G} \;\cong\; \pr_{2}^{*}\,\Omega_{G/k} \qquad \text{on } G \times_{k} G.
  \]
  Since \(\sigma_{\mathrm{hom}}\) is an iso over \(\pr_{1}\) (i.e.\ \(\pr_{1} \circ \sigma_{\mathrm{hom}} = \pr_{1}\), by construction of \(\sigma_{\mathrm{hom}} = \langle \pr_{1}, \mu \rangle\) via the binary-product universal property), \(\sigma\) acts as a self-iso of \(G \times_{k} G\) as a \(G\)-scheme via \(\pr_{1}\); functoriality of \(\Omega_{(\cdot)/G}\) under such self-isos then re-routes the Step 2 iso through \(\sigma\) to relate the differentials at the two projections, with \(\pr_{2} \circ \sigma_{\mathrm{hom}} = \mu\).

  \emph{Step 3: Restrict along the section \(s := \langle \mathrm{id}_{G},\, \eta_{G} \rangle \colon G \to G \times_{k} G\).} The section \(s\) satisfies \(\pr_{1} \circ s = \mathrm{id}_{G}\) and \(\pr_{2} \circ s = \eta_{G} \circ \pi_{G}\) in \(\Over\,(\Spec k)\), where \(\pi_{G} \colon G \to \Spec k\) is the structure map. Applying \(s^{*}\) to both sides of the Step 2 iso \(\Omega_{(G \times_{k} G)/G} \cong \pr_{2}^{*}\,\Omega_{G/k}\):
  \begin{itemize}
    \item The LHS \(s^{*}(\Omega_{(G \times_{k} G)/G})\) is, by the \(G\)-scheme structure given by \(\pr_{1}\) and the section identity \(\pr_{1} \circ s = \mathrm{id}_{G}\), canonically isomorphic to \(\Omega_{G/k}\) (restricting the relative cotangent of \(G \times_{k} G\) over \(G\) to a section of the \(G\)-projection recovers \(\Omega_{G/k}\) — a standard relative-cotangent identification for a section of a smooth morphism). Mathlib-wise this combines \texttt{PresheafOfModules.pullbackComp} with \texttt{PresheafOfModules.pullbackId} on the composition \(\pr_{1} \circ s = \mathrm{id}_{G}\).
    \item The RHS \(s^{*}(\pr_{2}^{*}\,\Omega_{G/k})\) is, by \cref{lem:GrpObj_omega_restrict_to_identity_section} applied to the section identity \(\pr_{2} \circ s = \eta_{G} \circ \pi_{G}\), canonically isomorphic to \(\pi_{G}^{*}(\eta_{G}^{*}\,\Omega_{G/k})\).
  \end{itemize}
  Composing the LHS and RHS identifications with the pulled-back Step 2 iso yields the displayed iso \(\Omega_{G/k} \cong \pi_{G}^{*}(\eta_{G}^{*}\,\Omega_{G/k})\). Estimated build: \(\sim 30\)--\(80\) LOC (per the iter-133 mathlib-analogist verdict Decision 4, step 3 of the LOC envelope). The role of \(\sigma\) in this proof is the Step 2 functorial book-keeping — it pins the "right-translation" interpretation of the trivialisation (the second-coordinate cotangent \(\pr_{2}^{*}\,\Omega_{G/k}\) is identified with the first-projection cotangent \(\Omega_{(G \times_{k} G)/G}\) via Step 2's helper, and \(\sigma\) realises the consistency of these two views as a single \(G\)-scheme structure on \(G \times_{k} G\)).

  \emph{Mathlib name summary.} Step 1 (shear iso): \cref{lem:GrpObj_shearMulRight} (Lean target \texttt{AlgebraicGeometry.GrpObj.shearMulRight}; closed iter-134 via the \texttt{lift\_*} calculus on \texttt{CategoryTheory.GrpObj.mulRight} and \texttt{CategoryTheory.GrpObj.isPullback} of \texttt{Mathlib.CategoryTheory.Monoidal.Grp\_.lean:277--281,\,293--323}). Step 2 (base-change of differentials): \cref{lem:GrpObj_omega_basechange_proj} (chains \texttt{KaehlerDifferential.tensorKaehlerEquiv} from \texttt{Mathlib.RingTheory.Kaehler.TensorProduct} with the presheaf-side \texttt{TopCat.Presheaf.pullback} from \texttt{Mathlib.Topology.Sheaves.Pullback}, plus the project's \texttt{relativeDifferentialsPresheaf\_obj\_kaehler}; NEEDS\_MATHLIB\_GAP\_FILL, \(\sim 150\)--\(300\) LOC). Step 3 (section restriction): \cref{lem:GrpObj_omega_restrict_to_identity_section} (applies \texttt{PresheafOfModules.pullbackComp} and \texttt{PresheafOfModules.pullbackId} from \texttt{Mathlib.Algebra.Category.ModuleCat.Presheaf.Pullback}; \(\sim 30\)--\(80\) LOC).

  \emph{Total LOC envelope for piece (i.b) under the recommended sheaf-level RHS}: \(\sim 210\)--\(440\) LOC (per the iter-133 mathlib-analogist verdict). The bridge to the value-level \(\texttt{cotangentSpaceAtIdentity}\,G\) — a chart-localisation identification of \(\sim 100\)--\(200\) LOC — is a structurally distinct artefact consumed downstream in piece (i.c) (\cref{lem:GrpObj_omega_free}), \textbf{not} inside this lemma's body.
\end{proof}

\begin{definition}\leanok
[(i.b-packaging) Adjunction-transpose compatibility of structure presheaves of rings along a scheme morphism]
  \label{def:GrpObj_schemeHomRingCompatibility}
  \lean{AlgebraicGeometry.GrpObj.schemeHomRingCompatibility}
  
  Let \(f \colon Y \to Z\) be a morphism of schemes. The \emph{adjunction-transpose compatibility morphism along \(f\)} packages the canonical morphism of presheaves of commutative rings on \(Y\)
  \[
    \varphi_f \;\colon\; f_{\mathrm{top}}^{-1}\,\mathcal O_Z \;\longrightarrow\; \mathcal O_Y,
  \]
  obtained as the image of the structure-sheaf morphism \(f^{\sharp} \colon \mathcal O_Z \to (f_{\mathrm{top}})_{*}\,\mathcal O_Y\) under the symmetric direction of the presheaf-pullback / pushforward adjunction \(f_{\mathrm{top}}^{-1} \dashv (f_{\mathrm{top}})_{*}\) on \(\mathrm{CommRingCat}\)-valued presheaves of the underlying topological map \(f_{\mathrm{top}}\). This is the precise compatibility-of-rings shape consumed by the project's \texttt{relativeDifferentialsPresheaf} on Lean targets whose ring-side input comes from a scheme morphism: the \texttt{Differentials.lean:52} construction of \(\varphi'\) from \(f.c\) via the same adjunction transpose is the running model. The helper is used in piece (i.b) wherever a scheme morphism's structure-presheaf data must be fed into \texttt{relativeDifferentialsPresheaf} in this transposed shape.
\end{definition}

\begin{remark}[\(\varphi_f\) vs the \texttt{PresheafOfModules.pullback} compatibility morphism]
  \label{rem:GrpObj_schemeHomRingCompatibility_vs_toRingCatSheafHom}
  
  The compatibility morphism \(\varphi_f\) of \cref{def:GrpObj_schemeHomRingCompatibility} is \emph{structurally distinct} from the compatibility morphism required by \texttt{PresheafOfModules.pullback}, which expects a morphism of presheaves of rings on the \emph{codomain} \(Y\) (informally \(\mathcal O_Y \to (f_{\mathrm{top}})_{*} \mathcal O_Z\), in the homEquiv shape inverse to \(\varphi_f\)) obtained via Mathlib's \texttt{Scheme.Hom.toRingCatSheafHom}. Per the iter-135 mathlib-analogist verdict (\texttt{analogies/phi-compatibility-morphisms.md}), the canonical Mathlib idiom for the latter is \texttt{(Scheme.Hom.toRingCatSheafHom \(f\)).hom}. The two conventions serve distinct downstream consumers inside piece (i.b): \(\varphi_f\) feeds \texttt{relativeDifferentialsPresheaf} on the source, while \texttt{Scheme.Hom.toRingCatSheafHom} feeds \texttt{PresheafOfModules.pullback} on a pulled-back module-presheaf. Both shapes appear in the closure of \cref{lem:GrpObj_omega_basechange_proj} and \cref{lem:GrpObj_omega_restrict_to_identity_section} below.
\end{remark}

\paragraph{Iter-145 EXCISE disposition (block kept as traceable iter-141 \(\to\) iter-143 decision history).}
The Lean target \texttt{relativeDifferentialsPresheaf\_basechange\_along\_proj\_two}, together with the four sister declarations \texttt{basechange\_along\_proj\_two\_inv\_derivation} (\cref{lem:GrpObj_omega_basechange_proj_inv_derivation}), \texttt{basechange\_along\_proj\_two\_inv} (\cref{lem:GrpObj_omega_basechange_proj_inv}), \texttt{basechange\_along\_proj\_two\_inv\_app\_isIso} (\cref{lem:GrpObj_basechange_along_proj_two_inv_app_isIso}), and \texttt{mulRight\_globalises\_cotangent} (\cref{lem:GrpObj_mulRight_globalises}), were \textbf{EXCISED from \texttt{AlgebraicJacobian/Cotangent/GrpObj.lean} iter-145} as part of the iter-144 chart-algebra pivot (the bundled-route piece (i.b) is DESCOPED from critical-path; the iter-145 prover lane removed the sorry-bodied scaffolds rather than keeping them as auditable record). The blueprint blocks below are kept verbatim for traceability of the iter-141 \(\to\) iter-143 Route (b'2) closure programme; the original \texttt{\textbackslash lean\{\ldots\}} hints are iter-146-stripped because they pointed at declarations that no longer exist in tree. Should a future strategy reversal re-prioritise the bundled route, these blocks supply the prover-ready recipe at the level of detail recorded in their proof bodies.

\begin{lemma}[(i.b-helper) Base-change of differentials along the second projection]
  \label{lem:GrpObj_omega_basechange_proj}
  \lean{AlgebraicGeometry.TODO.GrpObj_omega_basechange_proj}
  \uses{lem:relative_kaehler_presheaf_obj, lem:GrpObj_omega_basechange_proj_inv, lem:GrpObj_basechange_along_proj_two_inv_app_isIso}
  % Iter-146 disposition: \lean{...relativeDifferentialsPresheaf_basechange_along_proj_two}
  % stripped — Lean target EXCISED iter-145 (see EXCISE disposition
  % paragraph immediately above this block). Block kept as auditable
  % bundled-route history.
  % Lean signature stub (iter-134+ NEEDS_MATHLIB_GAP_FILL helper for
  % piece (i.b); the load-bearing piece of (i.b) per the iter-133
  % mathlib-analogist verdict Decision 2 at
  % `task_results/mathlib-analogist-mulright-globalises-iter133.md`,
  % ~150-300 LOC). Chains the algebra-side
  % `KaehlerDifferential.tensorKaehlerEquiv` with the presheaf-side
  % `TopCat.Presheaf.pullback` to upgrade the algebra-side base-change-of-Ω
  % iso to a sheaf-level natural iso on `(G ×_k G).left`:
  %   noncomputable def relativeDifferentialsPresheaf_basechange_along_proj_two
  %       (G : Over (Spec (.of k))) [CategoryTheory.GrpObj G] :
  %       relativeDifferentialsPresheaf
  %           (CategoryTheory.Limits.prod.fst : (G ⊗ G).left ⟶ G.left).⋯ ≅
  %         (PresheafOfModules.pullback (φ_pr_two G)).obj
  %           (relativeDifferentialsPresheaf G.hom)
  % where `(G ⊗ G).left` is the underlying scheme of the categorical product
  % `G ⊗ G` in `Over (Spec (.of k))` (equivalently, the fibre product
  % `G.left ×_{Spec k} G.left`), the LHS uses the first projection
  % `pr_1 : G ⊗ G ⟶ G` as the structure morphism for the relative cotangent
  % of `G ⊗ G` viewed as a `G`-scheme, and `φ_pr_two` is the compatibility
  % morphism for the second projection `pr_2 : G ⊗ G ⟶ G` in the same shape
  % as `Differentials.lean:52`'s `φ'`. Note: the LHS index `pr_1`-vs.-`pr_2`
  % asymmetry follows from the binary-product universal property in
  % `Over (Spec k)`: the relative cotangent for `pr_1` is dual to pulling
  % back the cotangent of `G/Spec k` along `pr_2`, via
  % `tensorKaehlerEquiv`'s `Algebra.IsPushout` square.
  
  % NOTE iter-135: the Lean target
  % \lean{AlgebraicGeometry.GrpObj.relativeDifferentialsPresheaf_basechange_along_proj_two}
  % ships as an honest sorry-bodied scaffold; the signature is the intended
  % sheaf-level RHS displayed above (the compatibility morphism for
  % `PresheafOfModules.pullback` along the second projection is constructed
  % via Mathlib's `Scheme.Hom.toRingCatSheafHom`, per the iter-135
  % mathlib-analogist verdict in `analogies/phi-compatibility-morphisms.md`).
  % The proof body is `sorry`; this is the load-bearing \(\sim 150\)--\(300\) LOC
  % NEEDS_MATHLIB_GAP_FILL helper of piece (i.b). `sync_leanok` removes the
  % prior iter-134 `\leanok` from the proof block accordingly. Body closure
  % is iter-136+ work.
  % NOTE iter-140: \notready stripped — Lean target
  % `\lean{AlgebraicGeometry.GrpObj.relativeDifferentialsPresheaf_basechange_along_proj_two}`
  % is formalized at `Cotangent/GrpObj.lean:670` with a partial body
  % (Route (b'2) skeleton; per-open IsIso sub-sorry at L689). Per
  % CLAUDE.md marker vocab, `\notready` on a formalized statement is stale.
  % NOTE iter-146: the Lean target was subsequently EXCISED iter-145
  % from `AlgebraicJacobian/Cotangent/GrpObj.lean` under the iter-144
  % chart-algebra pivot (see EXCISE disposition paragraph at the head
  % of this block). The statement is therefore once again blueprint-only.
  \notready
  Let \(G : \Over\,(\Spec k)\) be as above, and consider the binary product \(G \times_{k} G\) in \(\Over\,(\Spec k)\), with first and second projections \(\pr_{1}, \pr_{2} \colon G \times_{k} G \to G\). View \(G \times_{k} G\) as a \(G\)-scheme via \(\pr_{1}\). Then there is a canonical natural isomorphism of presheaves of modules on \(G \times_{k} G\)
  \[
    \Omega_{(G \times_{k} G)/G} \;\cong\; \pr_{2}^{*}\,\Omega_{G/k},
  \]
  where the left-hand side is the relative cotangent of \(G \times_{k} G\) viewed as a \(G\)-scheme via \(\pr_{1}\), and the right-hand side is the pullback along \(\pr_{2}\) of \(\Omega_{G/k} = \texttt{relativeDifferentialsPresheaf}\,G.\mathrm{hom}\).
\end{lemma}

\begin{proof}
  
  % NOTE iter-139: the  marker on this proof block (immediately
  % below) may be a sync_leanok mis-mark. The Lean target
  % \lean{AlgebraicGeometry.GrpObj.relativeDifferentialsPresheaf_basechange_along_proj_two}
  % at \texttt{AlgebraicJacobian/Cotangent/GrpObj.lean:612} contains a
  % \texttt{letI : IsIso ... := sorry} pattern at line 624 (plus two
  % further sub-sorries inside the helper
  % \texttt{basechange\_along\_proj\_two\_inv\_derivation} at lines 581
  % and 585; see the iter-138 NOTE block below for the breakdown).
  % \texttt{sync\_leanok} may have mis-handled the \texttt{letI ... := sorry}
  % construction (a \texttt{let}-bound \texttt{sorry} inside a pure-term
  % definition body, distinct from a top-level \texttt{by sorry}). This is
  % flagged for the iter-139 plan agent's \texttt{doctor}-skill consult;
  % no agent edit to the \texttt{} line itself per the project
  % marker vocabulary (\texttt{} is \texttt{sync\_leanok}'s territory).
  %
  % \emph{Iter-141 update (pattern-citation only).} Iter-140 moved
  % the \texttt{letI : IsIso ... := sorry} pattern to a
  % \texttt{(fun \_ => sorry)} argument inside the iter-140-landed
  % helper \texttt{isIso\_of\_app\_iso\_module} (at
  % \texttt{Cotangent/GrpObj.lean:689}); the line/pattern citations
  % above are therefore dated. The underlying \texttt{sync\_leanok}
  % mis-mark concern persists — a sorry is still present in the
  % proof body, so the \texttt{} on this proof block remains
  % a mis-mark candidate until iter-142+ closes the per-open
  % \texttt{IsIso} sub-sorry.
  % NOTE iter-138: this proof block records the iter-138 prover-lane
  % closure shape for the Lean target
  % \lean{AlgebraicGeometry.GrpObj.relativeDifferentialsPresheaf_basechange_along_proj_two}.
  % The iter-138 prover lane chose \textbf{Route (b)
  % inverse-direction-via-adjunction-transpose} over the Route (a)
  % chart-unfolding-helper of the iter-137 5-step recipe (see the
  % iter-137 NOTE block immediately below), for the structural reason
  % recorded in the iter-138 lean-vs-blueprint-checker: Route (a)'s
  % chart-unfolding helper hits the same \texttt{pullback}-opacity
  % blocker as iter-137 (Mathlib's
  % \texttt{PresheafOfModules.pullback} is defined as
  % \texttt{(pushforward \(\varphi\)).leftAdjoint} with no chart-wise
  % unfolding lemma on \texttt{.obj}/\texttt{.map}), whereas Route (b)
  % admits typeable derivation construction without unfolding
  % \texttt{pullback}. The pushforward functor is transparent on
  % \texttt{.obj}/\texttt{.map} (see
  % \texttt{Mathlib/Algebra/Category/ModuleCat/Presheaf/Pushforward.lean:39, 86}),
  % so building the inverse direction's derivation pointwise via
  % \texttt{Derivation'.mk} + \texttt{ModuleCat.Derivation.mk} is
  % typeable end-to-end.
  %
  % Iter-138 introduced two top-level Lean helpers in
  % \texttt{AlgebraicJacobian/Cotangent/GrpObj.lean}, both pinned in
  % \cref{lem:GrpObj_omega_basechange_proj_inv_derivation} and
  % \cref{lem:GrpObj_omega_basechange_proj_inv} below:
  % \begin{itemize}
  %   \item \texttt{basechange\_along\_proj\_two\_inv\_derivation}
  %     (\texttt{Cotangent/GrpObj.lean:547}), the pointwise inverse-
  %     direction derivation
  %     \(D : ((\texttt{pushforward}\,\psi).\mathrm{obj}\,\Omega_{(G \times_k G)/G}).\texttt{Derivation'}\,\varphi_G\)
  %     whose pointwise rule at each open \(X\) is
  %     \(d_X(b) := \texttt{KaehlerDifferential.D}\,((\psi.\app\,X).\hom\,b)\).
  %   \item \texttt{basechange\_along\_proj\_two\_inv}
  %     (\texttt{Cotangent/GrpObj.lean:596}), the inverse morphism
  %     obtained by \texttt{(DifferentialsConstruction.isUniversal'\,\(\varphi_G\)).desc}
  %     on \(D\), then transposed along
  %     \texttt{(pullbackPushforwardAdjunction\,\(\psi\)).homEquiv.symm}.
  % \end{itemize}
  %
  % \textbf{Closed iter-138.} The additive (\texttt{d\_add}) and
  % Leibniz (\texttt{d\_mul}) laws of the pointwise derivation are
  % closed honestly via \texttt{RingHom.map\_add} / \texttt{RingHom.map\_mul}
  % applied to \texttt{(\(\psi\).app\,X).hom} composed with
  % \texttt{ModuleCat.Derivation.d\_add} / \texttt{d\_mul} on the
  % algebra-side \texttt{KaehlerDifferential.D}. The iso itself is
  % then assembled as \texttt{(asIso\,(basechange\_along\_proj\_two\_inv\,G)).symm}
  % once the iso property of the inverse is supplied.
  %
  % \textbf{Sub-sorry status (iter-143 status update; supersedes the
  % earlier iter-138 "three sub-sorries remain (iter-140 prover-lane
  % targets)" framing).}
  % \begin{enumerate}
  %   \item \texttt{d\_map}: \textbf{CLOSED iter-142 substantively}
  %     in \texttt{AlgebraicJacobian/Cotangent/GrpObj.lean:638--674}
  %     via the 3-step ALIGN\_WITH\_MATHLIB chase per
  %     \texttt{PresheafOfModules.pushforward\_obj\_map\_apply'} +
  %     \texttt{NatTrans.naturality\_apply} +
  %     \texttt{relativeDifferentials'\_map\_d} (the iter-141
  %     mathlib-analogist's prescribed pattern). The cross-open
  %     naturality square of the pointwise derivation is now a
  %     theorem with kernel-only axioms. See the \emph{d\_map closure
  %     recipe} iter-138 NOTE block immediately below for the
  %     historical recipe; the iter-142 closure followed it with the
  %     two empirical refinements codified in the iter-142 Rules 1--3
  %     NOTE above.
  %   \item \texttt{d\_app}: \textbf{PARTIAL iter-143}. The iter-143
  %     prover lane closed Step 3.a (the categorical equality
  %     $(\texttt{fst}\,G\,G).\mathrm{left} \circ G.\mathrm{hom} =
  %        (\texttt{snd}\,G\,G).\mathrm{left} \circ G.\mathrm{hom}$ in
  %     \(\Over\,(\Spec k)\)) via the 1-LOC
  %     \texttt{have hw := by rw [(fst G G).w, (snd G G).w]}
  %     derivation, but the residual chase 3.a \(\to\) 3.b
  %     (\texttt{PresheafedSpace.comp\_c\_app} lift) \(\to\) 3.c
  %     (\texttt{Adjunction.homEquiv\_naturality\_right\_symm}
  %     transpose) \(\to\) 3.d (\texttt{ModuleCat.Derivation.d\_map}
  %     discharge) is blocked at the \texttt{Pushforward.comp\_eq} +
  %     \texttt{eqToHom} type-coercion layer (see Rule 4 of the
  %     iter-143 empirical lessons NOTE below, just before the
  %     Step 3 sub-recipe block). Residual \texttt{sorry} at
  %     \texttt{Cotangent/GrpObj.lean} inside the
  %     \texttt{basechange\_along\_proj\_two\_inv\_derivation} body.
  %   \item \texttt{IsIso\,(basechange\_along\_proj\_two\_inv\,G)}:
  %     \textbf{Refactored iter-143} into a named sorry-bodied
  %     theorem
  %     \texttt{AlgebraicGeometry.GrpObj.basechange\_along\_proj\_two\_inv\_app\_isIso}
  %     at \texttt{Cotangent/GrpObj.lean:745--751} per the
  %     \texttt{lean-auditor-review142} MAJOR finding and the
  %     iter-143 \texttt{strategy-critic-iter143} sorry-must-be-
  %     named-declaration discipline. Body remains \texttt{sorry};
  %     the new named theorem is pinned in the blueprint as
  %     \cref{lem:GrpObj_basechange_along_proj_two_inv_app_isIso}
  %     adjacent to \cref{lem:GrpObj_omega_basechange_proj_inv} below.
  % \end{enumerate}
  %
  % \textbf{Iter-144 chart-algebra pivot disposition.} Under the
  % iter-144 chart-algebra pivot (see the \S "Iter-144 chart-algebra
  % pivot --- disposition" paragraph at the head of
  % \cref{subsec:RigidityKbar_piece_i_decomposition}), the residual
  % d\_app sub-sorry and the named IsIso theorem are both
  % \textbf{DESCOPED from critical-path}: their named Lean
  % declarations remain in-tree as auditable record of the bundled
  % Route (b) / Route (b'2) closure programme, but the iter-145+
  % prover lane targets the chart-algebra (\(\alpha\)) + (\(\beta\))
  % helpers of piece (ii) instead (see the \S "Iter-144 chart-algebra
  % envelope for piece (ii)" paragraph at the head of the same
  % subsection).
  %
  % \textbf{Iter-138 negative lesson (codified in
  % \texttt{.archon/.debug-feedback/debug\_feedback.md}).} Inside the
  % goals produced by \texttt{Derivation.mk}, \texttt{simp} with the
  % obvious lemma names (\texttt{map\_add}, \texttt{map\_mul},
  % \texttt{ModuleCat.Derivation.d\_add},
  % \texttt{ModuleCat.Derivation.d\_mul}) does \emph{not} fire because
  % the function passed to \texttt{Derivation.mk} appears as a
  % beta-redex (the lambda is not yet reduced when \texttt{simp}
  % attempts to match the redex's head). The working pattern is
  % the \texttt{have h : ... := ... ; change ... ; rw [h] ; exact ...}
  % four-step sequence: extract the equation as an explicit hypothesis,
  % \texttt{change} the goal into a form whose head matches that
  % hypothesis's left-hand side, rewrite by the hypothesis, and close
  % with the derivation law. Iter-138's closed \texttt{d\_add} and
  % \texttt{d\_mul} bodies (\texttt{Cotangent/GrpObj.lean:560--575})
  % exhibit this pattern verbatim; future iter-140 attacks on
  % \texttt{d\_app} and \texttt{d\_map} should mimic it rather than
  % attempt a one-shot \texttt{simp}.
  %
  % NOTE iter-138 (d\_app closure recipe). Let
  % \(\varphi_G : (G.\mathrm{hom}.\mathrm{base}^{-1}\,\mathcal O_{\Spec k}) \to G.\mathrm{left}.\texttt{presheaf}\)
  % be the compatibility morphism for the structure map
  % \(G.\mathrm{hom} : G.\mathrm{left} \to \Spec k\) (the iter-135
  % \texttt{Scheme.Hom.toRingCatSheafHom} idiom, in homEquiv-transpose
  % shape for the source-side of the pullback adjunction), and
  % \(\psi : G.\mathrm{left}.\texttt{presheaf} \to (G \otimes G).\mathrm{left}.\texttt{presheaf} \circ \mathrm{snd}^{-1}\)
  % the compatibility morphism for the second projection
  % \(\pr_2 = \texttt{snd}\,G\,G : G \otimes G \to G\). The
  % \texttt{d\_app} law of the pointwise derivation requires
  % showing, for each open \(X \in G.\mathrm{left}.\Opens^{\mathrm{op}}\)
  % and each section \(a \in (G.\mathrm{hom}.\mathrm{base}^{-1}\,\mathcal O_{\Spec k}).\mathrm{obj}\,X\),
  % the vanishing
  % \(\texttt{KaehlerDifferential.D}\,\_\,((\psi.\app\,X).\hom\,(\varphi_G.\app\,X\,a)) = 0\).
  %
  % This vanishing follows from the \textbf{categorical commutativity in
  % \(\Over\,(\Spec k)\)}: the first and second projections
  % \(\pr_1 = \texttt{fst}\,G\,G\), \(\pr_2 = \texttt{snd}\,G\,G\) of
  % \(G \otimes G\) satisfy
  % \(\pr_1.\mathrm{left} \circ G.\mathrm{hom} = \pr_2.\mathrm{left} \circ G.\mathrm{hom}\)
  % as morphisms \(G \otimes G \to \Spec k\) (both compositions equal
  % the structure map of \(G \otimes G\) in \(\Over\,(\Spec k)\) — a
  % direct consequence of the binary-product universal property in
  % \(\Over\,(\Spec k)\), which forces both projection morphisms to be
  % morphisms over \(\Spec k\) and hence to commute with the structure
  % maps). Translating this scheme-level identity to the presheaf
  % level via Mathlib's \texttt{Scheme.Hom.toRingCatSheafHom}
  % functorial: the composite \((\psi.\app\,X) \circ (\varphi_G.\app\,X)\)
  % factors through the source-presheaf morphism
  % \(\varphi_{\texttt{LHS}}.\app\,(\mathrm{snd}^{-1}\,X)\), where
  % \(\varphi_{\texttt{LHS}}\) is the compatibility morphism for the
  % structure map of \((G \otimes G)\) as a \(G\)-scheme via \(\pr_1\) (the
  % \texttt{LHS} of the displayed iso, i.e.\
  % \(\texttt{relativeDifferentialsPresheaf}\,(\texttt{fst}\,G\,G).\mathrm{left}\)).
  % By construction of the pointwise derivation (Mathlib's
  % \texttt{KaehlerDifferential.D}), the universal Kähler differential
  % vanishes on the image of the source-presheaf morphism — this is
  % the algebra-side \texttt{d\_app} law of \texttt{KaehlerDifferential.D}
  % packaged via \texttt{ModuleCat.Derivation.d\_app} at
  % \texttt{Mathlib/Algebra/Category/ModuleCat/Differentials/Basic.lean}.
  %
  % \textbf{Concretely}, the Lean closure should: (1) chase the
  % scheme-level equality
  % \(\pr_1.\mathrm{left} \circ G.\mathrm{hom} = \pr_2.\mathrm{left} \circ G.\mathrm{hom}\)
  % using the \(\Over\,(\Spec k)\) structural identities
  % (\texttt{Over.comp\_left}, \texttt{Over.toUnit\_left}, and the
  % universal property of the binary product \(G \otimes G\) packaged
  % by \texttt{CartesianMonoidalCategory.hom\_ext} on the
  % \(\texttt{toUnit}\colon G \otimes G \to \mathbf 1 = \Spec k\)
  % morphism); (2) translate via
  % \texttt{Scheme.Hom.toRingCatSheafHom}'s functorial action to the
  % presheaf-of-rings level on opens of \(G.\mathrm{left}\), obtaining
  % the factorisation
  % \((\psi.\app\,X) \circ (\varphi_G.\app\,X) = \varphi_{\texttt{LHS}}.\app\,(\mathrm{snd}^{-1}\,X) \circ (\text{source-side ring map})\);
  % (3) apply the algebra-side
  % \texttt{ModuleCat.Derivation.d\_app} (or the underlying
  % \texttt{Derivation.d\_app} on
  % \texttt{KaehlerDifferential.D}) to discharge the vanishing.
  %
  % \emph{Implementation note (iter-141).} The algebra-side closing
  % lemma is \texttt{ModuleCat.Derivation.d\_map} at
  % \texttt{Mathlib/Algebra/Category/ModuleCat/Differentials/Basic.lean:80}
  % (the \texttt{@[simp]} lemma \(D.d\,(f\,a) = 0\) for any
  % \(\texttt{ModuleCat.Derivation}\,D\) over a ring map \(f\)).
  % \textbf{NOT} \texttt{Derivation.map\_algebraMap}; the
  % \texttt{ModuleCat.Derivation.d\_map} packages the
  % \texttt{Algebra}/\texttt{Module}/\texttt{IsScalarTower} instance
  % discharge + the \texttt{map\_algebraMap} call into a single
  % \texttt{@[simp]} lemma. The streamlined closure shape, verified
  % standalone this iter via \texttt{lean\_run\_code} (analogy
  % \texttt{d-app-d-map-recipe-shape}, Decision 1), is:
  % \begin{verbatim}
  %   example (A B C : CommRingCat) (f1 : A ⟶ B) (g : C ⟶ B) (k : A ⟶ C)
  %       (hcomm : k ≫ g = f1) (a : A) :
  %       (CommRingCat.KaehlerDifferential.D g).d (f1.hom a) = 0 := by
  %     rw [show f1.hom a = g.hom (k.hom a) from by rw [← hcomm]; rfl]
  %     exact (CommRingCat.KaehlerDifferential.D g).d_map _
  % \end{verbatim}
  % This saves \(\sim 4\) LOC per call site over the iter-140-validated
  % \texttt{letI : Algebra ... ; letI : Module ... ; letI : IsScalarTower ... ;
  % exact (D : Derivation ...).map\_algebraMap \_} chain by bundling the
  % instance discharge inside \texttt{d\_map}. \textbf{Combined LOC estimate
  % for the d\_app closure (iter-142 empirical band, validated on the
  % d\_map closure that shares the same opacity-sensitive
  % \texttt{pushforward\(_0\)}/transposed adjunction shape): \(\sim 40\)--\(80\) LOC}
  % (\(\sim 20\)--\(40\) LOC for the Step 3 categorical-witness chase
  % \(h : \text{Source} \to \text{Target}\) with
  % $h \circ \varphi_{\texttt{LHS}}.\app\,(\mathrm{snd}^{-1}\,X)
  % = \varphi_G.\app\,X \circ \psi.\app\,X$
  % via \((\texttt{fst}\,G\,G).w\), \((\texttt{snd}\,G\,G).w\),
  % \texttt{LocallyRingedSpace.comp\_c\_app}, and
  % \texttt{pullbackPushforwardAdjunction.homEquiv.symm};
  % plus \(\sim 5\)--\(10\) LOC for the algebra-side discharge via
  % \texttt{d\_map}; plus \(\sim 15\)--\(30\) LOC for the explicit
  % \texttt{change}-skeleton wiring per iter-142 Rule 1 below). The
  % earlier iter-141 envelope \(\sim 50\)--\(90\) LOC is superseded by the
  % iter-142 band as the d\_map closure at
  % \texttt{Cotangent/GrpObj.lean:638--674} landed in the lower portion
  % of the predicted band.
  %
  % \textbf{NOTE iter-142 (empirical lessons; d\_map closure
  % transferable to d\_app).} The iter-142 prover lane closed the
  % d\_map sub-sorry substantively (the first strict-count closure on
  % the piece (i.b) route since iter-138; see
  % \texttt{task\_results/Cotangent\_GrpObj.lean.md} iter-142). The
  % closure pattern surfaced three empirical rules that transfer
  % directly to the d\_app closure, because the d\_app categorical
  % chase crosses the \emph{same} \texttt{pushforward\(_0\)}-annotated
  % definitions and the \emph{same}
  % \texttt{pullbackPushforwardAdjunction.homEquiv.symm}-transposed
  % term as d\_map. These rules supersede the iter-138/140 negative
  % lesson framing and the iter-141 \texttt{simp only [pushforward\_obj\_map\_apply']}
  % advisory:
  %
  % \emph{Rule 1 (fully-explicit \texttt{change}).} Any \texttt{change}
  % block whose goal crosses a \texttt{pushforward\(_0\)}-annotated
  % definition must spell both LHS and RHS \emph{fully explicitly};
  % \texttt{\_}-placeholders on either side of an equation crossing
  % \texttt{pushforward\(_0\)} trigger a deterministic \texttt{whnf}
  % timeout at \texttt{maxHeartbeats=200000} (the iter-140 d\_map
  % attempt; the iter-142 d\_map fix is to spell out both sides). The
  % rule applies regardless of whether the equation is
  % \texttt{\_ = 0}-shape or \texttt{\_ = \_}-shape: it is the
  % \texttt{set\_option backward.isDefEq.respectTransparency false}
  % annotation on \texttt{pushforward\(_0\)}
  % (\texttt{Mathlib/Algebra/Category/ModuleCat/Presheaf/Pushforward.lean:37, 55})
  % that blocks \texttt{whnf}/\texttt{change} from descending into the
  % placeholder when elaborating the target shape. Concretely for
  % d\_app, the explicit \texttt{change} must spell the
  % \texttt{(CommRingCat.KaehlerDifferential.D \_).d ((\psi.app\,X).hom\,((\varphi_G.app\,X).hom\,a)) = 0}
  % form rather than a \texttt{\_ = 0} placeholder. The iter-142 d\_app
  % skeleton already lands this explicit \texttt{change} at
  % \texttt{Cotangent/GrpObj.lean:611--618}; the residual gap is Step 3
  % (the witness construction below), not the \texttt{change}-skeleton.
  %
  % \emph{Rule 2 (\texttt{NatTrans.naturality\_apply} packaging).} A
  % bare \texttt{rw [NatTrans.naturality\_apply \(\psi\).hom f x]} fails
  % to match goals carrying \texttt{RingCat.Hom.hom}/\texttt{CommRingCat.Hom.hom}-form
  % terms, because the lemma produces an equation in
  % \texttt{ConcreteCategory.hom}-form whose kernel pattern does not
  % unify with the post-\texttt{change} \texttt{.hom}-form goal. The
  % working idiom is to package the naturality fact via
  % \texttt{rw [show \ldots = \ldots from NatTrans.naturality\_apply \ldots]},
  % forcing the rewrite pattern into the \texttt{.hom}-form that
  % matches the goal. Schematically, for d\_app's Step 3.b
  % \(\psi\)-naturality step:
  % \begin{verbatim}
  %   rw [show ((Scheme.Hom.toRingCatSheafHom (snd G G).left).hom.app Y).hom
  %             ((G.left.presheaf.map f).hom x)
  %         = ((G ⊗ G).left.presheaf.map ((Opens.map (snd G G).left.base).op.map f)).hom
  %             (((Scheme.Hom.toRingCatSheafHom (snd G G).left).hom.app X).hom x) from
  %         NatTrans.naturality_apply
  %           (Scheme.Hom.toRingCatSheafHom (snd G G).left).hom f x]
  % \end{verbatim}
  % The d\_app categorical chase also routes through \(\psi\)-naturality
  % (Step 3.b below), so this packaging is reusable verbatim modulo
  % renaming of \(f, x, X, Y\).
  %
  % \emph{Rule 3 (\texttt{pushforward\_obj\_map\_apply'} alternative).}
  % \texttt{PresheafOfModules.pushforward\_obj\_map\_apply'}
  % (\texttt{Mathlib/Algebra/Category/ModuleCat/Presheaf/Pushforward.lean:99--106},
  % located by the iter-141 mathlib-analogist
  % \texttt{d-app-d-map-recipe-shape}) is the named \texttt{@[simp]}
  % lemma for the kernel-form unfolding of
  % \texttt{((pushforward \(\psi\)).obj LHS).map f}. In the iter-142
  % d\_map closure the unfolding was achieved via the fully-explicit
  % \texttt{change} of Rule 1 (which obviates the need to invoke the
  % named lemma explicitly), but \texttt{pushforward\_obj\_map\_apply'}
  % remains the explicit alternative when the goal's surface form
  % does not align with a single \texttt{change} block. For d\_app,
  % both options are available; prefer the fully-explicit
  % \texttt{change} per Rule 1 when feasible.
  %
  % \emph{Rule 4 (iter-143 empirical lesson on
  % \texttt{Pushforward.comp\_eq} / \texttt{eqToHom} type-coercion).}
  % The iter-143 d\_app prover lane closed Step 3.a (the categorical
  % equality in \(\Over\,(\Spec k)\)) via the one-line
  % \texttt{have hw : (fst G G).left \(\circ\) G.hom = (snd G G).left \(\circ\) G.hom
  %    := by rw [(fst G G).w, (snd G G).w]}
  % derivation, but the residual chase 3.a \(\to\) 3.b
  % (\texttt{PresheafedSpace.comp\_c\_app} lift) \(\to\) 3.c
  % (\texttt{Adjunction.homEquiv\_naturality\_right\_symm} transpose)
  % \(\to\) 3.d (\texttt{ModuleCat.Derivation.d\_map} discharge) is
  % blocked at the \textbf{\texttt{Pushforward.comp\_eq} +
  % \texttt{eqToHom} type-coercion} layer. Specifically: identifying
  % $\texttt{pushforward}\,(\texttt{fst}).\mathrm{left}.\mathrm{base} \circ
  %    \texttt{pushforward}\,G.\mathrm{hom}.\mathrm{base}$
  % with \(\texttt{pushforward}\,(G \otimes G).\mathrm{hom}.\mathrm{base}\)
  % modulo \texttt{hw} is up-to-\texttt{rfl} at the functor level
  % (\texttt{TopCat.Presheaf.Pushforward.comp\_eq} is a definitional
  % \texttt{rfl} for \((f \circ g)_{*}\,\mathcal F = g_{*}\,(f_{*}\,\mathcal F)\)),
  % but applying this identification at a specific open \(X\) requires
  % explicit \texttt{eqToHom} insertions: the propositional equality
  % \texttt{hw} produces an \texttt{Eq.mpr}/\texttt{eqToHom} transport
  % when transposed through
  % \texttt{pullbackPushforwardAdjunction.homEquiv.symm}, and the
  % adjunction-transpose step does not collapse the \texttt{eqToHom}
  % insertion against \texttt{Pushforward.comp\_eq}'s definitional
  % \texttt{rfl}. The bespoke nat-trans chase through
  % \texttt{pullbackPushforwardAdjunction.homEquiv} must thread these
  % \texttt{eqToHom} rewrites alongside
  % \texttt{Adjunction.homEquiv\_naturality\_right\_symm}; the iter-143
  % prover lane could not fit this into its envelope. Under the
  % iter-144 chart-algebra pivot (see the \S "Iter-144 chart-algebra
  % pivot --- disposition" paragraph at the head of this
  % subsection), the Step 3 chase is \textbf{DESCOPED} (preserved
  % only as auditable record of the bundled route). Iter-145+ provers
  % should NOT re-attempt the bundled Step 3 chase; the chart-algebra
  % \(\alpha\) + \(\beta\) helpers (cf.\ piece (ii) chart-algebra
  % envelope) carry the M2.body-pile trajectory instead.
  %
  % \textbf{Step 3 (adjunction-transpose) sub-recipe (iter-142
  % decomposition; \(\sim 20\)--\(40\) LOC bespoke chase).} The iter-142
  % d\_app skeleton landed at \texttt{Cotangent/GrpObj.lean:611--618}
  % puts the goal in canonical form
  % \(\texttt{KD.d}\,((\psi.\app\,X).\hom\,((\varphi_G.\app\,X).\hom\,a)) = 0\);
  % the remaining work is Step 3 of the four-step recipe
  % (categorical-equality \(\to\) c-component lift \(\to\)
  % adjunction-transpose \(\to\) \texttt{d\_map} discharge). Sub-steps:
  % \begin{itemize}
  %   \item[3.a] \emph{Categorical equality in \(\Over\,(\Spec k)\)}
  %     (\(\sim 3\)--\(5\) LOC): from \((\texttt{fst}\,G\,G).w\) and
  %     \((\texttt{snd}\,G\,G).w\) (both equal \((G \otimes G).\mathrm{hom}\),
  %     i.e.\ the structure map of \(G \otimes G\) over \(\Spec k\)), derive
  %     $(\texttt{fst}\,G\,G).\mathrm{left} \circ G.\mathrm{hom} =
  %      (\texttt{snd}\,G\,G).\mathrm{left} \circ G.\mathrm{hom}$
  %     in the underlying \(\Scheme\) category.
  %   \item[3.b] \emph{Lift to c-components}
  %     (\(\sim 10\)--\(15\) LOC): apply
  %     \texttt{LocallyRingedSpace.comp\_c\_app} (or
  %     \texttt{PresheafedSpace.comp\_c\_app}, available transitively
  %     via \texttt{Mathlib.AlgebraicGeometry.PresheafedSpace}) to
  %     both sides of 3.a's identity to obtain the presheaf-of-rings
  %     equality
  %     $G.\mathrm{hom}.c.\app\,X \;\circ\; (\texttt{fst}\,G\,G).\mathrm{left}.c.\app\,(G.\mathrm{hom}^{-1}\,X)
  %       = G.\mathrm{hom}.c.\app\,X \;\circ\; (\texttt{snd}\,G\,G).\mathrm{left}.c.\app\,(G.\mathrm{hom}^{-1}\,X)$
  %     in \(\CommRingCat\). The \(\psi\)-naturality of Rule 2 above is
  %     invoked at this step, packaged via the
  %     \texttt{rw [show \ldots from NatTrans.naturality\_apply \ldots]}
  %     idiom.
  %   \item[3.c] \emph{Transpose through
  %     \texttt{pullbackPushforwardAdjunction.homEquiv.symm}}
  %     (\(\sim 5\)--\(15\) LOC of bespoke chase, no Mathlib shortcut).
  %     Construct the witness ring map
  %     \[
  %       h \;\colon\; ((\texttt{TopCat.Presheaf.pullback}\,G.\mathrm{hom}.\mathrm{base}).\obj\,(\Spec k).\texttt{presheaf}).\obj\,X
  %         \;\longrightarrow\;
  %         ((\texttt{TopCat.Presheaf.pullback}\,(\texttt{fst}\,G\,G).\mathrm{left}.\mathrm{base}).\obj\,G.\mathrm{left}.\texttt{presheaf}).\obj\,(\mathrm{snd}^{-1}\,X)
  %     \]
  %     satisfying
  %     $(\varphi_{\texttt{LHS}}.\app\,(\mathrm{snd}^{-1}\,X)) \circ h
  %       = (\psi.\app\,X) \circ (\varphi_G.\app\,X)$.
  %     Mathematically, this is the morphism between inverse-image
  %     presheaves induced by 3.b's c-component equality, presented
  %     concretely via the colimit description of
  %     \(\texttt{TopCat.Presheaf.pullback}\) (\texttt{Mathlib.Topology.Sheaves.Pullback}):
  %     a section of the source inverse-image presheaf at \(X\) is a
  %     colimit-class of sections of \((\Spec k).\texttt{presheaf}\) on
  %     opens containing \(G.\mathrm{hom}.\mathrm{base}\,X\), and 3.b's
  %     identity transports each such section to a coherent class of
  %     sections of \(G.\mathrm{left}.\texttt{presheaf}\) on opens
  %     containing \((\texttt{fst}\,G\,G).\mathrm{left}.\mathrm{base}\,(\mathrm{snd}^{-1}\,X)\).
  %     The verdicts of the iter-141 + iter-142 mathlib-analogists
  %     (\texttt{analogies/d-app-d-map-recipe-shape.md} Decision 2)
  %     are \texttt{NEEDS\_MATHLIB\_GAP\_FILL}: no Mathlib shortcut
  %     packages this transpose for arbitrary
  %     \texttt{pullbackPushforwardAdjunction}-transposed compatibility
  %     morphisms, so the witness \(h\) must be constructed term-mode in
  %     the body of the d\_app closure. This is the \textbf{load-bearing
  %     residual} of the d\_app closure; the iter-143 prover lane
  %     targets it as the single concrete sub-sorry remaining after
  %     the iter-142 \texttt{change}-skeleton (IsIso closure is
  %     deferred to iter-144+ as a separate prover round on the
  %     post-refactor named theorem; see Edit 2's NOTE in the IsIso
  %     block below).
  %   \item[3.d] \emph{\texttt{d\_map} discharge}
  %     (\(\sim 5\)--\(10\) LOC): with \(h\) in hand, the goal
  %     \(\texttt{KD.d}\,((\psi.\app\,X).\hom\,((\varphi_G.\app\,X).\hom\,a)) = 0\)
  %     rewrites via 3.c's factorisation
  %     $((\psi.\app\,X).\hom \circ (\varphi_G.\app\,X).\hom)\,a
  %       = (\varphi_{\texttt{LHS}}.\app\,(\mathrm{snd}^{-1}\,X)).\hom\,(h.\hom\,a)$
  %     to
  %     \(\texttt{KD.d}\,((\varphi_{\texttt{LHS}}.\app\,(\mathrm{snd}^{-1}\,X)).\hom\,(h.\hom\,a)) = 0\),
  %     which closes via the streamlined
  %     \(\texttt{(CommRingCat.KaehlerDifferential.D \_).d\_map \_}\)
  %     pattern (the \texttt{ModuleCat.Derivation.d\_map}
  %     \texttt{@[simp]} lemma at
  %     \texttt{Mathlib/Algebra/Category/ModuleCat/Differentials/Basic.lean:80}).
  %     A small \texttt{change} block may be needed to align
  %     argument shapes before the final \texttt{exact}; per Rule 1,
  %     this \texttt{change} must be fully explicit.
  % \end{itemize}
  %
  % \textbf{NOTE iter-142 status (d\_map sub-recipe block below):
  % CLOSED substantively in
  % \texttt{AlgebraicJacobian/Cotangent/GrpObj.lean:638--674}.} The
  % 3-step ALIGN\_WITH\_MATHLIB chase recipe in this block was
  % applied by the iter-142 prover lane with two empirical
  % refinements that are codified above in the iter-142 d\_app
  % empirical-lessons NOTE: (i) the explicit \texttt{change} must
  % spell both LHS \emph{and} RHS (not just LHS, contrary to the
  % iter-141 \texttt{simp only [pushforward\_obj\_map\_apply']}
  % framing which left the RHS implicit); (ii)
  % \texttt{NatTrans.naturality\_apply} is packaged via
  % \texttt{rw [show \ldots = \ldots from NatTrans.naturality\_apply \ldots]}
  % for kernel-form alignment. The iter-142 closure used the
  % fully-explicit \texttt{change} (Rule 1 above) to perform the
  % \texttt{((pushforward \(\psi\)).obj LHS).map f} unfolding
  % implicitly, without an explicit
  % \texttt{simp only [pushforward\_obj\_map\_apply']} step (Rule 3
  % above: both options are valid; the fully-explicit \texttt{change}
  % was preferred in iter-142). See the iter-142 d\_app empirical-
  % lessons NOTE above for the rule generalization; see
  % \texttt{task\_results/Cotangent\_GrpObj.lean.md} iter-142 for the
  % in-Lean closure details. Status is recorded here in-prose only;
  % \texttt{\textbackslash leanok} marker management is the
  % deterministic \texttt{sync\_leanok} phase's domain (CLAUDE.md
  % \S Blueprint Marker Vocabulary).
  %
  % NOTE iter-138 (d\_map closure recipe). The \texttt{d\_map} law of
  % the pointwise derivation requires showing, for each morphism
  % \(f : X \to Y\) of opens in \(G.\mathrm{left}.\Opens^{\mathrm{op}}\)
  % and each section \(x \in G.\mathrm{left}.\texttt{presheaf}.\mathrm{obj}\,X\),
  % the cross-open naturality square
  % \begin{align*}
  %   (d_Y).d\,(G.\mathrm{left}.\texttt{presheaf}.\mathrm{map}\,f\,x)
  %   &= ((\texttt{pushforward}\,\psi).\mathrm{obj}\,\texttt{LHS}).\mathrm{map}\,f\,((d_X).d\,x),
  % \end{align*}
  % where \(d_X(b) := \texttt{KaehlerDifferential.D}\,\_\,((\psi.\app\,X).\hom\,b)\)
  % is the pointwise derivation at open \(X\), and \texttt{LHS} denotes
  % \(\Omega_{(G \times_k G)/G} = \texttt{relativeDifferentialsPresheaf}\,(\texttt{fst}\,G\,G).\mathrm{left}\)
  % (the source side of the displayed iso, viewed at the
  % \texttt{PresheafOfModules}-level via the iter-138
  % \texttt{pushforward\,\(\psi\)} packaging).
  %
  % This square is a chase of two pieces of naturality:
  % \begin{enumerate}
  %   \item \textbf{\(\psi\).naturality \(f\)} (from
  %     \texttt{Scheme.Hom.c.naturality} applied to the second
  %     projection \(\texttt{snd}\,G\,G\)): translates the LHS's
  %     \((\psi.\app\,Y).\hom\,(G.\mathrm{left}.\texttt{presheaf}.\mathrm{map}\,f\,x)\)
  %     into
  %     \((G \otimes G).\mathrm{left}.\texttt{presheaf}.\mathrm{map}\,(\mathrm{snd}^{-1}.\mathrm{map}\,f)\,((\psi.\app\,X).\hom\,x)\).
  %   \item \textbf{\texttt{KaehlerDifferential.map\_d}} (Mathlib
  %     name \texttt{CommRingCat.KaehlerDifferential.map\_d}, available
  %     via \texttt{relativeDifferentials'\_map\_d} at
  %     \texttt{Mathlib/Algebra/Category/ModuleCat/Differentials/Presheaf.lean:201}):
  %     the universal Kähler derivation commutes with ring-map
  %     base-change, i.e.\ \(D'\,(f\,a) = M.\mathrm{map}\,f\,(D\,a)\)
  %     where \(D\) and \(D'\) are the universal Kähler derivations at
  %     the source and target of the ring map \(f\).
  % \end{enumerate}
  %
  % \textbf{Composing}: applying \(\texttt{KaehlerDifferential.D}\,\_\)
  % to both sides of the \(\psi\)-naturality identity (piece 1 above)
  % gives the required \texttt{d\_map} equation, after identifying
  % the right-hand side's
  % \(((\texttt{pushforward}\,\psi).\mathrm{obj}\,\texttt{LHS}).\mathrm{map}\,f\)
  % with the \texttt{LHS}-presheaf restriction
  % \(\texttt{LHS}.\mathrm{map}\,(\mathrm{snd}^{-1}.\mathrm{map}\,f)\).
  %
  % \emph{NOTE iter-141 (named lemma + \texttt{whnf}-disabled advisory).}
  % The unfolding $((\texttt{pushforward}\,\psi).\mathrm{obj}\,\texttt{LHS}).\mathrm{map}\,f
  % = \texttt{LHS}.\mathrm{map}\,(\mathrm{snd}^{-1}.\mathrm{map}\,f)$
  % is provided by the explicit \texttt{@[simp]}-tagged lemma
  % \texttt{PresheafOfModules.pushforward\_obj\_map\_apply'}
  % (\texttt{Mathlib/Algebra/Category/ModuleCat/Presheaf/Pushforward.lean:99--106}).
  % Use \texttt{simp only [pushforward\_obj\_map\_apply']} to perform
  % the unfolding; do \textbf{not} use \texttt{change} or other
  % \texttt{whnf}-based tactics, because \texttt{pushforward\(_0\)} is
  % annotated
  % \texttt{set\_option backward.isDefEq.respectTransparency false in}
  % (\texttt{Pushforward.lean:37, 55}) which explicitly disables the
  % transparency-based unfolding that \texttt{change}/\texttt{whnf}
  % rely on. The iter-140 prover validated this empirically: a
  % \texttt{change (CommRingCat.KaehlerDifferential.D \_).d \_ = \_}
  % attempt on the d\_map sub-goal hits a deterministic heartbeat
  % timeout at \texttt{maxHeartbeats=200000} (see
  % \texttt{task\_results/Cotangent\_GrpObj.lean.md} §"d\_map: detailed
  % gap" L111--L143).
  %
  % \textbf{Three-step ALIGN\_WITH\_MATHLIB chase} (per iter-141
  % \texttt{mathlib-analogist-d-app-d-map-iter141} Decision 4):
  % \begin{enumerate}
  %   \item \texttt{simp only [PresheafOfModules.pushforward\_obj\_map\_apply']}
  %     for the RHS unfolding.
  %   \item \texttt{NatTrans.naturality} for \(\psi\) — the natural
  %     transformation \((\texttt{Scheme.Hom.toRingCatSheafHom}\,(\texttt{snd}\,G\,G).\mathrm{left}).\mathrm{hom}\),
  %     which is \((\texttt{snd}\,G\,G).\mathrm{left}.c\) whiskered with
  %     \texttt{forget\(_2\)} per
  %     \texttt{Mathlib/AlgebraicGeometry/Modules/Presheaf.lean:42--45}.
  %   \item \texttt{PresheafOfModules.DifferentialsConstruction.relativeDifferentials'\_map\_d}
  %     (\texttt{Mathlib/Algebra/Category/ModuleCat/Differentials/Presheaf.lean:201--207})
  %     for the universal Kähler derivation commutation $D'\,(f\,a)
  %     = M.\mathrm{map}\,f\,(D\,a)$.
  % \end{enumerate}
  % LOC estimate for the d\_map closure: \(\sim 30\)--\(50\) LOC.
  %
  % \emph{Negative-lesson note (iter-141).} The iter-140 prover
  % attempted the iter-138 d\_add/d\_mul-style four-step
  % \texttt{have h := ... ; change (CommRingCat.KaehlerDifferential.D \_).d \_ = \_ ;
  % rw [h] ; exact ...} approach for d\_map. This approach IS valid
  % for d\_add/d\_mul (where the RHS is a pure
  % \texttt{KaehlerDifferential.D}-applied term and the \texttt{change}
  % only needs to align Kähler-derivation laws), but it
  % deterministically times out for d\_map, where the RHS carries
  % \(((\texttt{pushforward}\,\psi).\mathrm{obj}\,\texttt{LHS}).\mathrm{map}\,f\)
  % and \texttt{pushforward\(_0\)}'s transparency annotation explicitly
  % blocks \texttt{whnf} from unfolding it (see the named-lemma +
  % \texttt{whnf}-disabled advisory above). The d\_map pattern is
  % therefore distinct from d\_add/d\_mul: use
  % \texttt{simp only [pushforward\_obj\_map\_apply']} first to
  % bypass the opacity wall, \emph{then} apply \(\psi\)-naturality +
  % \texttt{relativeDifferentials'\_map\_d}. Future iters should
  % \textbf{not} re-attempt the \texttt{change}-first approach on
  % any \texttt{pushforward}-transposed goal.
  %
  % NOTE iter-137: the chart-by-chart prose recipe IMMEDIATELY BELOW is
  % preserved as informal mathematical motivation, NOT as the actual Lean
  % closure recipe. The iter-137 prover lane on the Lean target
  % \lean{AlgebraicGeometry.GrpObj.relativeDifferentialsPresheaf_basechange_along_proj_two}
  % returned PARTIAL with a load-bearing infrastructure blocker, identified
  % by the iter-137 mathlib-analogist verdict in
  % \texttt{analogies/kaehler-tensorequiv-presheafpullback.md} (Decision 4):
  %
  % \emph{The \texttt{PresheafOfModules.pullback} chart-opacity blocker.}
  % Mathlib defines \texttt{PresheafOfModules.pullback \(\varphi\)} as
  % \((\texttt{PresheafOfModules.pushforward}\,\varphi).\texttt{leftAdjoint}\)
  % (a left adjoint, packaged as the abstract left adjoint of the
  % pushforward functor between presheaves of modules over presheaves of
  % rings). Consequently, the right-hand side
  % \((\texttt{pullback}\,\psi).\mathrm{obj}\,\Omega_{G/k}\) has \emph{no}
  % computable chart-wise description on the
  % \texttt{.obj}/\texttt{.map} fields: those fields are opaque outputs of
  % the abstract adjunction, not direct constructions. The chart-by-chart
  % recipe written in the surrounding paragraph requires expressing the
  % right-hand-side sections
  % \(((\texttt{pullback}\,\psi).\mathrm{obj}\,\Omega_{G/k})\mathrm{.obj}\,V\)
  % as a tensor product
  % \(\Gamma((G \times_k G).\mathrm{left}, V) \otimes_{\Gamma(G.\mathrm{left},\,\mathrm{pr}_2(V))} \Omega[\Gamma(G.\mathrm{left},\,\mathrm{pr}_2(V)) \,\big|\, k]\)
  % for affine \(V\), which is \emph{not} a direct consequence of the
  % pullback's left-adjoint definition. The chart-affine-cover-and-glue
  % mental model also runs into a presheaf-vs-sheaf hazard: the project's
  % \(\Omega_{G/k} = \texttt{relativeDifferentialsPresheaf}\,G.\mathrm{hom}\)
  % is a presheaf (no sheafification), so a chart-wise identification on
  % affines does not automatically extend to non-affine opens, and
  % \texttt{KaehlerDifferential.tensorKaehlerEquiv}'s \texttt{Algebra.IsPushout}
  % hypothesis is not generally available on non-affine opens.
  %
  % \textbf{Two alternative closure paths} are identified by the
  % iter-137 mathlib-analogist (Decision 4 of
  % \texttt{analogies/kaehler-tensorequiv-presheafpullback.md}, refined LOC
  % envelope in Decision 5; see also \texttt{analogies/mulright-globalises-cotangent.md}
  % Decision 2 for the parent decomposition):
  %
  % \emph{Route (a): chart-unfolding-helper route} (\(\sim 30\)--\(60\) LOC for
  % the helper plus \(\sim 250\)--\(500\) LOC for the body, total
  % \(\sim 360\)--\(710\) LOC). Build a stand-alone Mathlib-shaped helper
  % \texttt{pullbackObjEquivTensor} that unfolds
  % \(((\texttt{pullback}\,\varphi).\mathrm{obj}\,M)\mathrm{.obj}\,V\) as a
  % tensor product
  % \(\Gamma(R, V) \otimes_{((F^{\mathrm{op}})_{*} S).\mathrm{obj}\,V} M.\mathrm{obj}(F^{\mathrm{op}}.\mathrm{obj}\,V)\)
  % via the \texttt{pullbackPushforwardAdjunction} unit/counit (using the
  % left-adjoint definition and the explicit tensor-of-modules
  % shape generated by \texttt{ModuleCat.restrictScalars} +
  % \texttt{TensorProduct} on each open). With this helper in hand,
  % execute the iter-137 universal-property-at-presheaf-level 5-step
  % recipe:
  % \begin{itemize}
  %   \item[(1)] Chart-level \texttt{Algebra.IsPushout} (\(\sim 80\)--\(150\) LOC):
  %     for affine \(V_1 \subseteq G.\mathrm{left}\), \(V_2 \subseteq G.\mathrm{left}\)
  %     over a common affine \(U \subseteq \Spec k\) with
  %     \(W = V_1 \times_U V_2 \subseteq (G \times_k G).\mathrm{left}\) cut
  %     out via \texttt{pullbackSpecIso}, derive \texttt{Algebra.IsPushout}
  %     \(\Gamma(\Spec k, U)\) \(\Gamma(G.\mathrm{left}, V_1)\)
  %     \(\Gamma(G.\mathrm{left}, V_2)\) \(\Gamma((G \times_k G).\mathrm{left}, W)\)
  %     from \texttt{CommRingCat.isPushout\_iff\_isPushout},
  %     \texttt{pullbackSpecIso}, and
  %     \texttt{isPullback\_SpecMap\_of\_isPushout}.
  %   \item[(2)] \texttt{PresheafOfModules.pullback} chart-unfolding via
  %     the new \texttt{pullbackObjEquivTensor} helper.
  %   \item[(3)] Chart-wise derivation \(D_V\) (\(\sim 50\)--\(80\) LOC): for
  %     each open \(V \subseteq (G \times_k G).\mathrm{left}\), define a
  %     \(\Gamma(G.\mathrm{left}, \mathrm{pr}_1(V))\)-derivation \(D_V\) of
  %     \(\Gamma((G \times_k G).\mathrm{left}, V)\) into the pulled-back
  %     module by chaining the chart-level value of
  %     \texttt{KaehlerDifferential.tensorKaehlerEquiv} with the unfolding
  %     helper from step (2).
  %   \item[(4)] Apply \texttt{KaehlerDifferential.lift} (\(\sim 30\)--\(50\) LOC):
  %     on each \(V\), the universal property of Kähler differentials produces
  %     a unique linear extension of \(D_V\) to
  %     $\Omega[\Gamma((G \times_k G).\mathrm{left}, V) \,\big|\,
  %       \Gamma(G.\mathrm{left}, \mathrm{pr}_1(V))]$, and naturality of
  %     \texttt{lift} glues these into a morphism of presheaves of modules.
  %   \item[(5)] Inverse + \texttt{PresheafOfModules.isoMk} assembly
  %     (\(\sim 80\)--\(150\) LOC): construct the inverse the same way and
  %     package the pair with \texttt{PresheafOfModules.isoMk}; the
  %     naturality square reduces to commutativity of \texttt{lift}-of-
  %     derivation with restriction maps, follows from
  %     \texttt{relativeDifferentials'\_map\_d}.
  % \end{itemize}
  %
  % \emph{Route (b): inverse-direction-via-adjunction-transpose route}
  % (\(\sim 100\)--\(200\) LOC for the inverse-direction derivation construction
  % alone; the forward direction is still required separately). The
  % iter-137 prover (Attempt 2 in
  % \texttt{task\_results/Cotangent\_GrpObj.lean.md}, lines 40--89) validated
  % a closure of the \emph{inverse} direction
  % $(\texttt{pullback}\,\psi).\mathrm{obj}\,\Omega_{G/k} \;\longrightarrow\;
  %    \Omega_{(G \times_k G)/G}$
  % that bypasses the \texttt{pullback} opacity entirely on the inverse
  % side. The route applies the
  % \texttt{pullbackPushforwardAdjunction} on
  % \texttt{PresheafOfModules} (\texttt{ModuleCat/Presheaf/Pullback.lean})
  % to transpose this morphism to a morphism
  % \(\Omega_{G/k} \;\longrightarrow\; (\texttt{pushforward}\,\psi).\mathrm{obj}\,\Omega_{(G \times_k G)/G}.\)
  % By the universal property of
  % $\Omega_{G/k} = \texttt{relativeDifferentialsPresheaf}\,G.\mathrm{hom}
  %   = \texttt{relativeDifferentials'}\,\varphi_G$
  % (\texttt{DifferentialsConstruction.isUniversal'},
  % \texttt{Mathlib/Algebra/Category/ModuleCat/Differentials/Presheaf.lean:189}),
  % such a morphism corresponds bijectively to a
  % \((\texttt{pushforward}\,\psi).\mathrm{obj}\,\Omega_{(G \times_k G)/G}\)-valued
  % derivation of \(\varphi_G\). Crucially, the pushforward
  % \((\texttt{pushforward}\,\psi).\mathrm{obj}\,(-)\) is \emph{transparent}
  % on \texttt{.obj}/\texttt{.map}: per
  % \texttt{Mathlib/Algebra/Category/ModuleCat/Presheaf/Pushforward.lean}
  % (lines 39 and 86), \texttt{PresheafOfModules.pushforward} is
  % \(\texttt{pushforward}_{0} \;\circ\; \texttt{restrictScalars}\), with
  % both pieces carrying explicit \texttt{@[simps]} definitions. The
  % derivation is therefore constructible pointwise at each opposite open
  % \(X'\) as the universal Kähler derivation of
  % \(\Gamma((G \times_k G).\mathrm{left},\,F^{\mathrm{op}}.\mathrm{obj}\,X')\)
  % post-composed with the structure ring map
  % \(\psi.\mathrm{app}\,X'\). The iter-137 prover validated this
  % inverse-direction skeleton as compiling-typeable via
  % \texttt{lean\_run\_code}, reducing closure of the inverse direction to
  % a single concrete sub-goal: constructing the derivation \(D\). The
  % \emph{forward} direction
  % \(\Omega_{(G \times_k G)/G} \to (\texttt{pullback}\,\psi).\mathrm{obj}\,\Omega_{G/k}\)
  % remains blocked by the same opacity as Route (a): producing a
  % derivation valued in the right-hand side at the source's universal
  % property requires the chart-unfolding helper. An alternative closing
  % the iso from the inverse alone — checking the iso property locally on
  % \texttt{PresheafOfModules} generators — is on the table but not yet
  % scoped.
  %
  % \textbf{Interim Lean-side record.} Until this NOTE block is reflected
  % in revised proof prose, the Lean docstring at
  % \texttt{AlgebraicJacobian/Cotangent/GrpObj.lean:479--499} carries the
  % authoritative iter-137 analysis (with concrete signature, both routes'
  % LOC envelopes, and the validated-typeable inverse-direction skeleton).
  % Future iter-138+ provers SHOULD consult both the blueprint prose
  % AND the Lean docstring before attempting closure: the docstring is
  % more recent than the surrounding chart-by-chart prose, and the chart-
  % by-chart recipe BELOW is informal motivation only.
  %
  % \textbf{NOTE iter-143 (Lean-shape refactor; in-Lean only, recipe
  % unchanged).} The in-Lean encoding of the IsIso residual has been
  % refactored iter-143 from the in-line
  % \texttt{letI := isIso\_of\_app\_iso\_module
  % (basechange\_along\_proj\_two\_inv G) (fun \_ => sorry)} pattern
  % (which hid the per-open sorry inside a \texttt{letI} body and
  % propagated a sorry-tainted iso to downstream \texttt{simp}
  % consumers; see \texttt{Cotangent/GrpObj.lean:719--721} pre-refactor)
  % into a named sorry-bodied theorem
  % \lean{AlgebraicGeometry.GrpObj.basechange_along_proj_two_inv_app_isIso}
  % at \texttt{Cotangent/GrpObj.lean}. The named theorem carries the
  % per-open obligation
  % \(\forall\,X,\;\texttt{IsIso}\,((\texttt{basechange\_along\_proj\_two\_inv}\,G).\app\,X)\)
  % directly; the consuming declaration
  % \texttt{relativeDifferentialsPresheaf\_basechange\_along\_proj\_two}
  % invokes the named theorem inside its \texttt{letI} chain rather
  % than carrying the \texttt{(fun \_ => sorry)} argument inline.
  % \textbf{The Route (b'2) closure recipe below (items 2--4) is
  % unchanged}: the iter-144+ prover lane closes the body of the
  % named theorem
  % \texttt{basechange\_along\_proj\_two\_inv\_app\_isIso} rather than
  % filling a \texttt{(fun \_ => sorry)} argument inside the consumer.
  % The refactor was performed per the
  % \texttt{lean-auditor-review142} MAJOR finding (the in-line sorry
  % shape obscures the obligation), the iter-143
  % \texttt{progress-critic-iter143} CHURNING primary corrective
  % (decouple the IsIso closure from the consumer to enable an
  % isolated iter-144+ prover round), and the iter-143
  % \texttt{strategy-critic-iter143} definition-level diagnostic on
  % the IsIso \texttt{letI} shape. The iter-143 prover lane targets
  % \emph{only} the d\_app sub-sorry (narrowed from the bundled
  % 3-sub-sorry shape of iter-142; d\_map was closed iter-142, the
  % named IsIso theorem is deferred to iter-144+ as a separate
  % prover round).
  %
  % \textbf{Route (b'2): local-iso check via
  % \texttt{PresheafOfModules.toPresheaf} +
  % \texttt{NatTrans.isIso\_iff\_isIso\_app} (\(\sim 195\)--\(365\) LOC;
  % per the iter-139 mathlib-analogist verdict
  % \texttt{mathlib-analogist-isiso-routes-iter139} at
  % \texttt{analogies/isiso-basechange-along-proj-two-inv.md}).} Iter-138
  % Route (b) landed the inverse-direction-via-adjunction-transpose
  % skeleton end-to-end (additive + Leibniz laws of the pointwise
  % derivation closed honestly; three concrete sub-sorries remain —
  % see the iter-138 NOTE block at the head of this proof). The iso
  % \(\texttt{IsIso}\,(\texttt{basechange\_along\_proj\_two\_inv}\,G)\) is
  % the third remaining concrete sorry at
  % \texttt{Cotangent/GrpObj.lean:624}; the iter-139 analogist
  % returned \textbf{PROCEED with Route (b'2)} as the cheapest closure
  % path for iter-140.
  %
  % \textbf{The 5-line iso-reflection bridge.} Route (b'2)'s key
  % infrastructure piece is a private helper
  % \texttt{isIso\_of\_app\_iso\_module} placed in
  % \texttt{Cotangent/GrpObj.lean} adjacent to
  % \texttt{basechange\_along\_proj\_two\_inv}:
  % \begin{verbatim}
  %   theorem isIso_of_app_iso_module {C : Type*} [Category C]
  %       {R : Cᵒᵖ ⥤ RingCat} {M N : PresheafOfModules R}
  %       (f : M ⟶ N) (h : ∀ X, IsIso (f.app X)) : IsIso f := by
  %     rw [← isIso_iff_of_reflects_iso _ (PresheafOfModules.toPresheaf R),
  %         NatTrans.isIso_iff_isIso_app]
  %     intro X
  %     exact Functor.map_isIso (forget₂ (ModuleCat _) AddCommGrpCat) (f.app X)
  % \end{verbatim}
  % It chains two Mathlib facts: (i) the forgetful functor
  % \texttt{PresheafOfModules.toPresheaf\,R : PresheafOfModules R \(\mathbin{\Rightarrow}\) C\(^{\mathrm{op}}\) \(\mathbin{\Rightarrow}\) Ab}
  % at \texttt{Mathlib/Algebra/Category/ModuleCat/Presheaf.lean:149}
  % reflects isomorphisms, via the priority-100 instance
  % \texttt{reflectsIsomorphisms\_of\_reflectsMonomorphisms\_of\_reflectsEpimorphisms}
  % at \texttt{Mathlib/CategoryTheory/Functor/ReflectsIso/Balanced.lean:31},
  % combining \texttt{Balanced (PresheafOfModules R)} +
  % \texttt{Faithful (toPresheaf R)} (already an instance at
  % \texttt{Presheaf.lean:164}) \(\Rightarrow\)
  % \texttt{ReflectsMonomorphisms} + \texttt{ReflectsEpimorphisms}; and
  % (ii) \texttt{NatTrans.isIso\_iff\_isIso\_app} at
  % \texttt{Mathlib/CategoryTheory/NatIso.lean:232}, reducing the
  % iso-check on a natural transformation to a pointwise iso-check at
  % every object. The import of
  % \texttt{Mathlib.CategoryTheory.Functor.ReflectsIso.Balanced} must
  % be present (transitively or directly) for the typeclass synthesis
  % to succeed; it is reachable transitively via
  % \texttt{Mathlib.Algebra.Category.ModuleCat.Presheaf.Sheafification}
  % which the project already imports.
  %
  % \textbf{Closure shape.} With the bridge in hand, the iso check
  % reduces to verifying that the per-open \texttt{ModuleCat}-morphism
  % \((\texttt{basechange\_along\_proj\_two\_inv}\,G).\app\,X\) is an
  % iso for every open \(X\). Identifying this per-open morphism with
  % \texttt{KaehlerDifferential.tensorKaehlerEquiv.symm} (from
  % \texttt{Mathlib/RingTheory/Kaehler/TensorProduct.lean}) modulo the
  % chart-unfolding helper \texttt{pullbackObjEquivTensor} (Route (a)'s
  % Step 2 of the 5-step recipe — also load-bearing here, since Route
  % (b'2) does \emph{not} escape the \texttt{pullback}-opacity issue:
  % describing the source \(((\texttt{pullback}\,\psi).\mathrm{obj}\,M_G).\mathrm{obj}\,X\)
  % per-open to compare with \texttt{tensorKaehlerEquiv}'s output
  % requires precisely the same chart-unfolding) and the chart-level
  % \texttt{Algebra.IsPushout} square (Route (a)'s Step 1 — shared)
  % discharges the goal. The companion simp lemma
  % \texttt{tensorKaehlerEquiv\_symm\_D\_tmul} handles the value
  % identity at the level of \(D\,b \mapsto 1 \otimes D\,b\).
  %
  % \textbf{Iter-140 prover gap items, in build order.}
  %
  % \textbf{Iter-140 update.} Item (1), the iso-reflection bridge
  % \texttt{isIso\_of\_app\_iso\_module}, is closed in
  % \texttt{AlgebraicJacobian/Cotangent/GrpObj.lean:544--550}. Items
  % (2)--(4) remain iter-141+ targets. The residual per-open sorry at
  % \texttt{Cotangent/GrpObj.lean:689} has type
  % \(\forall\,X,\;\texttt{IsIso}\,((\texttt{basechange\_along\_proj\_two\_inv}\,G).\app\,X)\)
  % (iter-140 narrowed the sorry from "whole iso" \texttt{letI := sorry}
  % to the per-open variant inside
  % \texttt{isIso\_of\_app\_iso\_module ... (fun \_ => sorry)}).
  % \begin{enumerate}
  %   \item \textbf{[iter-140 closed]} The 5-line
  %     \texttt{isIso\_of\_app\_iso\_module} helper (landed at
  %     \texttt{Cotangent/GrpObj.lean:544--550}; verified iter-140
  %     typecheck).
  %   \item \textbf{[iter-144+ target]} The chart-level
  %     \texttt{Algebra.IsPushout}-from-affine-product helper
  %     (\(\sim 80\)--\(150\) LOC; shared with Route (a)).
  %     \emph{Concrete recipe sketch (iter-143).} For an affine open
  %     \(W = \Spec B \subseteq (G \times_k G).\mathrm{left}\) cut out
  %     via the pullback square \(W = V_1 \times_U V_2\) with affine
  %     opens \(V_1, V_2 \subseteq G.\mathrm{left}\) over a common
  %     affine \(U \subseteq \Spec k\), the chart-level
  %     \texttt{Algebra.IsPushout} square chains three Mathlib facts:
  %     (i) \texttt{CommRingCat.isPushout\_iff\_isPushout}
  %     (\texttt{Mathlib.Algebra.Category.Ring.Pushout}; converts a
  %     \texttt{IsPushout} square of \texttt{CommRingCat}-morphisms
  %     into an \texttt{Algebra.IsPushout} square of the underlying
  %     ring-algebra morphisms); (ii) \texttt{pullbackSpecIso}
  %     (\texttt{Mathlib.AlgebraicGeometry.Pullbacks}; gives the
  %     canonical iso \(\Spec B \cong \Spec B_1 \times_{\Spec A} \Spec B_2\)
  %     identifying \(B\) with the ring-side tensor product
  %     \(B_1 \otimes_A B_2\)); (iii)
  %     \texttt{isPullback\_SpecMap\_of\_isPushout} (\texttt{Mathlib.AlgebraicGeometry.Pullbacks}
  %     — \texttt{Spec} sends pushout squares of rings to pullback
  %     squares of schemes, the converse statement is used here to
  %     derive the pushout shape on the ring side from a pullback
  %     shape on the scheme side). Per Rule 1 of the iter-142
  %     empirical lessons above, any \texttt{change}-based alignment
  %     used here must spell both LHS and RHS explicitly.
  %   \item \textbf{[iter-144+ target]} The
  %     \(((\texttt{pullback}\,\psi).\mathrm{obj}\,M).\mathrm{obj}\,X\)
  %     chart-unfolding helper \texttt{pullbackObjEquivTensor}
  %     (\(\sim 30\)--\(60\) LOC; shared with Route (a)).
  %     \emph{Concrete recipe sketch (iter-143).} The helper unfolds
  %     a section of the pulled-back presheaf-of-modules
  %     \(((\texttt{pullback}\,\varphi).\obj\,M).\obj\,V\) as the
  %     restricted tensor product
  %     $\Gamma(R, V) \otimes_{((F^{\mathrm{op}})_{*} S).\obj\,V}
  %       M.\obj(F^{\mathrm{op}}.\obj\,V)$
  %     using the unit/counit of
  %     \texttt{pullbackPushforwardAdjunction} on
  %     \texttt{PresheafOfModules}
  %     (\texttt{Mathlib/Algebra/Category/ModuleCat/Presheaf/Pullback.lean:50}).
  %     The shape is a definitional equality on the
  %     \texttt{pushforward.leftAdjoint.obj} construction: the
  %     forgetful \texttt{toPresheaf} sends the pullback to the
  %     pointwise tensor product modulo the
  %     \texttt{ModuleCat.restrictScalars}-coercion of the structure
  %     ring at each open. The helper is needed verbatim by both
  %     Routes (a) and (b'2); it does not escape the
  %     \texttt{pullback}-opacity issue but exposes the chart-level
  %     tensor shape that the iso check (item 4) requires.
  %   \item \textbf{[iter-144+ target]} The per-open identification
  %     with \texttt{tensorKaehlerEquiv.symm} (\(\sim 80\)--\(150\) LOC).
  %     \emph{Concrete recipe sketch (iter-143).} With items 2+3 in
  %     hand, the per-open iso
  %     $((\texttt{pullback}\,\psi).\obj\,\Omega_{G/k}).\obj\,X
  %       \;\xrightarrow{\;\sim\;}\;
  %       \Omega_{(G \times_k G)/G}.\obj\,X$
  %     for affine \(X = W\) is the value of
  %     \texttt{KaehlerDifferential.tensorKaehlerEquiv.symm}
  %     (\texttt{Mathlib.RingTheory.Kaehler.TensorProduct}) applied to
  %     the chart-level \texttt{Algebra.IsPushout} square from item 2,
  %     re-presented through the chart-unfolding equivalence
  %     \texttt{pullbackObjEquivTensor} from item 3. The
  %     algebra-side value identity is
  %     \texttt{tensorKaehlerEquiv\_symm\_D\_tmul}
  %     (\(D\,b \mapsto 1 \otimes D\,b\), where \(D\) is the universal
  %     Kähler derivation on each chart). Composition of items 2+3
  %     produces the per-open iso for affine opens; the iter-138
  %     Route (b) iso-reflection bridge
  %     \texttt{isIso\_of\_app\_iso\_module} (item 1, closed iter-140)
  %     then promotes the pointwise iso to a presheaf-of-modules iso
  %     without requiring a separate forward-direction construction.
  % \end{enumerate}
  %
  % \textbf{Mathlib API verified iter-139} (per the analogist's
  % \texttt{lean\_loogle} + \texttt{lean\_run\_code} typecheck pass):
  % \texttt{PresheafOfModules.toPresheaf} at
  % \texttt{Mathlib/Algebra/Category/ModuleCat/Presheaf.lean:149};
  % \texttt{NatTrans.isIso\_iff\_isIso\_app} at
  % \texttt{Mathlib/CategoryTheory/NatIso.lean:232};
  % \texttt{PresheafOfModules.pullback} at
  % \texttt{Mathlib/Algebra/Category/ModuleCat/Presheaf/Pullback.lean:44};
  % \texttt{PresheafOfModules.pullbackPushforwardAdjunction} at
  % \texttt{Mathlib/Algebra/Category/ModuleCat/Presheaf/Pullback.lean:50};
  % \texttt{KaehlerDifferential.tensorKaehlerEquiv} at
  % \texttt{Mathlib/RingTheory/Kaehler/TensorProduct.lean};
  % \texttt{isIso\_iff\_of\_reflects\_iso} at
  % \texttt{Mathlib/CategoryTheory/Functor/ReflectsIso/Basic.lean:49};
  % \texttt{reflectsIsomorphisms\_of\_reflectsMonomorphisms\_of\_reflectsEpimorphisms}
  % at \texttt{Mathlib/CategoryTheory/Functor/ReflectsIso/Balanced.lean:31}.
  %
  % \textbf{Why Route (b'2) over Route (a).} Route (a) requires
  % explicitly constructing the forward direction
  % \(\Omega_{(G \otimes G)/G} \to (\texttt{pullback}\,\psi).\mathrm{obj}\,\Omega_{G/k}\)
  % \emph{as well as} the inverse pair, paying \(\sim 80\)--\(200\) LOC of
  % extra work for the forward-direction construction and the
  % inverse-pair iso proof. Route (b'2) uses the 5-line iso-reflection
  % bridge \texttt{isIso\_of\_app\_iso\_module} to skip the forward
  % direction entirely (iter-138's Route (b) already supplied the
  % inverse direction), paying instead for a slightly heavier per-open
  % identification. Refined LOC envelopes per the iter-139 analogist:
  % Route (a) \(\sim 280\)--\(560\) LOC (helpers inclusive), Route (b'2)
  % \(\sim 195\)--\(365\) LOC (helpers inclusive); net savings
  % \(\sim 80\)--\(195\) LOC. Both routes share the heavy lifting on
  % \texttt{pullbackObjEquivTensor} and the chart-level
  % \texttt{Algebra.IsPushout}; these helpers are upstream-PR
  % candidates regardless of the route choice. The
  % \texttt{pullback}-opacity issue itself is \emph{not} escaped by
  % Route (b'2) — the chart-unfolding helper is needed by both
  % routes — but Route (b'2)'s savings come specifically from
  % avoiding the explicit forward-direction construction.
  %
  The construction chains two Mathlib idioms. On each affine chart \(W = \Spec B\) of the source \(G \times_{k} G\), the sections of both sides are Kähler-differential modules with appropriate base rings. The left-hand side \(\Omega_{(G \times_{k} G)/G}\) over \(W\) gives \(\Omega_{B/B_{1}}\), where \(B_{1}\) corresponds to the chart of \(G\) obtained by restricting \(\pr_{1}\) to \(W\). The right-hand side \(\pr_{2}^{*}\,\Omega_{G/k}\) over \(W\) gives \(B \otimes_{B_{2}} \Omega_{B_{2}/k}\), where \(B_{2}\) corresponds to the chart of \(G\) obtained by restricting \(\pr_{2}\). The binary-product universal property in \(\Over\,(\Spec k)\) supplies an \texttt{Algebra.IsPushout}\,\(k\)\,\(B_{1}\)\,\(B_{2}\)\,\(B\) square (the two projections \(B_{1}, B_{2} \to B\) exhibit \(B\) as the tensor product \(B_{1} \otimes_{k} B_{2}\) at chart level). The algebra-side base-change identity
  \[
    \texttt{KaehlerDifferential.tensorKaehlerEquiv} \colon \;\; B \,\otimes_{B_{2}}\, \Omega_{B_{2}/k} \;\xrightarrow{\;\sim\;}\; \Omega_{B/B_{1}}
  \]
  (\texttt{Mathlib.RingTheory.Kaehler.TensorProduct}; for an \texttt{Algebra.IsPushout R S A B} square, yields \(B \otimes_{A} \Omega[A/R] \simeq_{B} \Omega[B/S]\)) gives the chart-level value identity. Promoting this chart-by-chart equivalence to a natural iso of presheaves of modules uses the presheaf-side functor \texttt{TopCat.Presheaf.pullback} (\texttt{Mathlib.Topology.Sheaves.Pullback}) composed with the project-side chart-level value identity \texttt{relativeDifferentialsPresheaf\_obj\_kaehler} (\cref{lem:relative_kaehler_presheaf_obj}, \texttt{AlgebraicJacobian/Differentials.lean:60}); naturality follows from naturality of \texttt{tensorKaehlerEquiv} and of \texttt{TopCat.Presheaf.pullback}.

  \emph{Mathlib name summary (NEEDS\_MATHLIB\_GAP\_FILL, \(\sim 150\)--\(300\) LOC).} Algebra-side: \texttt{KaehlerDifferential.tensorKaehlerEquiv} (\texttt{Mathlib.RingTheory.Kaehler.TensorProduct}). Presheaf-side: \texttt{TopCat.Presheaf.pullback} (\texttt{Mathlib.Topology.Sheaves.Pullback}) and \texttt{PresheafOfModules.pullback} (\texttt{Mathlib.Algebra.Category.ModuleCat.Presheaf.Pullback}). Project-side: \texttt{relativeDifferentialsPresheaf\_obj\_kaehler} (\cref{lem:relative_kaehler_presheaf_obj}, \texttt{AlgebraicJacobian/Differentials.lean:60--66}). This is the load-bearing piece of (i.b); without it, the shear iso of Step 1 of \cref{lem:GrpObj_mulRight_globalises}'s proof cannot be turned into a cotangent-level iso. Mathlib has no scheme-level relative cotangent sheaf at all (per \texttt{analogies/mulright-globalises-cotangent.md} Decision 2), so this lemma is a genuine NEEDS\_MATHLIB\_GAP\_FILL contribution candidate. Compatible upstream Mathlib PR-shape: a sheaf-of-modules version of \texttt{tensorKaehlerEquiv} promoting \texttt{Algebra.IsPushout}-of-charts to a presheaf-level natural iso.
\end{proof}

\begin{lemma}[(i.b-helper) Inverse-direction derivation for the base-change iso]
  \label{lem:GrpObj_omega_basechange_proj_inv_derivation}
  \lean{AlgebraicGeometry.TODO.GrpObj_omega_basechange_proj_inv_derivation}
  \uses{def:relative_kaehler_presheaf, def:GrpObj_schemeHomRingCompatibility}
  % Iter-146 disposition: \lean{...basechange_along_proj_two_inv_derivation}
  % stripped — Lean target EXCISED iter-145 (see EXCISE disposition
  % paragraph above \cref{lem:GrpObj_omega_basechange_proj}). Block kept
  % as auditable bundled-route history.
  % Lean signature stub (iter-138 helper for piece (i.b); landed
  % substantively iter-138 with `d_add` and `d_mul` closed honestly,
  % `d_app` and `d_map` carried as concrete sub-sorries — iter-140
  % prover-lane targets per the iter-138 NOTE block in the proof of
  % \cref{lem:GrpObj_omega_basechange_proj}):
  %   noncomputable def basechange_along_proj_two_inv_derivation
  %       (G : Over (Spec (.of k))) [CategoryTheory.GrpObj G] :
  %       ((PresheafOfModules.pushforward
  %             (Scheme.Hom.toRingCatSheafHom (snd G G).left).hom).obj
  %           (Scheme.relativeDifferentialsPresheaf (fst G G).left)).Derivation'
  %         (((TopCat.Presheaf.pullbackPushforwardAdjunction CommRingCat
  %           G.hom.base).homEquiv _ _).symm G.hom.c)
  
  % NOTE iter-140: \notready stripped — Lean target
  % `\lean{AlgebraicGeometry.GrpObj.basechange_along_proj_two_inv_derivation}`
  % is formalized at `Cotangent/GrpObj.lean:573` (iter-138 helper; d_add
  % and d_mul laws closed; d_app L624 + d_map L643 sub-sorries iter-141+
  % targets). Per CLAUDE.md marker vocab, `\notready` on a formalized
  % statement is stale.
  % NOTE iter-146: the Lean target was subsequently EXCISED iter-145
  % from `AlgebraicJacobian/Cotangent/GrpObj.lean` under the iter-144
  % chart-algebra pivot; the statement is once again blueprint-only.
  \notready
  Let \(G : \Over\,(\Spec k)\) be as in \cref{lem:GrpObj_omega_basechange_proj}, and let
  \[
    \psi \;:=\; (\Scheme.\Hom.\texttt{toRingCatSheafHom}\,(\texttt{snd}\,G\,G).\mathrm{left}).\hom
  \]
  be the compatibility morphism for the second projection
  \(\pr_2 = \texttt{snd}\,G\,G : G \otimes G \to G\) (the iter-135
  \texttt{Scheme.Hom.toRingCatSheafHom} idiom, in the shape consumed by
  \texttt{PresheafOfModules.pullback}; see
  \cref{rem:GrpObj_schemeHomRingCompatibility_vs_toRingCatSheafHom}).
  There is a \(\varphi_G\)-derivation
  \[
    D \;\colon\; \bigl((\texttt{PresheafOfModules.pushforward}\,\psi).\mathrm{obj}\,\Omega_{(G \times_{k} G)/G}\bigr).\texttt{Derivation'}\,\varphi_G
  \]
  on the transparent pushforward of the left-hand-side relative
  cotangent along \(\psi\), where \(\varphi_G\) is the
  \(\Over\,(\Spec k)\)-side compatibility morphism for \(G.\mathrm{hom}\)
  obtained via the homEquiv-transpose
  \(((\texttt{TopCat.Presheaf.pullbackPushforwardAdjunction}\,\CommRingCat\,G.\mathrm{hom}.\mathrm{base}).\texttt{homEquiv}\,\_\,\_).\texttt{symm}\,G.\mathrm{hom}.c\).
  The derivation \(D\) is defined pointwise at each open
  \(X \in G.\mathrm{left}.\Opens^{\mathrm{op}}\) by
  \[
    d_X(b) \;:=\; \texttt{KaehlerDifferential.D}\,\_\,((\psi.\app\,X).\hom\,b),
  \]
  i.e.\ push the section \(b \in G.\mathrm{left}.\texttt{presheaf}.\mathrm{obj}\,X\)
  forward along \(\psi.\app\,X\) to a section of
  \((G \otimes G).\mathrm{left}.\texttt{presheaf}.\mathrm{obj}\,(\mathrm{snd}^{-1}\,X)\),
  then take its universal Kähler differential in the source-side
  presheaf of relative differentials of
  \((\texttt{fst}\,G\,G).\mathrm{left}\) over the pulled-back open.

  The Lean construction uses \texttt{PresheafOfModules.Derivation'.mk}
  at the per-open level via \texttt{ModuleCat.Derivation.mk} with the
  rule above; the additive (\texttt{d\_add}) and Leibniz
  (\texttt{d\_mul}) laws follow from the \texttt{RingHom}-ness of
  \(\psi.\app\,X\) composed with the algebra-side derivation laws
  \texttt{ModuleCat.Derivation.d\_add} / \texttt{d\_mul} on
  \texttt{KaehlerDifferential.D}. The \texttt{d\_app}
  (algebraic-coherence, zero on \(\varphi_G\)-image) and \texttt{d\_map}
  (cross-open naturality) sub-goals remain \texttt{sorry}-bodied in
  iter-138; see the iter-138 NOTE block of
  \cref{lem:GrpObj_omega_basechange_proj}'s proof above for the
  closure recipes.
\end{lemma}

\begin{proof}
  
  Direct construction via \texttt{PresheafOfModules.Derivation'.mk} +
  \texttt{ModuleCat.Derivation.mk} at the per-open level, with the
  pointwise rule
  \(d_X(b) := \texttt{KaehlerDifferential.D}\,\_\,((\psi.\app\,X).\hom\,b)\).
  The four derivation laws are: \texttt{d\_add} and \texttt{d\_mul}
  (closed iter-138 via the \texttt{RingHom}-ness of \(\psi.\app\,X\)
  composed with the algebra-side
  \texttt{ModuleCat.Derivation.d\_add} / \texttt{d\_mul}, using the
  \texttt{have\,h ... ; change ... ; rw [h] ; exact ...}
  beta-redex-aware pattern codified in
  \texttt{.archon/.debug-feedback/debug\_feedback.md}); \texttt{d\_app}
  and \texttt{d\_map} (iter-140 prover-lane targets, closure recipes
  in the iter-138 NOTE block of
  \cref{lem:GrpObj_omega_basechange_proj}'s proof above).
\end{proof}

\begin{lemma}[(i.b-helper) Inverse-direction morphism for the base-change iso]
  \label{lem:GrpObj_omega_basechange_proj_inv}
  \lean{AlgebraicGeometry.TODO.GrpObj_omega_basechange_proj_inv}
  \uses{lem:GrpObj_omega_basechange_proj_inv_derivation, def:relative_kaehler_presheaf}
  % Iter-146 disposition: \lean{...basechange_along_proj_two_inv}
  % stripped — Lean target EXCISED iter-145 (see EXCISE disposition
  % paragraph above \cref{lem:GrpObj_omega_basechange_proj}). Block kept
  % as auditable bundled-route history.
  % Lean signature stub (iter-138 helper for piece (i.b); the inverse
  % map for the base-change iso, assembled via universal-property
  % `desc` on the iter-138 derivation followed by the
  % `pullbackPushforwardAdjunction` transpose. The iso property
  % `IsIso (basechange_along_proj_two_inv G)` is the third concrete
  % sub-sorry of \cref{lem:GrpObj_omega_basechange_proj}'s proof block,
  % iter-140 prover-lane target via Route (b'2) per
  % `analogies/isiso-basechange-along-proj-two-inv.md`):
  %   noncomputable def basechange_along_proj_two_inv
  %       (G : Over (Spec (.of k))) [CategoryTheory.GrpObj G] :
  %       (PresheafOfModules.pullback
  %           (Scheme.Hom.toRingCatSheafHom (snd G G).left).hom).obj
  %         (Scheme.relativeDifferentialsPresheaf G.hom) ⟶
  %       Scheme.relativeDifferentialsPresheaf (fst G G).left
  
  % NOTE iter-140: \notready stripped — Lean target
  % `\lean{AlgebraicGeometry.GrpObj.basechange_along_proj_two_inv}`
  % is formalized at `Cotangent/GrpObj.lean:654` (iter-138 helper;
  % sorry-free as a definition; the iso property is the third concrete
  % sub-sorry of `\cref{lem:GrpObj_omega_basechange_proj}`'s proof block).
  % Per CLAUDE.md marker vocab, `\notready` on a formalized statement is stale.
  % NOTE iter-146: the Lean target was subsequently EXCISED iter-145
  % from `AlgebraicJacobian/Cotangent/GrpObj.lean` under the iter-144
  % chart-algebra pivot; the statement is once again blueprint-only.
  \notready
  With \(G, \psi\) as in
  \cref{lem:GrpObj_omega_basechange_proj_inv_derivation}, the
  universal property of
  \(\Omega_{G/k} = \texttt{relativeDifferentialsPresheaf}\,G.\mathrm{hom}\)
  (Mathlib's \texttt{PresheafOfModules.DifferentialsConstruction.isUniversal'}
  at \texttt{Mathlib/Algebra/Category/ModuleCat/Differentials/Presheaf.lean:189})
  applied to the derivation \(D\) of
  \cref{lem:GrpObj_omega_basechange_proj_inv_derivation}, then
  transposed along the
  \texttt{PresheafOfModules.pullbackPushforwardAdjunction\,\(\psi\)}
  (at \texttt{Mathlib/Algebra/Category/ModuleCat/Presheaf/Pullback.lean:50}),
  produces an inverse morphism
  \[
    \texttt{basechange\_along\_proj\_two\_inv}\,G \;\colon\;
    (\texttt{PresheafOfModules.pullback}\,\psi).\mathrm{obj}\,\Omega_{G/k}
    \;\longrightarrow\; \Omega_{(G \times_{k} G)/G}.
  \]
  This morphism is the inverse of the iso of
  \cref{lem:GrpObj_omega_basechange_proj}; the iso property
  \(\texttt{IsIso}\,(\texttt{basechange\_along\_proj\_two\_inv}\,G)\)
  is the third concrete sub-sorry of
  \cref{lem:GrpObj_omega_basechange_proj}'s proof block (iter-140
  prover-lane target, Route (b'2) per
  \texttt{analogies/isiso-basechange-along-proj-two-inv.md} — see the
  Route (b'2) sub-paragraph of the IsIso NOTE block in that proof).
\end{lemma}

\begin{proof}
  
  Direct definition; no proof body. The Lean construction reads
  \[
    \texttt{((pullbackPushforwardAdjunction}\,\psi).\texttt{homEquiv}\,M_G\,\texttt{LHS}).\texttt{symm}\,((\texttt{DifferentialsConstruction.isUniversal'}\,\varphi_G).\texttt{desc}\,D),
  \]
  where \(D\) is the derivation of
  \cref{lem:GrpObj_omega_basechange_proj_inv_derivation}. The iter-138
  Lean target \texttt{AlgebraicGeometry.GrpObj.basechange\_along\_proj\_two\_inv}
  is sorry-free as a \emph{definition} (the construction itself
  compiles end-to-end); the only remaining work for piece (i.b)
  closure is the iso property
  \(\texttt{IsIso}\,(\texttt{basechange\_along\_proj\_two\_inv}\,G)\),
  carried as a sub-sorry of
  \cref{lem:GrpObj_omega_basechange_proj}'s proof and attacked via
  Route (b'2) iter-140.
\end{proof}

\begin{lemma}\leanok
[(i.b-helper) Per-open IsIso of the inverse-direction morphism]
  \label{lem:GrpObj_basechange_along_proj_two_inv_app_isIso}
  \lean{AlgebraicGeometry.TODO.GrpObj_basechange_along_proj_two_inv_app_isIso}
  \uses{lem:GrpObj_omega_basechange_proj_inv, lem:isIso_of_app_iso_module}
  % Iter-146 disposition: \lean{...basechange_along_proj_two_inv_app_isIso}
  % stripped — Lean target EXCISED iter-145 (see EXCISE disposition
  % paragraph above \cref{lem:GrpObj_omega_basechange_proj}). Block kept
  % as auditable bundled-route history. (\leanok carried over from
  % iter-143 era is left for sync_leanok to reconcile against the
  % now-deleted target in the next sync run.)
  % Lean signature stub (iter-143 Wave 2 refactor extracted the
  % per-open IsIso obligation from an inline
  % `letI := isIso_of_app_iso_module ... (fun _ => sorry)` pattern
  % inside `relativeDifferentialsPresheaf_basechange_along_proj_two`
  % into a named sorry-bodied theorem in
  % `Cotangent/GrpObj.lean:745--751`. The refactor restores audit
  % transparency per the iter-143 STRATEGY.md Edit 1 sorry-must-be-
  % named-declaration discipline, and decouples the IsIso closure
  % from its consumer to enable an isolated iter-144+ prover round.
  % Body remains `sorry`; under the iter-144 chart-algebra pivot
  % this body closure is DESCOPED from critical-path (auditable
  % record only):
  %   theorem basechange_along_proj_two_inv_app_isIso
  %       (G : Over (Spec (.of k))) [CategoryTheory.GrpObj G] :
  %       ∀ X : (G ⊗ G).leftᵒᵖ,
  %         IsIso ((basechange_along_proj_two_inv G).app X)
  
  \notready
  Let \(G : \Over\,(\Spec k)\) be as in \cref{lem:GrpObj_omega_basechange_proj}, and let
  \[
    \texttt{basechange\_along\_proj\_two\_inv}\,G \;\colon\;
    (\texttt{PresheafOfModules.pullback}\,\psi).\mathrm{obj}\,\Omega_{G/k}
    \;\longrightarrow\; \Omega_{(G \times_{k} G)/G}
  \]
  denote the inverse-direction morphism of \cref{lem:GrpObj_omega_basechange_proj_inv} (with \(\psi\) the second-projection compatibility morphism of \cref{def:GrpObj_schemeHomRingCompatibility}). Then for every open \(X \in (G \otimes G).\mathrm{left}.\Opens^{\mathrm{op}}\), the application
  \[
    (\texttt{basechange\_along\_proj\_two\_inv}\,G).\app\,X
    \;\;\colon\;\;
    \bigl((\texttt{PresheafOfModules.pullback}\,\psi).\mathrm{obj}\,\Omega_{G/k}\bigr).\mathrm{obj}\,X
    \;\longrightarrow\; \Omega_{(G \times_{k} G)/G}.\mathrm{obj}\,X
  \]
  is an isomorphism of \(\Gamma((G \otimes G).\mathrm{left},\,\mathrm{snd}^{-1}\,X)\)-modules.

  Composed with the iso-reflection bridge \texttt{isIso\_of\_app\_iso\_module} (\texttt{Cotangent/GrpObj.lean:544--550}, closed iter-140; reflects pointwise iso-of-app to a presheaf-of-modules iso via \texttt{PresheafOfModules.toPresheaf} + \texttt{NatTrans.isIso\_iff\_isIso\_app}), this per-open conclusion supplies the third concrete sub-sorry of \cref{lem:GrpObj_omega_basechange_proj}'s proof, completing the base-change iso of piece (i.b) Step 2.
\end{lemma}

\begin{proof}
  
  Per Route (b'2) items 2--4 of the IsIso closure programme detailed inside \cref{lem:GrpObj_omega_basechange_proj}'s proof block (the \S "Route (b'2): local-iso check via \texttt{PresheafOfModules.toPresheaf} + \texttt{NatTrans.isIso\_iff\_isIso\_app}" NOTE; per the iter-139 mathlib-analogist verdict \texttt{analogies/isiso-basechange-along-proj-two-inv.md}, \textasciitilde\(195\)--\(365\) LOC), the per-open iso at each affine \(X = W \subseteq (G \otimes G).\mathrm{left}\) is constructed as the composition of three chart-level helpers:
  \begin{enumerate}
    \item \textbf{Chart-level \texttt{Algebra.IsPushout}-from-affine-product helper} (\textasciitilde\(80\)--\(150\) LOC; shared with Route (a) of the iter-137 5-step recipe). For an affine open \(W = \Spec B \subseteq (G \times_{k} G).\mathrm{left}\) cut out via the pullback square \(W = V_{1} \times_{U} V_{2}\) with affine opens \(V_{1}, V_{2} \subseteq G.\mathrm{left}\) over a common affine \(U \subseteq \Spec k\), the chart-level \texttt{Algebra.IsPushout} square chains three Mathlib facts: \texttt{CommRingCat.isPushout\_iff\_isPushout} (converts a \texttt{CommRingCat}-pushout into an \texttt{Algebra.IsPushout} of underlying algebras); \texttt{pullbackSpecIso} (the canonical iso \(\Spec B \cong \Spec B_{1} \times_{\Spec A} \Spec B_{2}\) identifying \(B\) with the ring-side tensor product \(B_{1} \otimes_{A} B_{2}\)); and \texttt{isPullback\_SpecMap\_of\_isPushout} (\(\Spec\) sends pushout squares of rings to pullback squares of schemes; used in the converse direction here).
    \item \textbf{The \texttt{pullbackObjEquivTensor} chart-unfolding helper} (\textasciitilde\(30\)--\(60\) LOC; shared with Route (a)). Unfolds the section of the pulled-back presheaf-of-modules \(((\texttt{PresheafOfModules.pullback}\,\psi).\mathrm{obj}\,M).\mathrm{obj}\,X\) as the restricted tensor product \(\Gamma(R, X) \otimes_{((F^{\mathrm{op}})_{*}\,S).\mathrm{obj}\,X} M.\mathrm{obj}(F^{\mathrm{op}}.\mathrm{obj}\,X)\) via the unit/counit of \texttt{PresheafOfModules.pullbackPushforwardAdjunction} (\texttt{Mathlib/Algebra/Category/ModuleCat/Presheaf/Pullback.lean:50}). The helper exposes the chart-level tensor shape that the iso check (item 3) requires; it does not escape the \texttt{pullback}-opacity issue but bypasses the need to manipulate the abstract left-adjoint definition.
    \item \textbf{Per-open identification with \texttt{tensorKaehlerEquiv.symm}} (\textasciitilde\(80\)--\(150\) LOC). With items 1+2 in hand, the per-open iso \(((\texttt{PresheafOfModules.pullback}\,\psi).\mathrm{obj}\,\Omega_{G/k}).\mathrm{obj}\,X \xrightarrow{\sim} \Omega_{(G \times_{k} G)/G}.\mathrm{obj}\,X\) for affine \(X = W\) is the value of \texttt{KaehlerDifferential.tensorKaehlerEquiv.symm} (\texttt{Mathlib.RingTheory.Kaehler.TensorProduct}; for an \texttt{Algebra.IsPushout R S A B} square, yields \(B \otimes_{A} \Omega[A/R] \simeq_{B} \Omega[B/S]\)) applied to the chart-level \texttt{Algebra.IsPushout} square from item 1, re-presented through the chart-unfolding equivalence \texttt{pullbackObjEquivTensor} of item 2. The algebra-side value identity \(D\,b \mapsto 1 \otimes D\,b\) is packaged by the companion simp lemma \texttt{tensorKaehlerEquiv\_symm\_D\_tmul}.
  \end{enumerate}

  \paragraph{Iter-144 chart-algebra pivot --- disposition.} Under the iter-144 chart-algebra pivot (see the \S "Iter-144 chart-algebra pivot --- disposition" paragraph at the head of \cref{subsec:RigidityKbar_piece_i_decomposition}), this lemma is \textbf{DESCOPED from critical-path}; its named Lean declaration \texttt{AlgebraicGeometry.GrpObj.basechange\_along\_proj\_two\_inv\_app\_isIso} remains in-tree as a sorry-bodied auditable record of the bundled Route (b'2) closure programme, but the iter-145+ prover lane targets the chart-algebra (\(\alpha\)) + (\(\beta\)) helpers of piece (ii) instead (see the \S "Iter-144 chart-algebra envelope for piece (ii)" paragraph for the chart-algebra envelope). The Route (b'2) sketch above is preserved for traceability and is the prover-ready recipe should a future strategy reversal re-prioritise the bundled route.
\end{proof}

\begin{lemma}[(i.b-helper) Cotangent at the identity via restriction along the canonical section]
  \label{lem:GrpObj_omega_restrict_to_identity_section}
  \lean{AlgebraicGeometry.GrpObj.relativeDifferentialsPresheaf_restrict_along_identity_section}
  \uses{def:relative_kaehler_presheaf, lem:section_snd_eq_identity_struct}
  % Lean signature stub (iter-134+ helper for piece (i.b); ~30-80 LOC per
  % the iter-133 mathlib-analogist verdict step 3 of the piece (i.b) LOC
  % envelope at
  % `task_results/mathlib-analogist-mulright-globalises-iter133.md`):
  %   noncomputable def relativeDifferentialsPresheaf_restrict_along_identity_section
  %       (G : Over (Spec (.of k))) [CategoryTheory.GrpObj G]
  %       {n : ℕ} [SmoothOfRelativeDimension n G.hom]
  %       [IsProper G.hom] [GeometricallyIrreducible G.hom] :
  %       (PresheafOfModules.pullback (φ_section G)).obj
  %           ((PresheafOfModules.pullback (φ_pr_two G)).obj
  %             (relativeDifferentialsPresheaf G.hom)) ≅
  %         (PresheafOfModules.pullback (φ_str G)).obj
  %           ((PresheafOfModules.pullback (φ_η G)).obj
  %             (relativeDifferentialsPresheaf G.hom))
  % where `φ_section`/`φ_pr_two`/`φ_str`/`φ_η` are the compatibility
  % morphisms for `s = ⟨𝟙_G, η_G⟩`, `pr_2`, `π_G : G ⟶ Spec k`, and `η_G`
  % respectively; the iso witnesses that the two ways of composing pullbacks
  % implementing `pr_2 ∘ s = η_G ∘ π_G` agree.
  
  % NOTE iter-136 review: the Lean target
  % \lean{AlgebraicGeometry.GrpObj.relativeDifferentialsPresheaf_restrict_along_identity_section}
  % is now substantively closed (no `sorry`; kernel-only axioms). The proof
  % uses `PresheafOfModules.pullbackComp` on both sides plus a new private
  % helper `section_snd_eq_identity_struct` (~5 LOC; captures the categorical
  % identity `s.left ≫ pr_2.left = G.hom ≫ η[G].left` at the underlying
  % scheme level via `lift_snd` + `Over.toUnit_left`). The four compatibility
  % morphisms for `PresheafOfModules.pullback` along `s`, `pr_2`, `π_G`,
  % `η_G` are constructed inline via Mathlib's `Scheme.Hom.toRingCatSheafHom`,
  % per the iter-135 mathlib-analogist verdict in
  % `analogies/phi-compatibility-morphisms.md`. Total ~27 LOC closure
  % (helper 5 LOC + body 22 LOC), comfortably inside the predicted
  % \(\sim 30\)--\(80\) LOC envelope. `sync_leanok` will (re-)add `\leanok` to
  % the proof block accordingly. Prior `\notready` removed.
  Let \(s := \langle \mathrm{id}_{G},\, \eta_{G} \rangle \colon G \to G \times_{k} G\) denote the canonical section of the first projection \(\pr_{1}\) in \(\Over\,(\Spec k)\), and let \(\pi_{G} \colon G \to \Spec k\) denote the structure map. The categorical identity \(\pr_{2} \circ s = \eta_{G} \circ \pi_{G}\) in \(\Over\,(\Spec k)\) — which holds by direct evaluation of the binary-product second projection on the \(\langle \mathrm{id}_{G}, \eta_{G} \rangle\) pair — yields a canonical natural isomorphism of presheaves of modules on \(G\):
  \[
    s^{*}\!\bigl(\pr_{2}^{*}\,\Omega_{G/k}\bigr) \;\cong\; \pi_{G}^{*}\!\bigl(\eta_{G}^{*}\,\Omega_{G/k}\bigr).
  \]
\end{lemma}

\begin{proof}
  
  \leanok
  The categorical identity \(\pr_{2} \circ s = \eta_{G} \circ \pi_{G}\) in \(\Over\,(\Spec k)\) holds by direct evaluation: \(s = \langle \mathrm{id}_{G}, \eta_{G} \rangle\) has second coordinate \(\eta_{G}\) (viewed as a \(\Spec k\)-valued point of \(G\), pre-composed with the structure map \(\pi_{G} \colon G \to \Spec k\) to obtain a \(G\)-morphism to \(G\)), so post-composing with \(\pr_{2}\) extracts that coordinate. Applying the Mathlib functoriality lemma \texttt{PresheafOfModules.pullbackComp} (\texttt{Mathlib.Algebra.Category.ModuleCat.Presheaf.Pullback}; for composable morphisms \(f, g\) in the underlying category, \((f \circ g)^{*} \cong g^{*} \circ f^{*}\) as functors on \texttt{PresheafOfModules}) to both sides of the composition identity, applied to the presheaf-of-modules \(\Omega_{G/k} = \texttt{relativeDifferentialsPresheaf}\,G.\mathrm{hom}\), yields the displayed iso. Estimated build: \(\sim 30\)--\(80\) LOC.
\end{proof}

\begin{lemma}[(i.c) Freeness of the relative cotangent of \(G\)]
  \label{lem:GrpObj_omega_free}
  \lean{AlgebraicGeometry.TODO.GrpObj_omega_free}
  \uses{lem:GrpObj_mulRight_globalises, lem:GrpObj_cotangentSpace}
  % Iter-146 disposition: \lean{AlgebraicGeometry.GrpObj.omega_free}
  % stripped — DESCOPED iter-144 chart-algebra pivot; no Lean
  % declaration of that name ever materialised in tree (chart-algebra
  % routes the freeness fact through chart-level
  % \texttt{Algebra.IsStandardSmooth.free\_kaehlerDifferential} per-chart,
  % not through a global sheaf-of-modules construction). Block kept as
  % traceable iter-126 \(\to\) iter-144 decision history; remains \notready.
  
  \notready
  The relative cotangent presheaf \(\Omega_{G/k}\) is a free \(\mathcal O_{G}\)-module.
\end{lemma}

\begin{proof}
  
  By \cref{lem:GrpObj_mulRight_globalises}, \(\Omega_{G/k}\) is canonically isomorphic to the pullback along the structure map \(G \to \Spec k\) of the \(k\)-vector space \(\mathfrak g^{\vee} = \eta_{G}^{*}\,\Omega_{G/k}\). By \cref{lem:GrpObj_cotangentSpace}, \(\mathfrak g^{\vee}\) is a finitely generated free \(k\)-module. The pullback of a finitely generated free \(k\)-module along the structure morphism is a finitely generated free \(\mathcal O_{G}\)-module.
\end{proof}

\begin{lemma}[(i.c) Rank of the relative cotangent of \(G\)]
  \label{lem:GrpObj_omega_rank_eq_dim}
  \lean{AlgebraicGeometry.TODO.GrpObj_omega_rank_eq_dim}
  \uses{lem:GrpObj_omega_free, lem:GrpObj_lieAlgebra_finrank}
  % Iter-146 disposition: \lean{AlgebraicGeometry.GrpObj.omega_rank_eq_dim}
  % stripped — DESCOPED iter-144 chart-algebra pivot (same reason as
  % \cref{lem:GrpObj_omega_free}); no Lean declaration of that name
  % materialised in tree. Block kept as traceable history; remains
  % \notready.
  
  \notready
  The (constant) rank of \(\Omega_{G/k}\) as a free \(\mathcal O_{G}\)-module equals \(n\), where \(n\) is the relative dimension supplied by \([\mathtt{SmoothOfRelativeDimension}\,n\,G.\mathtt{hom}]\).
\end{lemma}

\begin{proof}
  
  The rank of the pulled-back bundle equals the \(k\)-dimension of its fibre. The fibre is \(\mathfrak g^{\vee}\), whose \(k\)-rank equals \(n\) by \cref{lem:GrpObj_lieAlgebra_finrank}.
\end{proof}

\begin{remark}
  \label{rem:piece_i_first_target}
  
  Sub-step (i.a) --- \cref{lem:GrpObj_cotangentSpace} (the definition) plus the rank lemma \cref{lem:GrpObj_lieAlgebra_finrank} --- is the natural first target for the iter-128+ prover lane (closed iter-128 \(\to\) iter-132 over a 5-iter span; see STRATEGY.md \S Sequencing piece (i.a) row for the empirical cost trace). Under Replacement (B) the iter-131 strategy-critic Q4 trio\(\to\)duo collapse demoted the bridge lemma \cref{lem:GrpObj_cotangent_bridge} to a vestigial-on-live-path deferred alternative (see rank-lemma proof Step 3 + footer at \cref{lem:GrpObj_lieAlgebra_finrank}). Sub-steps (i.b) and (i.c) are staged on top of (i.a) and are scheduled for iter-133+ / iter-137+ respectively.
\end{remark}

\paragraph{Iter-131 \texttt{Classical.choose}-chain body shape (Lean encoding note).} The iter-131 fix-up refactored the body of \texttt{AlgebraicGeometry.GrpObj.cotangentSpaceAtIdentity} from an iter-130 \texttt{by}-tactic body that left \texttt{Classical.choice} as its outermost elaborated head symbol — opaque to downstream \texttt{unfold}/\texttt{whnf} — into a pure-term \texttt{noncomputable def}. Concretely: a \texttt{let}-chain of \texttt{Classical.choose} / \texttt{Classical.choose\_spec} on \cref{thm:smooth_locally_free_omega}'s existential introduces the affine charts \(U \subseteq \Spec k\) and \(V \subseteq G.\mathrm{left}\) (around the identity-section image \(x_0 = \eta_G(\mathrm{pt})\)), the appLE inclusion \(e : V \leq G.\mathrm{hom}^{-1}\,U\), the membership witness \(h_{xV} : x_0 \in V\), and the ring map \(\psi_V : \Gamma(G, V) \to k\) (identity section restricted to \(V\), composed with \texttt{Scheme.\(\Gamma\)SpecIso}); the outer expression of the body is then the explicit
\[
  \texttt{(ModuleCat.extendScalars\,\(\psi_V\).hom).obj\,(ModuleCat.of\,\(\Gamma\)(G, V)\,\(\Omega\)[\(\Gamma\)(G, V)\,/\,\(\Gamma\)(Spec k, U)])},
\]
with no top-level \texttt{Classical.choice} head symbol. The companion lemma \cref{lem:GrpObj_cotangentSpace_extendScalars_witness} (\texttt{cotangentSpaceAtIdentity\_eq\_extendScalars}) witnesses this structural shape via an existential (the chart triple \((U, V, e)\) plus a top-inclusion \(h_{\mathrm{top}}\) and the explicit equation displayed in that lemma), closed essentially by \texttt{rfl} once the same \texttt{Classical.choose}-chain is unpacked on the right-hand side.

Two downstream rewrite patterns are available, both mathematically equivalent (the body's underlying \texttt{TensorProduct} carrier and the \texttt{extendScalars}-of-\texttt{ModuleCat.of} form share the same \texttt{Classical.choose} extraction, hence are definitionally equal once the \texttt{ModuleCat.ExtendScalars.obj'} carrier identity is unfolded):

\begin{itemize}
  \item \textbf{Direct \texttt{change}-based route (primary; used by the iter-132 rank-lemma closure).} The iter-132 close of \cref{lem:GrpObj_lieAlgebra_finrank} (in-tree as \texttt{AlgebraicGeometry.GrpObj.cotangentSpaceAtIdentity\_finrank\_eq}) uses a direct \texttt{change Module.finrank k (TensorProduct \(\Gamma\)(G, V) k \(\Omega\)[\ldots]) = n} step that exposes the body's underlying \texttt{TensorProduct} carrier via the definitional unfolding of \texttt{ModuleCat.extendScalars} (whose \texttt{obj'} is \texttt{TensorProduct \(\Gamma\)(G, V) k M} up to \texttt{ModuleCat} coercion, per Mathlib's \texttt{Mathlib.Algebra.Category.ModuleCat.ChangeOfRings:396}). The chart witnesses \(U, V, e, h_{xV}, \texttt{hfree}, \texttt{hrank}\) are extracted directly by re-running the same \texttt{Classical.choose}-chain on \cref{thm:smooth_locally_free_omega} inside the consumer's proof body (and \(\psi_V\) reconstructed via \texttt{appLE} composition with \texttt{Scheme.\(\Gamma\)SpecIso}). The closure then reads \texttt{rw [Module.finrank\_baseChange]; exact Module.finrank\_eq\_of\_rank\_eq hrank}. \textbf{This is the shorter route, and what the rank lemma's iter-132 body actually does.} It is the recommended pattern for downstream consumers whose goal can be massaged via \texttt{change} into the underlying \texttt{TensorProduct} carrier.
  \item \textbf{\texttt{obtain}+\texttt{rw [heq]} alternative (preserved for consumers needing an explicit rewrite handle).} The companion lemma \cref{lem:GrpObj_cotangentSpace_extendScalars_witness} supports a fielded destructure \texttt{obtain \(\langle\)U, V, e, \(h_\mathrm{top}\), heq\(\rangle\) := cotangentSpaceAtIdentity\_eq\_extendScalars G}, after which \texttt{rw [heq]} exposes the explicit \texttt{extendScalars} form to the goal; the chart-side \texttt{hfree}/\texttt{hrank} are re-extracted on the same \texttt{Classical.choose}-chain from \cref{thm:smooth_locally_free_omega} and fed into \texttt{Module.finrank\_baseChange}. This route makes the rewrite step explicit and is preferred when the consumer's goal does not unify with the body's underlying \texttt{TensorProduct} carrier via \texttt{change} alone (e.g., when the goal is phrased at the \texttt{ModuleCat k}-object level rather than at the underlying-type level).
\end{itemize}

Note: an alternative interface that would let consumers destructure with a single \texttt{let} (a \texttt{\(\Sigma'\)}-bundled \texttt{smooth\_locally\_free\_omega'} returning the chart data as a sigma type rather than an existential) is a deferred cleanup opportunity flagged in \texttt{analogies/cotangent-body-shape.md} \S "Future cleanup"; it is out of scope until \texttt{Differentials.lean} comes off the do-not-touch list.

\subsection{Chart-algebra piece (ii) first-class decomposition}
\label{subsec:RigidityKbar_piece_ii_chart_algebra_decomposition}

The five-bullet envelope of \S "Iter-144 chart-algebra envelope for piece (ii)" at the head of \cref{subsec:RigidityKbar_piece_i_decomposition} (post-pivot piece (ii) total \textasciitilde\(600\)--\(1050\) LOC) is lifted here into five first-class declaration blocks, with pre-committed Lean target names supplied by the iter-145 planner. The dependency order between the blocks is bottom-up: sub-piece (\(\alpha\)) supplies the chart-level \texttt{Algebra.IsPushout} square; sub-pieces (\(\beta\)-aux integrally-closed-constants) and (\(\beta\)-core ring-side) feed into (\(\beta\)-core); the (\(\beta\)-core) per-chart helper is consumed by the scheme-level lift. The blocks below cite each other via \texttt{\textbackslash uses} as needed.

\paragraph{Strategy-critic Q3 absorption (iter-145).} The iter-145 strategy-critic flagged an internal contradiction in the M2.a \(\mathrm df = 0\) derivation chain articulated in STRATEGY.md (and partly mirrored in the chart-algebra envelope above): the chain reads as both \emph{using} a global-section-vanishing argument on \(\Omega_{C/k}\) (mathematically equivalent to \(H^0(C, \Omega_{C/k}) = 0\) on a genus-\(0\) curve) and \emph{not} invoking that vanishing. The honest articulation, recorded in this subsection's load-bearing block (\cref{lem:chart_algebra_df_zero_factors_through_constant_on_chart}), is option (a) of the iter-145 strategy-critic's Q3 dichotomy: the chart-algebra (\(\beta\)-core) helper invokes the cohomological content \(H^0(C, \Omega_{C/k}^{\oplus n}) = 0\) on a genus-\(0\) curve \(C\) via a two-chart \v Cech / Mayer--Vietoris computation directly on \(\Omega_{C/k}^{\oplus n}\), reusing the project's existing chart-\v Cech Mayer--Vietoris infrastructure (\cref{chap:Cohomology_MayerVietoris}, in particular \cref{thm:Scheme_AffineCoverMVSquare_HModule_prime_sequence_exact}; already consumed by \texttt{Genus.lean}'s \(H^1(C, \mathcal O_C) = 0\) computation on a genus-\(0\) curve, run here on \(\Omega_{C/k}^{\oplus n}\) instead of \(\mathcal O_C\)). The named theorem ``Serre duality on a smooth proper curve'' (piece (iv) of \cref{sec:RigidityKbar_shared_pile}) is \emph{not} invoked as a packaged citation; the disclaimer ``without Serre duality'' in the iter-144 chart-algebra envelope (third item of the (\(\beta\)) three-layer chain) narrows to ``the named Serre-duality theorem is not invoked; the cohomological content \(H^0(C, \Omega_{C/k}) = 0\) on a genus-\(0\) curve IS invoked, via a two-chart \v Cech computation matching the project's \(H^1(C, \mathcal O_C)\) idiom''. Piece (iv) remains a deferred named-gap for any future M2.d-alt route that would invoke Serre duality by name, but is not on the critical path for the chart-algebra (\(\beta\)-core) build.

\emph{Iter-155 scope clarification.} The statement ``piece (iv) is not on the critical path for the \(\beta\)-core build'' is true only because the \(\beta\)-core build (\cref{lem:chart_algebra_df_zero_factors_through_constant_on_chart}) is the \textbf{converse} ``\(\mathrm df = 0 \Rightarrow\) constant'' and takes \(\mathrm df = 0\) as a hypothesis. It does \emph{not} entail that the chapter's keystone (\cref{thm:rigidity_over_kbar}) avoids Serre duality: \emph{producing} \(\mathrm df = 0\) is a separate, gated step that --- per \texttt{analogies/df-zero-production-iter155.md} --- is a global-sections fact equivalent (given the gap-(i) trivialisation \(f^{*}\Omega_{A/k} \cong \mathcal O_C^{\oplus g}\)) to \(H^0(C, \Omega_{C/k}) = 0\), i.e.\ exactly piece (iv). The chart-by-chart machinery cannot detect \(\mathrm df = 0\) because \(\Omega_{C/k}(V)\) is free of rank \(1\) (nonzero) on every affine chart \(V\). So \(H^0(C, \Omega_{C/k}) = 0\) \textbf{is} on the critical path for the keystone (under route (a)); see the corrected chapter-introduction disposition and \cref{sec:RigidityKbar_shared_pile} piece (iv).

\begin{lemma}\leanok
[(ii-\(\alpha\)) Algebra.IsPushout from the affine product of two charts]
  \label{lem:chart_algebra_isPushout_of_affine_product}
  \lean{AlgebraicGeometry.GrpObj.algebra_isPushout_of_affine_product}
  % Iter-146: stray `` stripped — the
  % statement is purely algebra-level (commutative-ring pushout square
  % on the chart presentations), with no relative Kähler presheaf in
  % either the statement or the proof. The dependency-graph noise is
  % cleaned up here.
  % NOTE (iter-146 review): the iter-146 Lean closure discharges this
  % lemma via a single `inferInstance` after re-enabling Mathlib's
  % local `Algebra.TensorProduct.rightAlgebra` instance — the
  % `TensorProduct.isPushout` instance in `Mathlib.RingTheory.IsTensorProduct`
  % collapses the three-step `pullbackSpecIso` / `isPullback_SpecMap_of_isPushout`
  % / `CommRingCat.isPushout_iff_isPushout` chain into one resolution
  % step at the algebra level. The three-step chain remains the honest
  % derivation at the scheme level; the algebra-level lemma is what
  % the chart-algebra consumers (β-core + KDM + lift) actually need.
  % NOTE (iter-147 review): iter-147 prover lane carry-over — the
  % `inferInstance` closure stands unchanged; the entire content of
  % this lemma is now an instance-resolution wrapper over Mathlib's
  % `Algebra.IsPushout` infrastructure. The proof block is
  % informational (the explicit three-step chain remains the honest
  % blueprint derivation at the scheme level) rather than
  % load-bearing for the Lean side.
  Let \(X_1, X_2\) be schemes over \(\Spec k\), and let \(U_1 \subseteq X_1\), \(U_2 \subseteq X_2\) be affine open subsets with \(U_1 = \Spec B_1\) and \(U_2 = \Spec B_2\). Let \(W \subseteq X_1 \times_k X_2\) denote the affine open obtained as the fibre product \(U_1 \times_k U_2\), identified with \(\Spec B\) where \(B = B_1 \otimes_k B_2\) via the canonical chart presentation. Then the square of commutative ring maps
  \[
    \begin{array}{ccc}
      k & \longrightarrow & B_1 \\
      \downarrow & & \downarrow \\
      B_2 & \longrightarrow & B
    \end{array}
  \]
  carries an \texttt{Algebra.IsPushout}\,\(k\)\,\(B_1\)\,\(B_2\)\,\(B\) structure: \(B\) is the algebra-level tensor product of \(B_1\) and \(B_2\) over \(k\), packaged as a pushout square in \(\mathtt{CommRingCat}\) / \texttt{Algebra}.
\end{lemma}

\begin{proof}

\leanok
  Three Mathlib facts chain to deliver the pushout structure:
  \begin{enumerate}
    \item \emph{Affine pullback identifies \(W\) with \(\Spec(B_1 \otimes_k B_2)\).} The Mathlib lemma \texttt{pullbackSpecIso} (\texttt{Mathlib.AlgebraicGeometry.Pullbacks}) supplies the canonical iso \(\Spec(B_1 \otimes_k B_2) \cong \Spec B_1 \times_{\Spec k} \Spec B_2\) as schemes, together with the simp companions \texttt{pullbackSpecIso\_hom\_fst} / \texttt{\_hom\_snd} / \texttt{\_inv\_fst} / \texttt{\_inv\_snd} characterising its action against the two projections. Composing with the chart inclusion \(W \hookrightarrow X_1 \times_k X_2\) identifies \(\Gamma(X_1 \times_k X_2, W) \cong B_1 \otimes_k B_2\) as a \(k\)-algebra; since both \(W\) and \(\Spec(B_1 \otimes_k B_2)\) are affine and their structure-sheaf global sections agree, \(W \cong \Spec(B_1 \otimes_k B_2)\) as schemes over \(\Spec k\).
    \item \emph{\(\Spec\) sends ring pushout squares to scheme pullback squares.} The Mathlib lemma \texttt{isPullback\_SpecMap\_of\_isPushout} (\texttt{Mathlib.AlgebraicGeometry.Pullbacks}) gives the forward direction; in step 1 we already exhibited the pullback square on the scheme side, and the converse direction (a pullback of affine spectra over \(\Spec k\) is the \texttt{Spec} of a \(\mathtt{CommRingCat}\)-pushout square of the underlying rings) follows by reading the iso of step 1 back through \(\Spec\)'s contravariant action on pushout/pullback squares.
    \item \emph{Bridge from \(\mathtt{CommRingCat}\)-pushout to \texttt{Algebra.IsPushout}.} The lemma \texttt{CommRingCat.isPushout\_iff\_isPushout} (\texttt{Mathlib.Algebra.Category.Ring.Pushout}) translates the \(\mathtt{CommRingCat}\)-pushout statement of step 2 into the \texttt{Algebra.IsPushout}\,\(k\)\,\(B_1\)\,\(B_2\)\,\(B\) statement on the underlying \(k\)-algebras (the underlying class is \texttt{Algebra.IsPushout} from \texttt{Mathlib.RingTheory.IsTensorProduct}, the iter-145-verified Mathlib home for \texttt{Algebra.IsPushout} per the NOTE in \texttt{AlgebraicJacobian/Cotangent/ChartAlgebra.lean:10--15}; an earlier blueprint draft cited the file as \texttt{Mathlib.RingTheory.IsPushout}, which does not exist upstream).
  \end{enumerate}
  Composing the three steps yields the desired \texttt{Algebra.IsPushout} structure. The recipe is shared verbatim with item~(1) of the iter-137 Route~(a) chart-unfolding 5-step recipe inside \cref{lem:GrpObj_omega_basechange_proj}'s proof block, and with item~2 of the iter-143 Route~(b'2) per-open IsIso closure recipe inside \cref{lem:GrpObj_basechange_along_proj_two_inv_app_isIso}'s proof; reusing it here keeps the chart-algebra (\(\alpha\)) helper's LOC envelope tight. Estimated build: \(\sim 80\)--\(150\) LOC.

  \emph{Mathlib status.} All three component lemmas are [verified] in Mathlib snapshot \texttt{b80f227} by the iter-145 M3 audit refresh: \texttt{pullbackSpecIso} and \texttt{isPullback\_SpecMap\_of\_isPushout} in \texttt{Mathlib.AlgebraicGeometry.Pullbacks}; \texttt{CommRingCat.isPushout\_iff\_isPushout} in \texttt{Mathlib.Algebra.Category.Ring.Pushout}; the underlying \texttt{Algebra.IsPushout} class in \texttt{Mathlib.RingTheory.IsTensorProduct} (the iter-145-verified Mathlib home; the earlier draft's citation of \texttt{Mathlib.RingTheory.IsPushout} pointed at a file that does not exist upstream).
\end{proof}

\begin{theorem}\leanok
[(ii-\(\beta\)-core) Per-chart translation-invariance K\"ahler-derivation: \(\mathrm df = 0\) factors through the base field]
  \label{lem:chart_algebra_df_zero_factors_through_constant_on_chart}
  \lean{AlgebraicGeometry.GrpObj.df_zero_factors_through_constant_on_chart}
  \uses{lem:chart_algebra_isPushout_of_affine_product,
    lem:KaehlerDifferential_mem_range_algebraMap_of_D_eq_zero,
    lem:constants_integral_over_base_field,
    thm:Scheme_AffineCoverMVSquare_HModule_prime_sequence_exact}
  
  % NOTE (iter-148 review): the iter-147 Lean closure of this theorem
  % is a structurally simpler one-line delegate to
  % `\lem:KaehlerDifferential_mem_range_algebraMap_of_D_eq_zero` (KDM).
  % Concretely, the Lean signature drops the chart-pair data
  % \((W = \Spec B, V = \Spec R, f^{\sharp} : B \to R, A, f : C \to A, \mathrm df = 0)\)
  % and keeps only $(k, C \text{ with chart-of-proper-curve typeclasses},
  % B \text{ a finite-type } k\text{-algebra}, b \in B, hDb : D b = 0)$.
  % The four typeclasses on \(C\) (\texttt{IsProper}, \texttt{Smooth},
  % \texttt{IsReduced}, \texttt{GeometricallyIrreducible}) are
  % decorative under the current commitment — the body is a one-line
  % wrapper around KDM and consumes none of them. Mirroring the
  % iter-146/iter-147 \texttt{\% NOTE} discipline on
  % `\lem:Scheme_Over_ext_of_diff_zero`: this is a thin-wrapper
  % disposition pending substantive Steps 1, 2, 3, 5 refinement.
  % Iter-149+ refinement plan: inflate the Lean signature to also
  % carry \((A, f : C \to A, \mathrm df = 0, W, V, f^{\sharp})\), derive
  % the chart-local kernel vanishing of Steps 1--2 from \(\mathrm df = 0\),
  % run the 2-chart \v Cech Mayer--Vietoris promotion of Step 3, and
  % delegate to KDM + integrally-closed-constants in Steps 4--5.
  % The five-step recipe below remains the load-bearing target;
  % the current Lean commits only to its terminal Steps 4 + KDM call.
  % NOTE (iter-152 review): the iter-152 alg-closed pivot propagated
  % `[IsAlgClosed k]` + `[IsDomain B]` onto this theorem's Lean
  % signature (so the one-line delegation to KDM, which now requires
  % those hypotheses, typechecks). The statement-block NOTE history
  % above is otherwise frozen at iter-148; mathematically the new
  % hypotheses are forced by the KDM delegation, not an independent
  % strengthening.
  % NOTE (iter-154 review): KDM
  % (`\lem:KaehlerDifferential_mem_range_algebraMap_of_D_eq_zero`) was
  % CLOSED iter-154, axiom-clean. This delegate is therefore now
  % genuinely axiom-clean too: `lean_verify` on
  % `df_zero_factors_through_constant_on_chart` reports axioms
  % `{propext, Classical.choice, Quot.sound}` — no `sorryAx`. The
  % iter-152 NOTE here previously warned that this theorem laundered a
  % transitive `sorryAx` from KDM's open `sorry`; that warning is now
  % obsolete and has been removed. Both the result and its delegate are
  % fully discharged.
  Let \(C : \Over\,(\Spec k)\) be a smooth proper geometrically irreducible curve of genus \(0\), let \(A : \Over\,(\Spec k)\) be a smooth proper geometrically irreducible group scheme over \(k\) with identity element \(\eta_A \in A(k)\), and let \(f \colon C \to A\) be a morphism in \(\Over\,(\Spec k)\) with \(\mathrm df = 0\) (the scheme-level cotangent pull-back \(f^{*}\,\Omega_{A/k} \to \Omega_{C/k}\) is the zero morphism of presheaves of \(\mathcal O_C\)-modules). Fix:
  \begin{itemize}
    \item an affine chart \(W = \Spec B \subseteq A.\mathrm{left}\) around the image of \(\eta_A\) on which \(\Omega_{A/k}\) is presented as a free \(B\)-module of rank \(g = \dim A\) via the standard-smooth presentation (Mathlib's \texttt{Algebra.IsStandardSmoothOfRelativeDimension.rank\_kaehlerDifferential}), with chart generators \(b_1, \ldots, b_g \in B\) whose universal differentials \(\mathrm d b_1, \ldots, \mathrm d b_g\) form a \(B\)-basis of \(\Omega_{B/k}\);
    \item an affine chart \(V = \Spec R \subseteq C.\mathrm{left}\) contained in \(f^{-1}(W)\) on which the chart-restriction of \(f\) yields a \(k\)-algebra map \(f^{\sharp} \colon B \to R\).
  \end{itemize}
  Then for every \(b \in B\),
  \[
    f^{\sharp}(b) \;\in\; \operatorname{range}(\mathtt{algebraMap}\,k\,R) \;\subseteq\; R.
  \]
  Equivalently, the chart-level ring map \(f^{\sharp} \colon B \to R\) factors as \(B \to k \to R\).

  \paragraph{Strategy-critic Q3 honesty note.} The proof articulated below invokes the cohomological content \(H^0(C, \Omega_{C/k}^{\oplus g}) = 0\) on the genus-\(0\) curve \(C\) via a two-chart \v Cech / Mayer--Vietoris computation directly on \(\Omega_{C/k}^{\oplus g}\) (Step 3 of the proof), reusing the project's existing Mayer--Vietoris infrastructure on schemes (\cref{chap:Cohomology_MayerVietoris}; in particular \cref{thm:Scheme_AffineCoverMVSquare_HModule_prime_sequence_exact} for the affine-cover Mayer--Vietoris sequence). The named theorem ``Serre duality on a smooth proper curve'' (piece (iv) of \cref{sec:RigidityKbar_shared_pile}) is \emph{not} invoked as a packaged citation. Mathematically the content \(H^0(C, \Omega_{C/k}) = 0\) on a genus-\(0\) curve coincides with the Serre dual of \(H^1(C, \mathcal O_C) = 0\), but the proof here runs the chart-\v Cech computation by hand on \(\Omega_{C/k}^{\oplus g}\) using the same idiom that the project applies to compute \(H^1(C, \mathcal O_C) = 0\) in \texttt{Genus.lean}, rather than naming Serre duality. The ``without Serre duality'' disclaimer of the iter-144 chart-algebra envelope is therefore correct at the named-theorem level; it does \emph{not} mean ``without using \(H^0(C, \Omega_{C/k}) = 0\) at all'', and the proof must articulate the chart-\v Cech invocation explicitly (Step 3 below).
\end{theorem}

\begin{proof}
  
  \leanok
  The proof proceeds in five steps. Steps 1--2 establish that the chart-level universal K\"ahler derivation \(\mathrm dB \colon R \to \Omega_{R/k}\) satisfies \(\mathrm dB(f^{\sharp}(b)) = 0\) for every \(b \in B\), locally on a single affine chart \(V\), from the hypothesis \(\mathrm df = 0\) and the standard-smooth presentation of \(\Omega_{A/k}\) on \(W\). Step 3 promotes this chart-local kernel-vanishing to global section-vanishing on \(C\) via the chart-\v Cech Mayer--Vietoris layer. Steps 4--5 extract the conclusion \(f^{\sharp}(b) \in \operatorname{range}(\mathtt{algebraMap}\,k\,R)\) via the ring-side helper (\cref{lem:KaehlerDifferential_mem_range_algebraMap_of_D_eq_zero}) and the integrally-closed-constants helper (\cref{lem:constants_integral_over_base_field}).

  \emph{Step 1: chart-level translation of \(\mathrm df = 0\).} The hypothesis \(\mathrm df = 0\) on the chart pair \((W, V)\) is the statement that the induced map of \(R\)-modules
  \[
    R \,\otimes_B\, \Omega_{B/k} \;\longrightarrow\; \Omega_{R/k}, \qquad r \otimes \mathrm db \;\mapsto\; r \cdot \mathrm dB(f^{\sharp}(b)),
  \]
  is the zero map (the chart-level pullback \(f^{\sharp,*}\,\Omega_{A/k} \to \Omega_{C/k}\), packaged via the algebra-side base change of K\"ahler modules along \(f^{\sharp}\); here \(\mathrm dB \colon R \to \Omega_{R/k}\) is the universal K\"ahler derivation of \(R\) over \(k\)). Since \(\Omega_{B/k}\) is a free \(B\)-module of rank \(g\) with basis \(\mathrm db_1, \ldots, \mathrm db_g\) via the standard-smooth presentation on \(W\) (Mathlib's \texttt{Algebra.IsStandardSmooth.free\_kaehlerDifferential} and \texttt{Algebra.IsStandardSmoothOfRelativeDimension.rank\_kaehlerDifferential} of \texttt{Mathlib.RingTheory.Smooth.StandardSmoothCotangent}), the zero-map hypothesis unpacks as the system
  \[
    \mathrm dB(f^{\sharp}(b_i)) \;=\; 0 \qquad \text{in } \Omega_{R/k}, \qquad i = 1, \ldots, g.
  \]

  \emph{Step 2: extend to all \(b \in B\).} For arbitrary \(b \in B\), the universal K\"ahler derivation \(\mathrm db \in \Omega_{B/k}\) has a \(B\)-linear expansion \(\mathrm db = \sum_{i=1}^{g} c_i(b)\, \mathrm db_i\) with coefficients \(c_i(b) \in B\) obtained from the standard-smooth presentation. The functoriality of the K\"ahler derivation along the ring map \(f^{\sharp} \colon B \to R\) (Mathlib's \texttt{KaehlerDifferential.map} / \texttt{relativeDifferentials'\_map\_d}) gives \(\mathrm dB(f^{\sharp}(b)) = \sum_{i=1}^{g} f^{\sharp}(c_i(b)) \cdot \mathrm dB(f^{\sharp}(b_i))\), which vanishes by Step 1. Hence
  \[
    \mathrm dB(f^{\sharp}(b)) \;=\; 0 \qquad \text{in } \Omega_{R/k}, \quad \text{for every } b \in B,
  \]
  as a chart-local statement on \(V = \Spec R \subseteq C.\mathrm{left}\).

  \emph{Step 3.aux: 2-chart affine cover of \(C\) with affine intersection.} Step 3 below consumes a 2-chart affine cover \(\{V_1, V_2\}\) of \(C.\mathrm{left}\) with \(V_1 \cap V_2\) also affine. Existence is standard for a smooth proper curve over \(k\):
  \begin{itemize}
    \item \(C\) is quasi-compact and quasi-separated (proper \(\Rightarrow\) both; Mathlib's \texttt{AlgebraicGeometry.IsProper}-namespace instances combined with \texttt{IsSeparated.quasiSeparated}).
    \item Quasi-compactness plus the affine-cover API of \texttt{Mathlib.AlgebraicGeometry.AffineScheme} (the \texttt{IsAffineOpen} predicate and the \texttt{IsBasicOpen} / \texttt{LocalizationAway} family at \texttt{Mathlib.RingTheory.Localization.Away.Basic}) supplies a finite affine cover \(\{V_i\}_{i \in I}\) of \(C.\mathrm{left}\) (via \texttt{Scheme.affineCover} together with \texttt{Scheme.affineCover\_isFinite} for \(C\) Noetherian, available from quasi-compactness on a Noetherian base scheme).
    \item Refining to a two-chart cover with affine intersection follows the construction of Stacks Tag~0F8L: pick any non-empty affine \(V_1 \subseteq C.\mathrm{left}\); the complement \(C \setminus V_1\) is a closed subset of \(C.\mathrm{left}\) of codimension \(\geq 1\), i.e.\ (since \(C\) is one-dimensional) a finite set of closed points. A second affine open \(V_2 \subseteq C.\mathrm{left}\) containing \(C \setminus V_1\) exists by normality + irreducibility of \(C\) (a smooth proper curve is normal; for any finite set of closed points one can find an affine neighbourhood, e.g.\ via the \texttt{AffineCover} construction plus the \texttt{Scheme.basicOpen\_of\_isAffineOpen}-style basic-open API). The intersection \(V_1 \cap V_2 = V_1 \setminus (C \setminus V_2)\) is the complement in \(V_1\) of a finite closed subset of \(V_1\), hence a basic open of \(V_1\) — i.e.\ \(\Spec\) of a \texttt{LocalizationAway} construction on \(\Gamma(C, V_1)\), hence affine.
  \end{itemize}
  If the project does not already expose a \texttt{Scheme.affineCover\_card\_two}-shape lemma specialised to smooth proper curves, the construction lands in tree as a thin wrapper around the \texttt{AffineCover} + \texttt{LocalizationAway} APIs; estimated build \(\sim 20\)--\(40\) LOC.

  \emph{Step 3: chart-\v Cech promotion to global section-vanishing on \(C\).} The chart-local conclusion of Step 2 holds on the single chart \(V \subseteq C.\mathrm{left}\) originally chosen in the lemma statement. To upgrade ``\(\mathrm dB(f^{\sharp}(b))\) vanishes as a section of \(\Omega_{C/k}^{\oplus g}\) for every \(b \in B\)'' to global section-vanishing on \(C\) (the form consumed by Step 4 below; the \(\oplus g\) index records that each \(b \in B\) produces a \(g\)-tuple of section data via the standard-smooth basis \(\{\mathrm db_1, \ldots, \mathrm db_g\}\) of \(\Omega_{B/k}\)), we run a 2-chart \v Cech computation on the affine cover \(\{V_1, V_2\}\) of Step 3.aux.

  Concretely, the chart-\v Cech \(\mathrm d^{0}\)-complex of \(\Omega_{C/k}^{\oplus g}\) on \(\{V_1, V_2\}\) is the 2-term complex
  \[
    \mathrm d^{0} \colon\; \Gamma(V_1, \Omega_{C/k}^{\oplus g}) \,\oplus\, \Gamma(V_2, \Omega_{C/k}^{\oplus g}) \;\longrightarrow\; \Gamma(V_1 \cap V_2, \Omega_{C/k}^{\oplus g}), \qquad (\sigma_1, \sigma_2) \;\mapsto\; \sigma_1|_{V_1 \cap V_2} - \sigma_2|_{V_1 \cap V_2}.
  \]
  The chart-algebra (\(\alpha\)) helper \cref{lem:chart_algebra_isPushout_of_affine_product}, instantiated on the chart presentation of \(V_1 \cap V_2\) as an affine product of \(V_1\) and \(V_2\) over the common overlap (in chart-algebra terms: the structure-sheaf rings sit in an \texttt{Algebra.IsPushout} square \(k \to \Gamma(C, V_1)\) / \(\Gamma(C, V_2) \to \Gamma(C, V_1 \cap V_2)\); the chart-algebra (\(\alpha\)) helper supplies this square at the ring level, which globalises to the standard sheaf-gluing pattern for an affine cover with affine intersection), identifies the kernel of \(\mathrm d^{0}\) with the global sections via the standard sheaf-gluing pattern (a pair of chart sections agreeing on the intersection patches uniquely to a global section). Equivalently, the chart-side \(\mathrm d^{0}\) identification with global-section-restriction is the \(H^{0}\)-row of the abstract Mayer--Vietoris sequence \cref{thm:Scheme_AffineCoverMVSquare_HModule_prime_sequence_exact} instantiated on the 2-chart cover with sheaf \(\Omega_{C/k}^{\oplus g}\):
  \[
    0 \;\to\; \Gamma(C, \Omega_{C/k}^{\oplus g}) \;\to\; \Gamma(V_1, \Omega_{C/k}^{\oplus g}) \,\oplus\, \Gamma(V_2, \Omega_{C/k}^{\oplus g}) \;\xrightarrow{\mathrm d^{0}}\; \Gamma(V_1 \cap V_2, \Omega_{C/k}^{\oplus g}),
  \]
  so \(\ker \mathrm d^{0} = \Gamma(C, \Omega_{C/k}^{\oplus g})\).

  To verify the kernel element is zero, apply Steps 1--2 separately on each chart \(V_i\) — taking \(V := V_i\) in Step 1's chart-pair \((W, V)\) (chart-uniform existence is automatic: \(A.\mathrm{hom}\) is smooth at the image of \(\eta_A\), and the standard-smooth presentation on \(W\) propagates to all preimages of \(W\) in \(C\) with the same generators \(b_1, \ldots, b_g\); the chart-side ring map \(f^{\sharp}_i \colon B \to \Gamma(C, V_i)\) for each \(i\) then satisfies the chart-local kernel condition of Steps 1--2). Each chart-restriction \(\mathrm dB(f^{\sharp}(b))|_{V_i} \in \Gamma(V_i, \Omega_{C/k}^{\oplus g})\) vanishes by Step 2's chart-local conclusion. The pair \((\mathrm dB(f^{\sharp}(b))|_{V_1}, \mathrm dB(f^{\sharp}(b))|_{V_2})\) lies in \(\ker \mathrm d^{0}\) (automatic agreement on \(V_1 \cap V_2\) by sheaf restriction) and equals the zero pair. Hence \(\mathrm dB(f^{\sharp}(b)) = 0\) as a global section of \(\Omega_{C/k}^{\oplus g}\).

  \paragraph{Cohomological content; no named Serre duality.} The chart-\v Cech kernel-equals-global-sections identification supplied by the chart-algebra (\(\alpha\)) helper is the project's chart-algebraic substitute for the cohomological vanishing \(H^0(C, \Omega_{C/k}^{\oplus g}) = 0\) on a genus-\(0\) curve \(C/k\). Mathematically this content coincides with the Serre dual of \(H^1(C, \mathcal O_C^{\oplus g}) = 0\); per the iter-145 strategy-critic Q3 honesty note in the statement block above, the proof here does \emph{not} invoke Serre duality as a packaged named theorem, but the underlying \(H^0(C, \Omega_{C/k}) = 0\) \emph{is} invoked (in the equivalent chart-\v Cech kernel form). The chart-algebra (\(\alpha\)) gluing pattern + chart-local vanishing of Steps 1--2 is the running argument; the cohomological content is folded directly into the chart-\v Cech kernel computation rather than extracted as a separately-named theorem. (An earlier iter-145 draft cited \texttt{Genus.lean}'s \(H^{1}(C, \mathcal O_C) = 0\) computation as ``the running model''; that citation was inaccurate — \texttt{Genus.lean} is 45 LOC and only defines the genus, with no in-tree \(H^{1}\) vanishing computation — and is iter-146-corrected to point at the chart-algebra (\(\alpha\)) helper and the abstract MV theorem \cref{thm:Scheme_AffineCoverMVSquare_HModule_prime_sequence_exact} instead.)

  \emph{Step 4: ring-side kernel extraction via \cref{lem:KaehlerDifferential_mem_range_algebraMap_of_D_eq_zero}.} The ring-level helper packages the implication ``\(D b = 0 \Rightarrow b \in \operatorname{range}(\mathtt{algebraMap}\,k\,B)\)'' for any commutative \(k\)-algebra \(B\) where \(k\) is a field and \(D \colon B \to \Omega_{B/k}\) is the universal K\"ahler derivation. In characteristic \(0\) this is the direct kernel characterisation; in characteristic \(p > 0\) it incorporates a uniform Frobenius-iteration step using \texttt{RingHom.iterateFrobenius\_comm} (\texttt{Mathlib.Algebra.CharP.Frobenius}). Applied to \(R\) as a \(k\)-algebra, \(b := f^{\sharp}(b)\), and \(D := \mathrm dB\), using Step 3's vanishing \(\mathrm dB(f^{\sharp}(b)) = 0\), the helper produces
  \[
    f^{\sharp}(b) \;\in\; \operatorname{range}(\mathtt{algebraMap}\,k\,R)
  \]
  modulo the closure of the image of \(k\) inside the global sections of \(C\) under integral dependence over \(k\) — the residual question handled by Step 5.

  \emph{Step 5: integrally-closed-constants closure via \cref{lem:constants_integral_over_base_field}.} For \(C\) smooth proper geometrically irreducible over \(k\), the global sections \(\Gamma(C, \mathcal O_C)\) are integral over \(k\) and (by the smooth-proper-geometrically-irreducible fact) coincide with \(\operatorname{range}(\mathtt{algebraMap}\,k\,\Gamma(C, \mathcal O_C))\); the integrally-closed-constants helper packages this fact. The chart-restriction map \(\Gamma(C, \mathcal O_C) \hookrightarrow R\) (canonical sheaf restriction from \(C\) to \(V\), injective by sheaf separatedness on a reduced source) ensures that the image of \(k\) inside \(R\) also exhausts the integral closure of \(k\) in \(R\). Combining with Step 4 yields the conclusion \(f^{\sharp}(b) \in \operatorname{range}(\mathtt{algebraMap}\,k\,R)\).

  Estimated build: \(\sim 150\)--\(300\) LOC. Steps 1--2 (\(\sim 30\)--\(50\) LOC) consume the standard-smooth presentation of \(W\) via \texttt{Algebra.IsStandardSmooth.free\_kaehlerDifferential}. Step 3 (\(\sim 60\)--\(120\) LOC) is the dominant cost: instantiating the affine-cover Mayer--Vietoris square (\cref{thm:Scheme_AffineCoverMVSquare_HModule_prime_sequence_exact}) on \(\Omega_{C/k}^{\oplus g}\) and chasing the kernel. Steps 4--5 (\(\sim 60\)--\(130\) LOC) consume the two project-internal helpers below.

  \emph{Mathlib status.} The chart-side standard-smooth presentation: \texttt{Algebra.IsStandardSmooth.free\_kaehlerDifferential} and \texttt{Algebra.IsStandardSmoothOfRelativeDimension.rank\_kaehlerDifferential} are [verified] in \texttt{Mathlib.RingTheory.Smooth.StandardSmoothCotangent}. The K\"ahler-derivation functoriality \texttt{KaehlerDifferential.map} / \texttt{relativeDifferentials'\_map\_d} is [verified] in \texttt{Mathlib.RingTheory.Kaehler.Basic} / \texttt{Mathlib.Algebra.Category.ModuleCat.Differentials.Presheaf}. The chart-\v Cech kernel-equals-global-sections identification of Step 3 is supplied by the chart-algebra (\(\alpha\)) helper \cref{lem:chart_algebra_isPushout_of_affine_product} plus the abstract affine-cover Mayer--Vietoris sequence \cref{thm:Scheme_AffineCoverMVSquare_HModule_prime_sequence_exact} (\cref{chap:Cohomology_MayerVietoris}) instantiated on the 2-chart cover of Step 3.aux; the 2-chart cover construction itself consumes the \texttt{AffineCover} + \texttt{IsAffineOpen} + \texttt{LocalizationAway} API of \texttt{Mathlib.AlgebraicGeometry.AffineScheme} + \texttt{Mathlib.RingTheory.Localization.Away.Basic} (cf.\ Stacks Tag~0F8L for the construction). The two project-internal helpers \cref{lem:KaehlerDifferential_mem_range_algebraMap_of_D_eq_zero} and \cref{lem:constants_integral_over_base_field} discharge Steps 4 and 5 respectively.
\end{proof}

\begin{lemma}

\leanok
[(S3.sep.1) Smooth proper schemes have geometrically reduced global sections]
  \label{lem:S3_sep_1_smooth_geometrically_reduced_Gamma}
  \lean{AlgebraicGeometry.isGeometricallyReduced_Gamma_of_smooth}
  \uses{lem:S3_pi_1_Gamma_baseChange_iso_tensor_of_proper}
  % NOTE (iter-152, alg-closed pivot): DESCOPED under the alg-closed
  % pivot — general-over-k Mathlib-PR fodder, NOT on the M2.a critical
  % path. \cref{lem:constants_integral_over_base_field} now assumes
  % [IsAlgClosed k] and closes in three steps without this (S3.sep.1)
  % factorisation; its \uses no longer references this lemma. The label
  % and statement are retained for the general-over-k route.
  Let \(X\) be a smooth proper scheme over a field \(k\). Then the global-section \(k\)-algebra \(\Gamma(X, \mathcal O_X)\) is geometrically reduced over \(k\), i.e.\ \texttt{Algebra.IsGeometricallyReduced}\,\(k\)\,\(\Gamma(X, \mathcal O_X)\).

  This sub-claim is the (S3.sep.1) gateway substep of path (b) of the proof of \cref{lem:constants_integral_over_base_field}: a smart-proof reduction of \(\Gamma(X, \mathcal O_X) \cong k\) to the conjunction of separability and pure inseparability, where separability is delivered via geometric reducedness \(\Rightarrow\) separable on a finite field extension (see the off-critical-path remark below).
\end{lemma}

\begin{proof}
  
  \leanok
  % NOTE iter-151 render-fix: corrected Stacks Tag from `0334` (Nagata
  % rings, irrelevant) to `035U` (Section 33.6 "Geometrically reduced
  % schemes"); verbatim text now quoted from the bundled
  % references/stacks-varieties.tex (the `stacks-0334.md` typo-audit
  % file cited by the prior note never existed on disk).
  Smoothness of \(X / \Spec k\) implies that for any field extension \(K / k\) the base change \(X_K\) is reduced (smooth \(\Rightarrow\) formally smooth \(\Rightarrow\) geometrically reduced fibres; Stacks Tags~\textbf{035U} + \textbf{04QM}). The global-section ring \(\Gamma(X_K, \mathcal O_{X_K})\) injects into the stalks of \(X_K\) at its points (sheaf separatedness on a reduced source), each of which is reduced; hence \(\Gamma(X_K, \mathcal O_{X_K})\) is itself reduced as a \(K\)-algebra. By the flat-base-change identification \cref{lem:S3_pi_1_Gamma_baseChange_iso_tensor_of_proper} (whose content is the same as path (a) step~(e) of the proof of \cref{lem:constants_integral_over_base_field}, just re-packaged as a single named lemma), \(\Gamma(X, \mathcal O_X) \otimes_k K \cong \Gamma(X_K, \mathcal O_{X_K})\) is reduced for every field extension \(K / k\). That is, \(\Gamma(X, \mathcal O_X)\) is geometrically reduced over \(k\). Estimated build: \(\sim 80\)--\(150\) LOC.

  % NOTE iter-151 render-fix: cited tag changed from `0334` (Nagata,
  % irrelevant) to `035U` (the `stacks-0334.md` file never existed);
  % verbatim text quoted from the bundled references/stacks-varieties.tex.
  % SOURCE: [Stacks Project], Tag 035U (Definition 33.6.1, "geometrically
  % reduced") (read from references/stacks-varieties.tex, L330).
  % SOURCE QUOTE: "Let \(k\) be a field. Let \(X\) be a scheme over \(k\).
  % \begin{enumerate} \item Let \(x \in X\) be a point. We say \(X\) is
  % {\it geometrically reduced at \(x\)} if for any field extension \(k'/k\)
  % and any point \(x' \in X_{k'}\) lying over \(x\) the local ring
  % \(\mathcal{O}_{X_{k'}, x'}\) is reduced. \item We say \(X\) is
  % {\it geometrically reduced} over \(k\) if \(X\) is geometrically reduced
  % at every point of \(X\). \end{enumerate}"
  % SOURCE: [Stacks Project], Tag 04QM (Lemma, smooth ⇒ geometrically
  % regular/normal/reduced) (read from references/stacks-varieties.tex,
  % L4616, lemma-smooth-geometrically-normal).
  % SOURCE QUOTE: "Let \(X \to \Spec(k)\) be a smooth morphism where \(k\)
  % is a field. Then \(X\) is geometrically regular, geometrically normal,
  % and geometrically reduced over \(k\)."
  \textit{Source: Stacks Tags 035U (Definition 33.6.1) and 04QM.}
  \emph{Literature.} Stacks Tag~\textbf{035U} (geometrically reduced schemes, definition + properties); Stacks Tag~\textbf{04QM} (smooth morphisms have geometrically reduced fibres); Hartshorne III.10 (smoothness and reducedness).
\end{proof}

\begin{remark}
[(S3.sep.2) Geometrically reduced finite field extension is separable]
  This remark documents the off-critical-path fact (S3.sep.2): a finite field extension that is geometrically reduced over the base field is separable (Stacks Tag~\textbf{0BUG} part (4)). It is retained here as documentation for a possible future general-over-\(k\) route or upstream Mathlib PR, not as an active proof obligation.

  Let \(F / k\) be a finite field extension of fields. If \(F\) is geometrically reduced as a \(k\)-algebra (\texttt{Algebra.IsGeometricallyReduced}\,\(k\)\,\(F\), equivalently \(F \otimes_k \bar k\) is reduced), then \(F / k\) is separable (\texttt{Algebra.IsSeparable}\,\(k\)\,\(F\)).

  This sub-claim is the (S3.sep.2) closer of path (b)'s separability strand: the (S3.sep.1) gateway delivers geometric reducedness of \(\Gamma(X, \mathcal O_X)\) over \(k\); combined with finite-field-extension structure (from steps (a)--(d) of path (a) applied at the geometric level), this lemma upgrades geometric reducedness to separability of the finite field extension \(\Gamma(X, \mathcal O_X) / k\).

  By definition \texttt{Algebra.IsGeometricallyReduced}\,\(k\)\,\(F\) says \(F \otimes_k \bar k\) is reduced. Since \(F / k\) is finite, \(F \otimes_k \bar k\) is a finite-dimensional \(\bar k\)-algebra, hence Artinian, and decomposes uniquely as a finite product \(F \otimes_k \bar k \cong \prod_i F_i\) of Artinian local \(\bar k\)-algebras at its (finitely many) minimal primes. Reducedness of the product forces each Artinian local factor \(F_i\) to be a field, finite over \(\bar k\); algebraic closedness of \(\bar k\) then forces \(F_i = \bar k\).

  Therefore \(F \otimes_k \bar k \cong \bar k^{r}\) for \(r = \#\{\text{minimal primes of } F \otimes_k \bar k\}\). The \(\bar k\)-dimension on both sides gives \([F : k] = \dim_{\bar k}(F \otimes_k \bar k) = r\). On the other hand, \(r\) equals the number of distinct \(k\)-algebra embeddings \(F \hookrightarrow \bar k\) (the closed points of \(\Spec(F \otimes_k \bar k)\)). The equality \([F : k] = \#\{\text{embeddings } F \hookrightarrow \bar k\}\) is precisely the criterion for separability of a finite field extension. Estimated build: \(\sim 30\)--\(50\) LOC.

  % NOTE iter-151 render-fix: corrected Stacks Tag from `0BJF`
  % (discriminant ⇔ étale, irrelevant) to `0BUG` (part (4), the actual
  % "geom-reduced finite ⇒ separable" packaging); verbatim text quoted
  % from the bundled references/stacks-varieties.tex (the `stacks-0BJF.md`
  % and `stacks-0BUG.md` files cited by the prior note never existed).
  % SOURCE: [Stacks Project], Tag 0BUG (Lemma, global sections of a
  % proper scheme over a field), parts (3)-(4) (read from
  % references/stacks-varieties.tex, L1932,
  % lemma-proper-geometrically-reduced-global-sections).
  % SOURCE QUOTE: "Let \(k\) be a field. Let \(X\) be a proper scheme over
  % \(k\). \begin{enumerate} \item \(A = H^0(X, \mathcal{O}_X)\) is a finite
  % dimensional \(k\)-algebra, ... \item if \(X\) is reduced, then
  % \(A = \prod_{i = 1, \ldots, n} k_i\) is a product of fields, each a
  % finite extension of \(k\), \item if \(X\) is geometrically reduced, then
  % \(k_i\) is finite separable over \(k\), ..."
  \textit{Source: Stacks Tag 0BUG, part (4).}
  \emph{Literature.} Stacks Tag~\textbf{0BUG} part (4) (geometrically reduced + finite over \(k\) \(\Rightarrow\) separable, the load-bearing \(\Gamma\)-of-proper packaging); Bourbaki \emph{Alg\`ebre} Chap.~V \S 15 N\(^\circ\) 3 (separability criteria via tensor products with \(\bar k\)).
\end{remark}

\begin{lemma}

\leanok
[(S3.pi.1) Flat base change of \(\Gamma\) for proper schemes]
  \label{lem:S3_pi_1_Gamma_baseChange_iso_tensor_of_proper}
  \lean{AlgebraicGeometry.Gamma_baseChange_iso_tensor_of_proper}
  % NOTE (iter-152, alg-closed pivot): DESCOPED under the alg-closed
  % pivot — general-over-k Mathlib-PR fodder, NOT on the M2.a critical
  % path. This flat-base-change-of-Γ content was the substantive gap of
  % the OLD general-over-k constants proof (step (e)); under
  % [IsAlgClosed k] there is no base change to k̄, so
  % \cref{lem:constants_integral_over_base_field} no longer uses it.
  % Label + statement retained for the general-over-k route (this is the
  % genuine upstream-Mathlib-PR candidate). Supersedes the iter-150
  % HYBRID-DEFERRED disposition note below, which is now resolved YES.
  Let \(X\) be a proper scheme over a field \(k\), and let \(K / k\) be a field extension. Then the canonical comparison map
  \[
    \Gamma(X, \mathcal O_X) \otimes_k K \;\xrightarrow{\;\sim\;}\; \Gamma(X_K, \mathcal O_{X_K})
  \]
  is an isomorphism of \(K\)-algebras, where \(X_K := X \times_{\Spec k} \Spec K\) is the base change of \(X\) along \(\Spec K \to \Spec k\). This is the specialisation of ``cohomology along proper morphisms commutes with flat base change'' (Stacks Tag~\textbf{02KH}) to the \(H^0 = \Gamma\) row.

  This sub-claim is the (S3.pi.1) load-bearing substep of path (b) of the proof of \cref{lem:constants_integral_over_base_field}. Its content is the same content as path (a) step~(e), packaged as a single named lemma request rather than as a 7-step ad-hoc chain scattered across \(\Gamma_{\bar k}\)-cohomology infrastructure.

  % NOTE (iter-150 review): iter-150 mathlib-analogist (cross-domain
  % inspiration mode) surfaced a HYBRID part (A) consumer reformulation
  % over `\bar k` via `IsAlgClosed.algebraMap_bijective_of_isIntegral`
  % which, if accepted by the user gate on adding `[IsAlgClosed kbar]`
  % to `rigidity_over_kbar`'s signature, descopes BOTH (S3.pi.1) and
  % (S3.pi.2) from the M2.a critical path entirely. The Lean body at
  % `ChartAlgebraS3.lean:L220-242` carries an inline `HYBRID-DEFERRED`
  % block disclosing this disposition; the chapter prose above
  % (referring to (S3.pi.1) as "load-bearing substep of path (b)") was
  % accurate iter-149 and is now contingent on the user gate's NO
  % outcome. If the gate is YES, this lemma becomes deferred
  % indefinitely as upstream-Mathlib-PR work. See
  % `analogies/h1cotangent-vanishing-iter150.md` § "Top suggestion" and
  % `iter/iter-150/review.md` for the iter-150 lane outcome.
\end{lemma}

\begin{proof}

\leanok
  Cover \(X\) by finitely many affine opens \(U_i = \Spec A_i\) (quasi-compactness from properness). The \v Cech complex on the cover \(\{U_i\}\) computes \(\Gamma(X, \mathcal O_X)\) as the equalizer
  \[
    \Gamma(X, \mathcal O_X) \;=\; \mathrm{eq}\bigl(\textstyle\prod_i A_i \rightrightarrows \prod_{i, j} A_{ij}\bigr) \qquad \text{where } A_{ij} := \Gamma(U_i \cap U_j, \mathcal O_X).
  \]
  Tensoring with \(K\) over \(k\) preserves equalizers (flatness of \(K\) as a \(k\)-module). Chart-by-chart, the affine-side flat-base-change isomorphism \(A_i \otimes_k K \cong \Gamma((U_i)_K, \mathcal O_{(U_i)_K})\) holds by Stacks Tag~\textbf{00DS} (flat base change for affines), and similarly for the intersections \(A_{ij} \otimes_k K \cong \Gamma((U_i \cap U_j)_K, \mathcal O)\) (each \(U_i \cap U_j\) is itself quasi-affine and admits a finite affine cover whose flat-base-change identification reassembles to the intersection's flat-base-change isomorphism). Reassembling chart-by-chart with the same \v Cech equalizer description on the base-changed cover \(\{(U_i)_K\}\) identifies \(\Gamma(X, \mathcal O_X) \otimes_k K\) with \(\Gamma(X_K, \mathcal O_{X_K})\). Estimated build: \(\sim 150\)--\(250\) LOC (the dominant cost is the \v Cech equaliser + flat-tensor exchange chase, plus the verification that the affine intersection flat-base-change isomorphism reassembles to the global comparison).

  \emph{Honesty note.} The content of this lemma is the SAME content as path~(a) step~(e) of the 7-step chain in the proof of \cref{lem:constants_integral_over_base_field}, not bypassed by the smart-proof path~(b). Path~(b)'s genuine advantage is that it re-packages the flat-base-change content as ONE named Mathlib lemma request, rather than as a 7-step ad-hoc chain scattered across \(\Gamma_{\bar k}\)-cohomology infrastructure.

  % SOURCE: [Stacks Project], Tag 02KH (Lemma, "Flat base change"),
  % part (2) (read from references/stacks-coherent.tex, L947,
  % lemma-flat-base-change-cohomology).
  % SOURCE QUOTE: "Consider a cartesian diagram of schemes ... Let
  % \(\mathcal{F}\) be a quasi-coherent \(\mathcal{O}_X\)-module with pullback
  % \(\mathcal{F}' = (g')^*\mathcal{F}\). Assume that \(g\) is flat and that
  % \(f\) is quasi-compact and quasi-separated. For any \(i \geq 0\)
  % \begin{enumerate} \item the base change map ... is an isomorphism
  % \(\) g^*R^if_*\mathcal{F} \longrightarrow R^if'_*\mathcal{F}', \(\)
  % \item if \(S = \Spec(A)\) and \(S' = \Spec(B)\), then
  % \(H^i(X, \mathcal{F}) \otimes_A B = H^i(X', \mathcal{F}')\).
  % \end{enumerate}"
  \textit{Source: Stacks Tag 02KH, part (2) (\(H^0\) row).}
  \emph{Literature.} Stacks Tag~\textbf{02KH} (cohomology along proper morphisms commutes with flat base change); Hartshorne III.9.3 (flat base change of cohomology, the \(H^0\) row).
\end{proof}

\begin{remark}
[(S3.pi.2) Finite-dim \(k\)-algebra with unique-minimal-prime base change to \(\bar k\) is purely inseparable]
  This remark documents the off-critical-path fact (S3.pi.2): a finite extension of fields whose base change to \(\bar k\) has a unique minimal prime is purely inseparable (Stacks Tag~\textbf{030K}). It is retained here as documentation for a possible future general-over-\(k\) route or upstream Mathlib PR, not as an active proof obligation.

  Let \(F / k\) be a finite extension of fields, viewing \(F\) as a finite-dimensional \(k\)-algebra. Suppose that \(F \otimes_k \bar k\) has a unique minimal prime (equivalently, \(\Spec(F \otimes_k \bar k)\) is irreducible). Then \(F / k\) is purely inseparable, i.e.\ \texttt{Algebra.IsPurelyInseparable}\,\(k\)\,\(F\).

  This sub-claim is the (S3.pi.2) closer of path (b)'s pure-inseparability strand: the (S3.pi.1) flat-base-change identification supplies the unique-minimal-prime hypothesis (when applied to \(X\) smooth proper geometrically irreducible over \(k\), the geometric irreducibility of \(X_{\bar k}\) pins \(\Gamma(X_{\bar k}, \mathcal O_{X_{\bar k}}) \cong \Gamma(X, \mathcal O_X) \otimes_k \bar k\) to a local \(\bar k\)-algebra with a unique minimal prime); this lemma then upgrades the unique-minimal-prime hypothesis to pure inseparability of the finite field extension.

  \(F \otimes_k \bar k\) is a finite-dimensional \(\bar k\)-algebra, hence Artinian. Its nilradical reduction \((F \otimes_k \bar k) / \mathrm{Nil}\) has the same prime spectrum and, with the unique-minimal-prime hypothesis, is an Artinian integral domain finite over \(\bar k\) --- hence a field finite over \(\bar k\). Since \(\bar k\) is algebraically closed, that field equals \(\bar k\) itself. Therefore \((F \otimes_k \bar k) / \mathrm{Nil} \cong \bar k\), so \(F \otimes_k \bar k\) is a \emph{local} Artinian \(\bar k\)-algebra with residue field \(\bar k\).

  % NOTE iter-151 render-fix: corrected Stacks Tag from `05DH`
  % (universally injective, irrelevant) to `030K` (the
  % separable-then-purely-inseparable factorisation); verbatim text
  % quoted from the bundled references/stacks-fields.tex (the
  % `stacks-05DH.md` file cited by the prior note never existed).
  Equivalently: \(\Spec(F \otimes_k \bar k)\) is a single closed point. The standard criterion (Stacks Tag~\textbf{030K}: a finite field extension \(F / k\) is purely inseparable iff \(F \otimes_k \bar k\) is a local ring) immediately yields pure inseparability of \(F / k\). The dependence on \cref{lem:S3_pi_1_Gamma_baseChange_iso_tensor_of_proper} is at the consumer level: the unique-minimal-prime hypothesis on \(F \otimes_k \bar k\) is supplied in the proof of \cref{lem:constants_integral_over_base_field} via \(\Gamma(X, \mathcal O_X) \otimes_k \bar k \cong \Gamma(X_{\bar k}, \mathcal O_{X_{\bar k}})\) for \(X\) smooth proper geometrically irreducible (geometric irreducibility makes \(X_{\bar k}\) irreducible, so the right-hand side has a unique minimal prime). Estimated build: \(\sim 50\)--\(100\) LOC.

  % NOTE iter-151 render-fix: tag `05DH` (universally injective)
  % replaced with `030K`. CAVEAT: the bundled references/stacks-fields.tex
  % entry for Tag 030K is the separable-then-purely-inseparable
  % FACTORISATION lemma (lemma-separable-first), NOT a "finite extension
  % is purely inseparable iff F⊗_k k̄ is local" statement as the prior
  % render-fix note asserted. The verbatim quote below is the genuine
  % 030K text; the prose body's "single closed point" criterion is the
  % proof's own elementary argument, not a direct 030K citation. Flagged
  % to the plan agent (see report).
  % SOURCE: [Stacks Project], Tag 09HD (Definition, "purely inseparable")
  % (read from references/stacks-fields.tex, L1580,
  % definition-purely-inseparable).
  % SOURCE QUOTE: "Let \(F\) be a field of characteristic \(p > 0\). Let
  % \(K/F\) be an extension. \begin{enumerate} \item An element
  % \(\alpha \in K\) is {\it purely inseparable} over \(F\) if there exists a
  % power \(q\) of \(p\) such that \(\alpha^q \in F\). \item The extension \(K/F\)
  % is said to be {\it purely inseparable} if and only if every element
  % of \(K\) is purely inseparable over \(F\). \end{enumerate}"
  % SOURCE: [Stacks Project], Tag 030K (Lemma, separable-then-inseparable
  % factorisation) (read from references/stacks-fields.tex, L1703,
  % lemma-separable-first).
  % SOURCE QUOTE: "Let \(E/F\) be an algebraic field extension. There
  % exists a unique subextension \(E/E_{sep}/F\) such that \(E_{sep}/F\) is
  % separable and \(E/E_{sep}\) is purely inseparable."
  \textit{Source: Stacks Tags 09HD (Definition) and 030K.}
  \emph{Literature.} Stacks Tag~\textbf{09HD} (Section 9.14 ``Purely inseparable extensions'', the definition); Stacks Tag~\textbf{030K} (every algebraic extension factors uniquely as a separable extension followed by a purely inseparable one); Bourbaki \emph{Alg\`ebre} Chap.~V \S 7 N\(^\circ\) 7 (purely inseparable extensions via tensor products).
\end{remark}

\begin{lemma}\leanok
[(ii-\(\beta\)-aux) Constants of a smooth proper geometrically irreducible scheme are the base field]
  \label{lem:constants_integral_over_base_field}
  \lean{AlgebraicGeometry.constants_integral_over_base_field}
  \uses{lem:S3_pi_1_Gamma_baseChange_iso_tensor_of_proper}
  % NOTE (iter-146 review): the iter-146 Lean signature carries an
  % explicit `[IsReduced X]` typeclass hypothesis (in addition to the
  % `Smooth`, `IsProper`, `GeometricallyIrreducible` hypotheses listed
  % in the prose). Mathlib snapshot `b80f227` lacks a general
  % `Smooth (X → Spec k) ⇒ IsReduced X` lemma over a field (the closest
  % is `Scheme.Hom.dense_smoothLocus_of_perfectField`, perfect-base
  % only), so the project explicitly carries `[IsReduced X]` at every
  % call site. Once the Mathlib gap closes, the explicit-`IsReduced`
  % hypothesis will be derivable from `[Smooth (X → Spec k)]` and may
  % be dropped. Same discipline already documented for `Rigidity.lean`.
  % NOTE (iter-146 review): the iter-146 Lean closure formalises substeps
  % (1)–(2) of the three-substep recipe below (`IrreducibleSpace` +
  % `isIntegral_of_irreducibleSpace_of_isReduced` for (1);
  % `isField_of_universallyClosed` + `finite_appTop_of_universallyClosed`
  % for (2)). Substep (3) — the geometric-irreducibility base-change
  % step that pins `dim_k Γ(X, ⊤) = 1` — remains a structured `sorry`
  % at iter-146 close, scheduled for iter-147+ closure via the
  % `IsBaseChange`-namespace + Stacks Tag 02KH idiom plus
  % `Module.finrank_baseChange`.
  % NOTE (iter-147 review): the iter-147 prover lane reduces the
  % `RingHom.range (appTop.hom) = ⊤` goal to surjectivity of the
  % appTop ring hom via `RingHom.range_eq_top`, and documents the
  % full 7-step closure chain (a)–(g) as in-source comments inside
  % the body. The chain pulls back `X` along the algebraic closure
  % `\bar k`, applies `isField_of_universallyClosed` +
  % `finite_appTop_of_universallyClosed` on `X_{\bar k}`, then
  % `IsAlgClosed.algebraMap_bijective_of_isIntegral` to get
  % `Γ(X_{\bar k}, ⊤) ≅ \bar k`, then flat base change of `Γ` to
  % identify `Γ(X, ⊤) ⊗_k \bar k ≃ \bar k`, then
  % `Module.finrank_baseChange` + `Subalgebra.bot_eq_top_iff_finrank_eq_one`
  % to conclude. The substantive Mathlib gap — flat base change of
  % `Γ` for a proper scheme over a field (step (e)) — is concentrated
  % at the single residual `sorry`; the iter-148+ continuation lands a
  % thin in-tree wrapper around `AlgebraicGeometry.IsBaseChange` + the
  % `H^0` cohomology-base-change instance for the proper case.
  Let \(X\) be a smooth proper geometrically irreducible scheme over a field \(k\), with structure morphism \(\pi_X \colon X \to \Spec k\). Then the global sections form a finite-dimensional \(k\)-algebra that is integral over \(k\) (via the canonical \(\mathtt{algebraMap}\,k\,\Gamma(X, \mathcal O_X)\) from \(\pi_X\)); and the image of \(\mathtt{algebraMap}\,k\,\Gamma(X, \mathcal O_X)\) coincides with \(\Gamma(X, \mathcal O_X)\):
  \[
    \Gamma(X, \mathcal O_X) \;=\; \operatorname{range}(\mathtt{algebraMap}\,k\,\Gamma(X, \mathcal O_X)).
  \]
  Equivalently, \(\Gamma(X, \mathcal O_X) \cong k\) as \(k\)-algebras.

  \emph{Iter-152 alg-closed pivot.} The conclusion (\(\mathtt{RingHom.range}\,(\mathtt{appTop.hom}) = \top\)) is unchanged, but under the alg-closed pivot the proof now \textbf{assumes the base field \(k\) is algebraically closed} (\texttt{[IsAlgClosed k]}): the lemma is invoked only with \(k = \bar k\) in the alg-closed rigidity setting (\cref{thm:rigidity_over_kbar}). This collapses the proof to the three short steps below and descopes the entire (S3.\textsc{sep}/\textsc{pi}) base-change-to-\(\bar k\) chain.
\end{lemma}

\begin{proof}

\leanok
  % NOTE (iter-152, alg-closed pivot): the proof now ASSUMES the base
  % field k is algebraically closed ([IsAlgClosed k]); the lemma is
  % invoked only with k = k̄ in the alg-closed rigidity setting
  % (rigidity_over_kbar). Under [IsAlgClosed k] the former path-(a)
  % 7-step flat-base-change chain (a)–(g) and the path-(b) 4×(S3.*)
  % factorisation both COLLAPSE to the three short steps below: there is
  % no base change to k̄ (we already are over k̄), so flat-base-change
  % of Γ (the old substantive gap, step (e)) and the (S3.sep/pi)
  % separability/inseparability factorisation are no longer needed. The
  %  are pruned from this block accordingly; the (S3.*)
  % lemmas remain DEFINED above (descoped, general-over-k Mathlib-PR
  % fodder — see their statement-block NOTEs).
  Since \(k\) is algebraically closed (\texttt{[IsAlgClosed k]}), the proof is three steps:
  \begin{enumerate}
    \item[(1)] \emph{\(X\) is integral, so \(\Gamma(X, \mathcal O_X)\) is a field.} Smoothness gives \texttt{IsReduced}\,\(X\) and the \texttt{GeometricallyIrreducible} hypothesis gives \texttt{IrreducibleSpace}\,\(X\) (over the singleton base \(\Spec k\), via \texttt{GeometricallyIrreducible.irreducibleSpace\_of\_subsingleton}); together (\texttt{isIntegral\_of\_irreducibleSpace\_of\_isReduced}) \(X\) is integral. Properness gives universally closed, and \texttt{isField\_of\_universallyClosed} then shows \(\Gamma(X, \mathcal O_X)\) is a field.
    \item[(2)] \emph{\(\Gamma(X, \mathcal O_X)\) is a finite field extension of \(k\).} Properness (\(\Rightarrow\) universally closed) makes the structure-morphism appTop ring hom \(k \to \Gamma(X, \mathcal O_X)\) finite (\texttt{finite\_appTop\_of\_universallyClosed}); in particular \(\Gamma(X, \mathcal O_X)\) is integral over \(k\). Combined with (1), \(\Gamma(X, \mathcal O_X)\) is a finite (hence integral) field extension of \(k\).
    \item[(3)] \emph{Algebraic closure absorbs the finite field extension.} Since \(k\) is algebraically closed, an integral extension is trivial: apply \texttt{IsAlgClosed.algebraMap\_bijective\_of\_isIntegral} (\texttt{Mathlib.FieldTheory.IsAlgClosed.Basic}, with \(\{k\,K\}\) \texttt{[Field k] [Ring K] [IsDomain K] [IsAlgClosed k] [Algebra k K] [Algebra.IsIntegral k K]} concluding \texttt{Function.Bijective (algebraMap k K)}) to \(K := \Gamma(X, \mathcal O_X)\). Bijectivity of \(\mathtt{algebraMap}\,k\,\Gamma(X, \mathcal O_X)\) gives in particular surjectivity, i.e.\ \(\mathtt{RingHom.range}\,(\mathtt{appTop.hom}) = \top\), which is the displayed conclusion (equivalently \(\Gamma(X, \mathcal O_X) \cong k\)).
  \end{enumerate}
  Estimated build: \(\sim 30\)--\(60\) LOC; the dominant work is the integral-structure assembly of step (1). No flat-base-change machinery is required under the alg-closed hypothesis.

  \emph{Mathlib status.} \texttt{isField\_of\_universallyClosed} and \texttt{finite\_appTop\_of\_universallyClosed} (steps (1)--(2)) are [verified] in \texttt{Mathlib.AlgebraicGeometry} via the \texttt{IsProper} / \texttt{UniversallyClosed}-namespace lemmas. \texttt{GeometricallyIrreducible.irreducibleSpace\_of\_subsingleton} (step (1)) is [verified] in \texttt{Mathlib.AlgebraicGeometry.GeometricallyIrreducible}, and \texttt{isIntegral\_of\_irreducibleSpace\_of\_isReduced} in \texttt{Mathlib.AlgebraicGeometry.Properties}. \texttt{IsAlgClosed.algebraMap\_bijective\_of\_isIntegral} (step (3)) is [verified] in \texttt{Mathlib.FieldTheory.IsAlgClosed.Basic}. \texttt{RingHom.range\_eq\_top} (converting the displayed range goal to surjectivity) is [verified] in \texttt{Mathlib.RingTheory.RingHom}.

  \emph{Literature.} Stacks Tag~\textbf{0BUG} (``for \(X\) proper geometrically integral over \(k\), \(\Gamma(X, \mathcal O_X)\) is a finite separable field extension of \(k\)'' --- the global structural statement; over an algebraically closed \(k\) the finite extension is trivial); Hartshorne III.5.2 (\(\Gamma\) of proper integral schemes over a field). The general-over-\(k\) route (the path-(a) flat-base-change chain and the path-(b) (S3.\textsc{sep}/\textsc{pi}) factorisation) is preserved for reference in the now-descoped (S3.*) lemma blocks above.
\end{proof}

\begin{lemma}
[(ii-\(\beta\)-aux-chart) Image of the base field is integrally closed in a chart of a smooth proper geometrically irreducible scheme --- \textbf{off-path record (iter-155)}]
  \label{lem:KaehlerDifferential_constants_in_chart_of_proper_curve}
  \lean{AlgebraicGeometry.TODO.KaehlerDifferential_constants_in_chart_of_proper_curve}
  \uses{lem:constants_integral_over_base_field}
  % NOTE (iter-155): DEMOTED to an off-path historical record. The live
  % (FT.1)–(FT.3) route that closed KDM
  % (`lem:KaehlerDifferential_mem_range_algebraMap_of_D_eq_zero`) at
  % iter-154 consumes NONE of this lemma — the chart-to-function-field
  % integral-closure argument here was part of the superseded char-p /
  % (p3) exploration. The KDM statement-block \uses was pruned of this
  % label at iter-155 accordingly. The block is retained for the record;
  % it is not on the critical path for any live result.
  \paragraph{Iter-155 disposition.} This lemma is \textbf{off the critical path}: the live KDM closure (\cref{lem:KaehlerDifferential_mem_range_algebraMap_of_D_eq_zero}, route (FT.1)--(FT.3)) consumes neither it nor \cref{lem:chart_algebra_isPushout_of_affine_product}, and the stale KDM statement-block \texttt{\textbackslash uses} pointing here has been pruned. The block is preserved as an auditable record of the superseded chart-to-function-field integral-closure exploration only.

  Let \(X\) be a smooth proper geometrically irreducible scheme over a field \(k\), let \(V \subseteq X.\mathrm{left}\) be a non-empty affine open subset, and let \(B := \Gamma(X, V)\) denote the chart-side ring (a \(k\)-algebra via the structure map \(X \to \Spec k\), restricted to \(V\)). Then the image of \(k\) inside \(B\) via the canonical algebra map is integrally closed in \(B\):
  \[
    \operatorname{range}(\mathtt{algebraMap}\,k\,B) \;\;\text{is integrally closed in } B.
  \]
  Equivalently, any \(b \in B\) satisfying a monic polynomial relation over \(\operatorname{range}(\mathtt{algebraMap}\,k\,B)\) already lies in \(\operatorname{range}(\mathtt{algebraMap}\,k\,B)\).

  \paragraph{Iter-146 scope note.} This lemma extracts the chart-of-smooth-proper-geometrically-irreducible-scheme hypothesis out of the proof of \cref{lem:KaehlerDifferential_mem_range_algebraMap_of_D_eq_zero} (where step (p3) of that proof previously appealed to it implicitly), concentrating the curve-specific (more precisely, smooth-proper-geom-irr-specific) input at a named lemma the iter-146 KDM ring-side prover lane can dispatch as a black box. The KDM lemma's signature is left at its current finite-type-\(k\)-algebra generality; consumers (notably \cref{lem:chart_algebra_df_zero_factors_through_constant_on_chart}) supply the chart-of-proper-curve hypothesis at their site by invoking this lemma where the descent step requires it.
\end{lemma}

\begin{proof}
  
  By \cref{lem:constants_integral_over_base_field} the global-section ring \(\Gamma(X, \mathcal O_X)\) coincides with the image of the canonical \(\mathtt{algebraMap}\,k\,\Gamma(X, \mathcal O_X)\); in particular \(\Gamma(X, \mathcal O_X)\) is a copy of \(k\) inside \(X\)'s global structure sheaf.

  The chart restriction \(\rho_V \colon \Gamma(X, \mathcal O_X) \to B = \Gamma(X, V)\) is the canonical sheaf-restriction map. On a reduced source (smoothness of \(X\) over \(k\) implies reducedness, Mathlib's \texttt{AlgebraicGeometry.SmoothScheme}-namespace lemmas combined with the standard reducedness-from-regular-local-rings instance), the structure-sheaf restriction along an open immersion to a non-empty affine open is injective (sheaf separatedness on a reduced source: \texttt{AlgebraicGeometry.IsOpenImmersion} family + \texttt{Scheme.Spec\_germ\_injective}-shape lemmas). Hence \(\rho_V\) embeds the field \(\Gamma(X, \mathcal O_X) = k\) inside \(B\) as \(\operatorname{range}(\mathtt{algebraMap}\,k\,B)\).

  For the integral-closure property: suppose \(b \in B\) is integral over \(\operatorname{range}(\mathtt{algebraMap}\,k\,B)\), i.e.\ satisfies \(b^{n} + a_{n-1} b^{n-1} + \cdots + a_0 = 0\) in \(B\) with each \(a_i \in \operatorname{range}(\mathtt{algebraMap}\,k\,B)\). Smoothness plus geometric irreducibility of \(X\) over \(k\) gives \(X\) integral (smooth \(\Rightarrow\) reduced; geometrically irreducible \(\Rightarrow\) irreducible), so the affine chart \(V = \Spec B\) embeds into the function field \(K(X)\) of \(X\) via the canonical generic-point localisation (\texttt{Mathlib.AlgebraicGeometry.FunctionField}; an open immersion of integral schemes induces an injection on rings into \(K(X)\)). The same monic relation holds in \(K(X)\), with coefficients \(a_i \in k \subseteq K(X)\), so \(b\) is integral over \(k\) inside \(K(X)\). The integral closure of \(k\) in \(K(X)\) is the field of constants of \(X\), which equals \(\Gamma(X, \mathcal O_X)\) for \(X\) smooth proper geometrically irreducible over \(k\) (a standard fact, Stacks Tag~0BUG; the integral elements over \(k\) in \(K(X)\) correspond to regular functions on \(X\), which by properness are global sections, which by \cref{lem:constants_integral_over_base_field} are exactly \(k\)). Hence \(b \in k\) inside \(K(X)\), and pulling back along \(\rho_V\) exhibits \(b \in \operatorname{range}(\mathtt{algebraMap}\,k\,B)\).

  Estimated build: \(\sim 80\)--\(150\) LOC. The dominant cost is the chart-to-function-field embedding \(B \hookrightarrow K(X)\) (an open-immersion-induced localisation idiom; chase via \texttt{AlgebraicGeometry.Scheme.Hom.appLE} on the generic point) and the "integral closure of \(k\) in \(K(X)\) equals \(\Gamma(X, \mathcal O_X)\)" identification (a thin wrapper around the constants-vanishing of \cref{lem:constants_integral_over_base_field} plus the standard normality-of-smooth-schemes fact).

  \emph{Mathlib status.} The chart-of-reduced-source injectivity for \(\rho_V\) is [verified] via \texttt{AlgebraicGeometry.IsOpenImmersion}-namespace lemmas on the structure sheaf. The function-field embedding \(B \hookrightarrow K(X)\) for \(B\) an affine chart of an integral scheme is [verified in idiom] in \texttt{Mathlib.AlgebraicGeometry.FunctionField}. The "integral closure of \(k\) in \(K(X)\) equals constants" identification for a smooth proper geometrically irreducible scheme is the standard Stacks Tag~0BUG; in Mathlib snapshot \texttt{b80f227} it is [gap]-fill via a thin in-tree wrapper that chains \cref{lem:constants_integral_over_base_field} with the normality-of-smooth-\(k\)-schemes instance (\texttt{AlgebraicGeometry.Smooth.normal}).
\end{proof}

\begin{lemma}[(FT.3 base case) The universal derivation does not kill the indeterminate in the rational function field]
  \label{lem:ratfunc_D_X_ne_zero}
  \lean{AlgebraicGeometry._ratfunc_D_X_ne_zero}
  Let \(k\) be a field. In the K\"ahler module \(\Omega_{\operatorname{Frac}(k[X])/k}\) of the rational function field over \(k\), the image of the indeterminate under the universal derivation is nonzero:
  \[
    D_{\operatorname{Frac}(k[X])}\bigl(\mathtt{algebraMap}\,(k[X])\,(\operatorname{Frac} k[X])\,X\bigr) \;\neq\; 0 .
  \]
\end{lemma}

\begin{proof}
  Proved directly in Lean.
\end{proof}

\begin{lemma}[(FT.3 closer) Over an algebraically closed base, an algebraic element lies in the image of the base field]
  \label{lem:algebraic_mem_range}
  \lean{AlgebraicGeometry._algebraic_mem_range}
  Let \(k\) be an algebraically closed field and \(K\) a field that is a \(k\)-algebra. If \(b \in K\) is algebraic over \(k\) (\(\mathtt{IsAlgebraic}\,k\,b\)), then \(b \in \operatorname{range}(\mathtt{algebraMap}\,k\,K)\).
\end{lemma}

\begin{proof}
  Proved directly in Lean.
\end{proof}

\begin{theorem}\leanok
[(ii-\(\beta\)-core ring-side) Kernel of the universal K\"ahler derivation: \(D b = 0\) implies \(b\) in the image of the base field]
  \label{lem:KaehlerDifferential_mem_range_algebraMap_of_D_eq_zero}
  \lean{AlgebraicGeometry.KaehlerDifferential.mem_range_algebraMap_of_D_eq_zero}
  \uses{lem:ratfunc_D_X_ne_zero, lem:algebraic_mem_range}
  % iter-155: stale statement-block \uses pruned. The live (FT.1)–(FT.3)
  % route (proof block below) consumes NEITHER
  % `lem:chart_algebra_isPushout_of_affine_product` NOR
  % `lem:KaehlerDifferential_constants_in_chart_of_proper_curve` — it is a
  % self-contained assembly of shipped Mathlib lemmas (the proof-block
  % \uses was already pruned at iter-154). Removing the now-stale
  % statement-block \uses keeps the dependency graph honest;
  % `lem:KaehlerDifferential_constants_in_chart_of_proper_curve` is thereby
  % demoted to an off-path historical record (see its statement-block note).
  % NOTE (iter-152, alg-closed pivot): the iter-151 lane proved this lemma is
  % FALSE as a bare B-only algebra lemma. Two counterexamples satisfy the old
  % hypotheses ([Field k] [CharZero k] [FiniteType k B]
  % [IsStandardSmoothOfRelativeDimension n k B]) yet have ker D ⊋ range(algebraMap):
  %   (CE1) B = k×k, n=0 (finite étale, so D=0; (1,0) ∉ diagonal range);
  %   (CE2) k=ℚ, B=ℚ(√2), n=0 (finite separable, so D=0; range = ℚ ⊊ ℚ(√2)).
  % The CORRECTED statement (this block) adds the JOINT hypotheses
  % [IsAlgClosed k] + [IsDomain B] (the char-zero hypothesis is retained). These
  % are BOTH needed and each excludes one counterexample:
  %   - [IsDomain B] kills CE1 (k×k is not a domain);
  %   - [IsAlgClosed k] kills CE2 (over an algebraically closed base, ℚ(√2)-type
  %     finite nontrivial separable extensions of the base inside Frac(B) cannot
  %     occur — k is algebraically closed in Frac(B)).
  % The iter-151 NOTE that "[IsDomain B] alone fails" only considered the
  % hypotheses one-at-a-time; the JOINT addition is sound. The missing geometric
  % content ("k algebraically closed in B") that was lost when the scheme-level
  % GeometricallyIrreducible+IsReduced data was reduced to a B-only algebra lemma
  % is exactly what [IsAlgClosed k] + [IsDomain B] restore. Do NOT re-weaken this
  % statement without re-introducing both counterexamples.
  Let \(k\) be an algebraically closed field of characteristic zero (\texttt{[Field k] [IsAlgClosed k] [CharZero k]}), and let \(B\) be a finite-type standard-smooth \(k\)-algebra of relative dimension \(n\) that is moreover an integral domain (\texttt{[Algebra.FiniteType k B]}, \texttt{[Algebra.IsStandardSmoothOfRelativeDimension n k B]}, \texttt{[IsDomain B]}). Let \(D \colon B \to \Omega_{B/k}\) denote the universal K\"ahler derivation (Mathlib's \texttt{KaehlerDifferential.D} as a \(B\)-linear additive Leibniz-rule map). Then for every \(b \in B\),
  \[
    D b \;=\; 0 \quad\Longrightarrow\quad b \;\in\; \operatorname{range}(\mathtt{algebraMap}\,k\,B).
  \]
\end{theorem}

\begin{proof}

\leanok
  % iter-154: \uses pruned to reflect the live route. The single-element /
  % H1Cotangent / perfect-field FT route below is a self-contained assembly of
  % shipped Mathlib lemmas and depends on NO blueprint label; it supersedes the
  % char-p (p1) and Differential.ContainConstants (p2) explorations that
  % consumed lem:chart_algebra_isPushout_of_affine_product and
  % lem:KaehlerDifferential_constants_in_chart_of_proper_curve (those \cref's
  % survive only inside the demoted historical records below and create no live
  % dependency edge).
  \textbf{Live route (iter-154, verified Mathlib assembly).} The reverse inclusion ``\(b \in \operatorname{range}(\mathtt{algebraMap}\,k\,B) \Rightarrow D b = 0\)'' is the standard \texttt{Derivation.map\_algebraMap} (\texttt{Mathlib.RingTheory.Derivation.Basic}) and holds with no extra hypotheses. The forward direction ``\(D b = 0 \Rightarrow b \in \operatorname{range}(\mathtt{algebraMap}\,k\,B)\)'' is established under the corrected hypotheses (\texttt{[IsAlgClosed k] [CharZero k]} on \(k\), \texttt{[IsDomain B]} on \(B\)) by the three-step \emph{single-element / perfect-field / Jacobi--Zariski \(H^1\)-cotangent} route (FT.1)--(FT.3) below. This route is a \emph{self-contained assembly of shipped Mathlib lemmas} (snapshot \texttt{b80f227}): the iter-154 \texttt{mathlib-analogist} consult verified that every step composes by compilation against the project toolchain (eight type-checking \texttt{example} blocks in \texttt{analogies/ftthree-kernel-iter154.md}), \emph{overturning} the iter-149--153 ``genuine Mathlib gap / bright-line STOP'' verdict on (FT.3). It needs neither a separating transcendence basis nor the polynomial-ring free-case helpers (C.a)--(C.c), and it makes the (p1)/(p2) explorations unnecessary: all of (C.a)--(C.c), (p1), (p2) are retained \textbf{below the (FT) list only as an auditable historical record} and are \textbf{not on the critical path}. (Only \texttt{[IsDomain B]} \(+\) \texttt{[FiniteType k B]} are actually consumed by the kernel argument --- to build \(K = \operatorname{Frac} B\) and \(\mathtt{EssFiniteType}\,k\,K\); the standard-smooth/relative-dimension-\(n\) hypotheses are retained in the signature but are harmless slack for this proof.)

  \emph{The three-step route (FT.1)--(FT.3).} Write \(D_R\) for the universal K\"ahler derivation of a \(k\)-algebra \(R\).
  \begin{itemize}
    \item[(FT.1)] \emph{Push to the fraction field.} \(B\) is an integral domain, so the localisation map \(B \hookrightarrow K := \operatorname{Frac}(B)\) is injective (\texttt{IsFractionRing}, \texttt{Mathlib.RingTheory.Localization.FractionRing}). Functoriality of the universal derivation (\texttt{KaehlerDifferential.map\_D}; the project's \texttt{\_hFunct} reduction) carries the chart hypothesis \(D_B b = 0\) to \(D_K b = 0\) in \(\Omega_{K/k}\). Since \(B \hookrightarrow K\) is injective, it suffices to prove that \(b\), viewed in \(K\), lies in \(\operatorname{range}(\mathtt{algebraMap}\,k\,K)\); pulling back along the injection then yields \(b \in \operatorname{range}(\mathtt{algebraMap}\,k\,B)\).
    \item[(FT.2)] \emph{Reduce to ``\(b\) algebraic over \(k\)'' (the perfect-field / \(H^1\)-cotangent core).} Argue by contradiction: suppose \(b\) is transcendental over \(k\). Then the \(k\)-algebra map \(k[X] \to K\), \(X \mapsto b\), is injective, so by the universal property of localisation (\texttt{IsFractionRing.lift}) the rational function field \(F := \operatorname{RatFunc} k \cong \operatorname{Frac}(k[X])\) embeds as a subfield of \(K\), giving a tower \(k \to F \to K\) with \(\mathtt{IsScalarTower}\,k\,F\,K\). Because \(\operatorname{char} k = 0\), \(F\) is a \emph{perfect} field (\texttt{PerfectField.ofCharZero}, \texttt{Mathlib.FieldTheory.Perfect}); and \(K\) is essentially of finite type over \(F\), since it is essentially of finite type over \(k\) --- \(\mathtt{EssFiniteType}\,k\,K\) holds because \(K = \operatorname{Frac} B\) is a localisation of the finite-type \(k\)-algebra \(B\) (\texttt{Algebra.EssFiniteType.comp}\,\(k\,B\,K\)) --- and one descends along the tower (\texttt{Algebra.EssFiniteType.of\_comp}, \texttt{Mathlib.RingTheory.EssentialFiniteType}). A field extension of a perfect field that is essentially of finite type is formally smooth: \(\mathtt{Algebra.FormallySmooth}\,F\,K\) (\texttt{Algebra.FormallySmooth.of\_perfectField}, \texttt{Mathlib.RingTheory.Smooth.Field}). Formal smoothness makes the first cotangent homology \(H^1(L_{K/F})\) a subsingleton (\texttt{instSubsingletonH1CotangentOfFormallySmooth}, \texttt{Mathlib.RingTheory.Smooth.Basic}). The Jacobi--Zariski exact sequence of the tower \(k \to F \to K\),
    \[
      H^1(L_{K/F}) \xrightarrow{\ \delta\ } K \otimes_F \Omega_{F/k} \xrightarrow{\ \mathtt{mapBaseChange}\,k\,F\,K\ } \Omega_{K/k},
    \]
    is exact at the middle term (\texttt{Algebra.H1Cotangent.exact\_\(\delta\)\_mapBaseChange}, \texttt{Mathlib.RingTheory.Kaehler.CotangentComplex}); with \(H^1(L_{K/F})\) a subsingleton the image of \(\delta\) is \(\{0\}\), so \(\mathtt{mapBaseChange}\,k\,F\,K\) is \emph{injective}. The base-change/derivation identity \(\mathtt{mapBaseChange}\,k\,F\,K\,(1 \otimes D_F x) = D_K(\mathtt{algebraMap}\,F\,K\,x)\) (from \texttt{KaehlerDifferential.mapBaseChange\_tmul}, \texttt{one\_smul} and \texttt{KaehlerDifferential.map\_D}, \texttt{Mathlib.RingTheory.Kaehler.Basic}), evaluated at the element \(b \in F\) (the image of the indeterminate \(X\) under \(\mathtt{algebraMap}\,(k[X])\,F\)), gives \(\mathtt{mapBaseChange}(1 \otimes D_F b) = D_K b = 0\); injectivity forces \(1 \otimes D_F b = 0\) in \(K \otimes_F \Omega_{F/k}\). Since the field extension \(F \to K\) is faithfully flat (\texttt{Module.FaithfullyFlat}, automatic for a field extension), the map \(\omega \mapsto 1 \otimes \omega\) on \(\Omega_{F/k}\) is injective (\texttt{Module.FaithfullyFlat.one\_tmul\_eq\_zero\_iff}, \texttt{Mathlib.RingTheory.Flat.FaithfullyFlat.Basic}), whence \(D_F b = 0\). This contradicts the base case (FT.3). Hence \(b\) is algebraic over \(k\).
    \item[(FT.3)] \emph{Transcendental base case, then the algebraic-closure closer.}
    \begin{itemize}
      \item \emph{Base case} (the only genuinely new helper, \({\sim}20\) LOC, verified): \(D_{\operatorname{Frac}(k[X])} X \neq 0\) --- the universal derivation does not kill the indeterminate in the rational function field. A localisation is formally \'etale (\texttt{Algebra.FormallyEtale.of\_isLocalization}, \texttt{Mathlib.RingTheory.Etale.Basic}), so \(\Omega_{\operatorname{Frac}(k[X])/k}\) is the localisation of \(\Omega_{k[X]/k}\) at the non-zero-divisors of \(k[X]\): the comparison map \(\mathtt{map}\,k\,k\,(k[X])\,(\operatorname{Frac} k[X])\) is a localised-module map (\texttt{KaehlerDifferential.isLocalizedModule\_map}, \texttt{Mathlib.RingTheory.Kaehler.Localization}). If \(D_{\operatorname{Frac}(k[X])} X = 0\), this map sends \(D_{k[X]} X\) to \(0\), so by \texttt{IsLocalizedModule.eq\_zero\_iff} (\texttt{Mathlib.Algebra.Module.LocalizedModule.Basic}) there is a non-zero-divisor \(c \in k[X]\) with \(c \cdot D_{k[X]} X = 0\). The \(k[X]\)-linear isomorphism \texttt{KaehlerDifferential.polynomialEquiv} (\(\Omega_{k[X]/k} \cong k[X]\)) carries \(D_{k[X]} X\) to the formal derivative \(X' = 1\) (\texttt{polynomialEquiv\_D} with \texttt{Polynomial.derivative\_X}, \texttt{Mathlib.RingTheory.Kaehler.Polynomial}); applying it gives \(c \cdot 1 = 0\) in \(k[X]\), contradicting \texttt{nonZeroDivisors.coe\_ne\_zero}.
      \item \emph{Closer} (\({\sim}10\) LOC, verified): from ``\(b\) algebraic over \(k\)'' conclude \(b \in \operatorname{range}(\mathtt{algebraMap}\,k\,K)\). The simple extension \(k(b)\) is finite-dimensional over \(k\) (\texttt{IntermediateField.adjoin.finiteDimensional}, applied to the integrality witnessed by \(\mathtt{IsAlgebraic}\,k\,b\)), hence integral over \(k\) (\texttt{Algebra.IsIntegral.of\_finite}); since \(k\) is algebraically closed, \(\mathtt{algebraMap}\,k\,k(b)\) is bijective (\texttt{IsAlgClosed.algebraMap\_bijective\_of\_isIntegral}, \texttt{Mathlib.FieldTheory.IsAlgClosed.Basic}). The generator of \(k(b)\) --- and therefore \(b\) itself inside \(K\) --- thus lies in \(\operatorname{range}(\mathtt{algebraMap}\,k\,K)\). This is the \emph{same} Mathlib lemma the project already invokes in the proof of \cref{lem:constants_integral_over_base_field}.
    \end{itemize}
    Combining the two: (FT.2) makes \(b\) algebraic over \(k\), the closer puts \(b \in \operatorname{range}(\mathtt{algebraMap}\,k\,K)\), and (FT.1) pulls this back to \(b \in \operatorname{range}(\mathtt{algebraMap}\,k\,B)\), closing the forward direction.
  \end{itemize}
  % NOTE (iter-154, live route — replaces the iter-152 "residual content" note):
  % The iter-154 mathlib-analogist consult OVERTURNED the iter-149–153 "genuine
  % Mathlib gap / bright-line STOP" verdict on (FT.3). (a) The single-element /
  % H1Cotangent / perfect-field route (FT.1)–(FT.3) above was verified to compile
  % end-to-end against snapshot b80f227 (eight type-checking `example` blocks in
  % analogies/ftthree-kernel-iter154.md, which is the source of truth for this
  % route). (b) The separating-transcendence-basis route is DISCARDED: Mathlib
  % ships no IsTranscendenceBasis-keyed freeness-of-Ω lemma (only the
  % SubmersivePresentation / IsStandardSmoothOfRelativeDimension rank lemmas and
  % Algebra.trdeg), and the single-element route makes a transcendence basis
  % unnecessary. (c) The `_mvPoly_*` free-case helpers (C.a)–(C.c) are now DEAD
  % code superseded by this route and may be removed from the Lean tree; they
  % survive in prose immediately below ONLY as an auditable historical record.

  \medskip
  \noindent\textbf{--- Historical record (NOT on the critical path; superseded by the live (FT) route above). ---}
  \medskip

  \emph{Superseded polynomial-ring layer (C.a)--(C.c) --- DEAD code.} (Superseded by the live (FT.1)--(FT.3) route above, which never leaves the fraction field \(K\) and uses no polynomial-ring free-case helper; retained as an auditable record of the iter-150 deposit only. The \texttt{\_mvPoly\_*} \texttt{private} lemmas below are now dead code and may be removed from \texttt{AlgebraicJacobian/Cotangent/ChartAlgebra.lean}.) These three sub-steps were closed in the iter-150 Lean tree and transferred the kernel question from \(B\) to its free polynomial presentation \(P.\mathrm{Ring} = \mathrm{MvPolynomial}\,\iota\,k\), where the kernel-equals-constants fact is elementary, pushing the conclusion back along the standard-smooth presentation map \(\pi \colon P.\mathrm{Ring} \to B\):
  \begin{itemize}
    \item[(C.a)] \emph{FREE-CASE.} Over \(k\) of characteristic \(0\), if the universal K\"ahler derivation \(D\) kills \(f \in \mathrm{MvPolynomial}\,\sigma\,k\) (i.e.\ \(D f = 0\)), then \(f \in (\mathrm{MvPolynomial.C}).\mathrm{range}\) (\(f\) is a constant). Proof: \(D f = 0\) forces every partial \(\mathrm{pderiv}_i f = 0\) (via the explicit free basis \texttt{KaehlerDifferential.mvPolynomialBasis}, with \(\mathrm{repr}(D f)\,i = \mathrm{pderiv}_i f\)); the coefficient identity \(\mathrm{coeff}_{u - \mathrm{single}_i 1}(\mathrm{pderiv}_i f) = \mathrm{coeff}_u f \cdot (u\,i)\) for \(u\,i \geq 1\) then forces every non-constant coefficient to vanish in characteristic \(0\) (where the integer \(u\,i\) is invertible). \textbf{Closed in Lean} (now DEAD): project lemmas \texttt{\_mvPoly\_coeff\_pderiv\_at\_shifted}, \texttt{\_mvPoly\_mem\_range\_C\_of\_pderiv\_eq\_zero}, \texttt{\_mvPoly\_mem\_range\_C\_of\_D\_eq\_zero} (\texttt{AlgebraicJacobian/Cotangent/ChartAlgebra.lean}; \texttt{private}, not in Mathlib).
    \item[(C.b)] \emph{Standard-smooth presentation.} Extract a \texttt{SubmersivePresentation}\,\(P\) from the \texttt{Algebra.IsStandardSmoothOfRelativeDimension}\,\(n\)\,\(k\)\,\(B\) hypothesis (via \texttt{Algebra.IsStandardSmooth.out}), giving a polynomial ring \(P.\mathrm{Ring} = \mathrm{MvPolynomial}\,\iota\,k\) together with a surjective \(k\)-algebra map \(\pi = \mathtt{algebraMap}\,P.\mathrm{Ring}\,B\) (surjectivity by \texttt{Algebra.Generators.algebraMap\_surjective}). \textbf{Closed in Lean} (now DEAD).
    \item[(C.c)] \emph{Lift \(+\) functoriality.} Lift \(b \in B\) to \(\tilde b \in P.\mathrm{Ring}\) along the surjection \(\pi\); the functoriality identity \texttt{KaehlerDifferential.map\_D} turns the hypothesis \(D_B b = 0\) into \((\mathtt{map}\,k\,k\,P.\mathrm{Ring}\,B)(D_A \tilde b) = 0\), i.e.\ \(D_A \tilde b \in \ker(\mathtt{map}\,k\,k\,P.\mathrm{Ring}\,B)\). \textbf{Closed in Lean} (the \texttt{\_hFunct} step; the \texttt{map\_D} functoriality is reused by (FT.1) above, but the polynomial-ring lift is DEAD).
  \end{itemize}

  \emph{Superseded path (p2) --- characteristic-\(0\) via \texttt{Differential.ContainConstants}.} (Superseded by the live single-element / \(H^1\)-cotangent (FT.1)--(FT.3) route above; retained as an auditable record of the iter-145--149 exploration. The iter-150 Lean body does \emph{not} register the \texttt{Differential B} instance this route relies on, and the iter-154 verified route makes the entire (BR.1)--(BR.5) bridge unnecessary.) When \(\mathtt{CharZero}\,k\) (or more generally when the chart-algebra has trivial \(p\)-torsion in its K\"ahler differentials --- so that the integer coefficient \(\alpha_i\) in the \(D\)-expansion of (p1.b) below never vanishes), the kernel of \(D \colon B \to \Omega_{B/k}\) is exactly \(\operatorname{range}(\mathtt{algebraMap}\,k\,B)\). The Mathlib idiom is the \texttt{Differential.ContainConstants} typeclass (\texttt{Mathlib.RingTheory.Derivation.DifferentialRing}, lines 62--70; one Mathlib consumer at \texttt{FieldTheory.Differential.Liouville}): a differential ring satisfies \texttt{ContainConstants} when the kernel of the chosen derivation equals the image of its constant subring. The typeclass is positioned for a specific derivation \(B \to B\) (a \texttt{Differential B} instance), \emph{not} for the universal K\"ahler derivation \(D \colon B \to \Omega_{B/k}\), so the bridge from \(D\) to a coefficient-derivation family \(\{\partial_i \colon B \to B\}_{i = 1}^{n}\) is project-internal infrastructure (\emph{not} Mathlib infrastructure). The (BR.1)--(BR.5) decomposition below structures the bridge, mirroring the (p1.a)--(p1.f) structuring of the alternative char-\(p\) path:
  \begin{itemize}
    \item[(BR.1)] \emph{Signature inflation.} Inflate the KDM signature with the typeclass conjunction \([\mathtt{CharZero}\,k] + [\mathtt{Algebra.IsStandardSmoothOfRelativeDimension}\,k\,B\,n]\). This is a signature-level refactor only; no new lemma, no Mathlib work. Estimated build: signature edit only.
    \item[(BR.2)] \emph{Basis selection of \(\Omega_{B/k}\).} Apply \texttt{Algebra.IsStandardSmooth.free\_kaehlerDifferential} (\texttt{Mathlib.RingTheory.Smooth.StandardSmoothCotangent}; [verified] in Mathlib snapshot \texttt{b80f227}) to extract a \(B\)-basis \(\{\mathrm dx_1, \ldots, \mathrm dx_n\}\) of \(\Omega_{B/k}\) from the standard-smooth chart presentation, with the rank-\(n\) identification supplied by \texttt{Algebra.IsStandardSmoothOfRelativeDimension.rank\_kaehlerDifferential}. No project work; the Mathlib lemma supplies the basis directly. Estimated build: existing Mathlib lemma; \({\sim}0\) LOC.
    \item[(BR.3)] \emph{Coefficient-derivation extraction \(\partial_i \colon B \to B\).} For each \(i \in \{1, \ldots, n\}\), project the universal derivation \(D\) onto the \(\mathrm dx_i\)-coefficient under the basis of (BR.2): \(\partial_i \colon B \to B,\ b \mapsto\) the \(\mathrm dx_i\)-coefficient in the expansion \(D b = \sum_{i = 1}^{n} (\partial_i b)\, \mathrm dx_i\). This is a project-internal construction: no off-the-shelf Mathlib lemma converts a universal K\"ahler derivation into a family of basis-coefficient \(B\)-self-derivations. Suggested Lean target name \texttt{KaehlerDifferential.coordinateDeriv}. Estimated build: \(\sim 30\)--\(50\) LOC of project work.
    \item[(BR.4)] \emph{\texttt{Differential}\,\(B\) instance per \(\partial_i\).} Each \(\partial_i\) is \(k\)-linear by construction (the basis-projection of the \(k\)-linear universal derivation \(D\)), hence in particular \(\mathbb Z\)-linear, and inherits the Leibniz rule from \(D\). Package it as a \texttt{Derivation \(\mathbb Z\) B B} (Mathlib's \texttt{RingTheory.Derivation.Basic}) and register the \texttt{Differential}\,\(B\) instance carried by that derivation. Suggested Lean target name \texttt{KaehlerDifferential.coordinateDerivDifferentialInstance}. Estimated build: \(\sim 10\)--\(20\) LOC; assemblable from \texttt{RingTheory.Derivation.Basic} once (BR.3) lands.
    \item[(BR.5)] \emph{\texttt{Differential.ContainConstants} instance for the chosen \(\partial_i\) in \(\mathtt{CharZero}\,k + \mathtt{IsStandardSmooth}\,k\,B\).} Mathlib defines the \texttt{Differential.ContainConstants} class (\texttt{Mathlib.RingTheory.Derivation.DifferentialRing} L62--70) but does \emph{not} ship an instance for the case relevant here (standard-smooth \(k\)-algebra in characteristic \(0\), with the coordinate derivation \(\partial_i\)). The project supplies the missing instance by hand: in characteristic \(0\), the kernel of the \(i\)-th partial derivative on a polynomial ring \(k[x_1, \ldots, x_n]\) is the \(k\)-subalgebra generated by all variables other than \(x_i\); the standard-smooth presentation \(B \cong k[x_1, \ldots, x_n, y_1, \ldots, y_m] / (f_1, \ldots, f_m)\) propagates this via the Jacobian invertibility (the \(\mathrm dy_j\)'s are \(B\)-linear combinations of \(\mathrm dx_i\)'s), so the \(\mathtt{ContainConstants}\) kernel of \(\partial_i\) on \(B\) is the \(k\)-subalgebra generated by all \(x_j\) (\(j \neq i\)) and all \(y_j\). Intersecting the \(n\) coordinate \(\mathtt{ContainConstants}\) kernels over \(i = 1, \ldots, n\) collapses to \(\operatorname{range}(\mathtt{algebraMap}\,k\,B)\). Suggested Lean target name \texttt{KaehlerDifferential.coordinateDeriv\_containConstants\_of\_charZero}. Estimated build: \(\sim 40\)--\(80\) LOC; the in-tree workaround constructs the instance from \texttt{RingTheory.Derivation.Basic} plus the standard-smooth presentation, since Mathlib has only the class definition, not an instance for this case.
  \end{itemize}
  % NOTE (iter-152 writer): the polynomial-ring layer (C.a)-(C.c) remains
  % first-class prose above, but the old (C.d) "transfer step" has been
  % REPLACED by the corrected field-of-fractions argument (FT.1)-(FT.3)
  % under the alg-closed pivot (the bare B-only goal that (C.d) targeted
  % is false). The (BR.3) coordinate-derivation construction is documented
  % but NOT used; the (BR.4) `Differential B` instance is NOT registered.
  % This whole (p2) block is a SUPERSEDED auditable record only. Provenance:
  % `analogies/h1cotangent-vanishing-iter150.md` § "Top suggestion (C)";
  % `iter/iter-150/review.md` § "Lane 2 KDM (BR.5)";
  % `task_results/lean-vs-blueprint-checker-chartalgebra-iter150.md`
  % (iter-150 must-fix on the prose-vs-Lean divergence, addressed here).

  Composing (BR.1)--(BR.5) closes the (p2) forward direction: the standard-smooth chart hypothesis from (BR.1) makes \(\Omega_{B/k}\) free of finite rank \(n\) via (BR.2); the (BR.3) coefficient-derivation family \(\{\partial_i\}\) + (BR.4) \texttt{Differential}\,\(B\) instances per \(\partial_i\) bridge from the universal derivation \(D\) to a family of \(B\)-self-derivations covered by the \texttt{Differential.ContainConstants} interface; and the (BR.5) instance in characteristic \(0\) then gives \(D b = 0 \Rightarrow (\forall i)\,\partial_i b = 0 \Rightarrow b \in \operatorname{range}(\mathtt{algebraMap}\,k\,B)\) directly, with no Frobenius detour. The (p2) path does \emph{not} consume \cref{lem:KaehlerDifferential_constants_in_chart_of_proper_curve}: the \texttt{ContainConstants} typeclass closes the forward direction by itself, eliminating the chart-of-proper-curve dependency that the (p1) char-\(p\) alternative carries. Aggregate (p2) build: \(\sim 80\)--\(150\) LOC for the full (BR.1)--(BR.5) bridge body. The (BR.3)--(BR.5) sub-steps may be extracted to per-sub-step Lean lemmas locally by a prover that wants finer-grained statements; they are not promoted to first-class blueprint declarations here because they live inside the (p2) bridge proof and their statements depend on the (BR.1) signature inflation.

  % NOTE iter-151 render-fix: corrected Stacks Tag from `07F4`
  % (Lemma 16.8.3, smoothing ring maps, irrelevant) to `00T7`
  % (Lemma 10.137.6, the standard-smooth ⇒ Ω-free citation);
  % verbatim text in references/stacks-algebra.tex (L37258); the
  % `stacks-07F4.md` file cited by the prior note never existed.
  \emph{Superseded alternative path (p1) --- characteristic-\(p\) via standard-smooth chart presentation.} (Superseded by the live single-element / \(H^1\)-cotangent (FT.1)--(FT.3) route above, which commits to \texttt{[CharZero k]} in the Lean signature and runs \emph{no} Frobenius iteration; the char-\(p\) chain below is retained as an auditable record only and is off the critical path.) When \(\mathtt{CharP}\,k\,p\) for \(p > 0\), the kernel of \(D\) is strictly larger than \(\operatorname{range}(\mathtt{algebraMap}\,k\,B)\) in general: any \(b = c^p\) with \(c \in B\) has \(D(c^p) = p c^{p-1} D c = 0\), so \(\ker D \supseteq B^p\) (the image of the absolute Frobenius on \(B\)). The recovery follows the Cartier-direction characterisation \(\ker D = B^p\) on a standard-smooth chart (Stacks Tag~00T7), then a Frobenius-power iteration, then descent through the chart-of-proper-curve integrally-closed-constants helper. Since Mathlib snapshot \texttt{b80f227} does not ship the chart-level \(\ker D = B^p\) as an off-the-shelf lemma for the standard-smooth-chart case, the project's in-tree build decomposes the chart-level step into four sub-steps (p1.a)--(p1.d), followed by the Frobenius-iteration step (p1.e) and the chart-of-proper-curve descent (p1.f):
  \begin{itemize}
    \item[(p1.a)] \emph{Subring structure of \(\ker D\).} The kernel \(\ker D \subseteq B\) is closed under addition (\texttt{KaehlerDifferential.D\_add} of \texttt{Mathlib.RingTheory.Kaehler.Basic}: \(D(b_1 + b_2) = D b_1 + D b_2\)) and under multiplication via the Leibniz rule (\texttt{Derivation.leibniz} / \texttt{KaehlerDifferential.D\_mul}: \(D(b_1 b_2) = b_1 \cdot D b_2 + b_2 \cdot D b_1\); when both factors lie in \(\ker D\), the right-hand side vanishes). It contains \(\operatorname{range}(\mathtt{algebraMap}\,k\,B)\) (\texttt{Derivation.map\_algebraMap}). Hence \(\ker D\) is a \(k\)-sub-algebra of \(B\) --- a \texttt{KaehlerDifferential.D\_ker\_subring}-shape statement, assemblable from the two named Mathlib lemmas plus the standard \texttt{Subring}-builder boilerplate of \texttt{Mathlib.RingTheory.Subring.Basic}.
    \item[(p1.b)] \emph{Standard-smooth chart presentation and \(D\)-coefficient calculus.} The standard-smooth-of-relative-dimension-\(n\) chart hypothesis (\texttt{Algebra.IsStandardSmoothOfRelativeDimension} of \texttt{Mathlib.RingTheory.Smooth.StandardSmooth}, instantiated on the chart of \cref{lem:chart_algebra_isPushout_of_affine_product}) supplies a presentation \(B \cong k[x_1, \ldots, x_n, y_1, \ldots, y_m] / (f_1, \ldots, f_m)\) with the Jacobian \(\det(\partial f_i / \partial y_j)\) invertible in \(B\). The universal K\"ahler module \(\Omega_{B/k}\) is free of rank \(n\) on the basis \(\{\mathrm dx_1, \ldots, \mathrm dx_n\}\) (\texttt{Algebra.IsStandardSmooth.free\_kaehlerDifferential} of \texttt{Mathlib.RingTheory.Smooth.StandardSmoothCotangent}; the rank-\(n\) identification via \texttt{Algebra.IsStandardSmoothOfRelativeDimension.rank\_kaehlerDifferential}). For any \(b = \sum_{\alpha \in \mathbb N^n} c_\alpha\, x^\alpha\;+\;(\text{relations})\), the universal derivation expands as
    \[
      D b \;=\; \sum_{\alpha} c_\alpha\!\!\sum_{i=1}^{n} \alpha_i\, x_i^{\alpha_i - 1}\!\!\prod_{j \neq i} x_j^{\alpha_j}\, \mathrm dx_i \qquad \pmod{\mathrm df_1, \ldots, \mathrm df_m}.
    \]
    In characteristic \(p\), the integer coefficient \(\alpha_i\) vanishes precisely when \(p \mid \alpha_i\). Modulo the standard-smooth relations (the Jacobian condition expresses each \(\mathrm dy_j\) as a \(B\)-linear combination of the \(\mathrm dx_i\)'s, eliminating the \(\mathrm dy_j\)'s coherently), \(D b = 0\) unpacks via the freeness of the \(\{\mathrm dx_i\}\)-basis into the coefficient-wise system \(\alpha_i\, c_\alpha = 0\) for every \(\alpha\) and every \(i\).
    \item[(p1.c)] \emph{Conclusion: \(b\) is a polynomial in \(x_i^p\) modulo relations.} The system \(\alpha_i\, c_\alpha = 0\) in characteristic \(p\) forces \(c_\alpha = 0\) unless every \(\alpha_i\) is divisible by \(p\). Writing \(\alpha_i = p \beta_i\), the polynomial unpacks as
    \[
      b \;=\; \sum_{\beta \in \mathbb N^n} c_{p \beta}\, x^{p \beta} \;+\; (\text{relations}) \;=\; \sum_{\beta} c_{p \beta}\,(x^{\beta})^{p} \;+\; (\text{relations}).
    \]
    Freeness of \(\Omega_{B/k}\) on \(\{\mathrm dx_i\}\) from (p1.b) makes the coefficient-extraction step rigorous: separating the \(\mathrm dx_i\)-components of \(D b\) shows that \(D b = 0\) is genuinely the conjunction of the per-\(\mathrm dx_i\) component vanishings, not a coincidental cancellation across \(i\).
    \item[(p1.d)] \emph{Lift back to \(b \in B^p\).} The polynomial relation \(b = \sum_\beta c_{p \beta}\,(x^\beta)^{p}\,\pmod{\text{relations}}\) exhibits \(b\) as a \(p\)-th power up to the standard-smooth relations. Setting \(c := \sum_\beta c_{p \beta}^{(p)} \cdot x^\beta + (\text{relations})\), where \(c_{p \beta}^{(p)} \in B\) is a \(p\)-th root of \(c_{p \beta}\) supplied at the chart-coefficient level (the coefficients \(c_{p \beta} \in k \cdot (\text{constants part of } B)\) admit \(p\)-th roots via the Frobenius-iteration depth bound of (p1.e) below, propagating to the constants part of \(B\), where the absolute Frobenius is the identity on the \(p\)-th-power-closed constants; for \(k\) perfect this is direct, and for non-perfect \(k\) the (p1.e) depth absorbs the residual data), one verifies \(c^{p} = b\) modulo the relations via the characteristic-\(p\) Frobenius-power expansion. The expansion uses that the iterated Frobenius is a ring homomorphism commuting with \(\mathtt{algebraMap}\), packaged in Mathlib as \texttt{RingHom.iterateFrobenius\_comm} (\texttt{Mathlib.Algebra.CharP.Frobenius}): this single lemma supplies the additivity + multiplicativity of \(\mathrm{Frobenius}\) (and of every iterate \(\mathrm{Frobenius}^{e}\)) needed to compute \(c^p = \sum_\beta c_{p \beta}\,(x^\beta)^{p}\), matching \(b\) up to relations. Together with the free-module structure \(\Omega_{B/k} \cong B^{n}\) from \texttt{Algebra.IsStandardSmooth.free\_kaehlerDifferential} (which guarantees the coefficient-extraction is unambiguous), this delivers \(b \in B^p\). The cleanest in-tree concretisation pattern follows Stacks Tag~00T7 (the standard-smooth presentation supplies a free \(\Omega_{B/k}\) with basis \(\{\mathrm dx_i\}\); in characteristic \(p\) the Cartier-direction inclusion \(\ker D \supseteq B^p\) for a smooth \(k\)-algebra follows from this basis identification); the explicit construction of \(c\) is the contribution of this sub-step. Estimated build for the (p1.a)--(p1.d) chain: \(\sim 60\)--\(100\) LOC.
    \item[(p1.e)] \emph{Frobenius iteration via \texttt{RingHom.iterateFrobenius\_comm}.} Mathlib's \texttt{RingHom.iterateFrobenius\_comm} (\texttt{Mathlib.Algebra.CharP.Frobenius}) gives the commutation \(\mathrm{Frobenius}^n \circ f = f \circ \mathrm{Frobenius}^n\) for any ring map \(f \colon B \to B\) in characteristic \(p\). Applied iteratively on the chain \(b = c_1^p\) from (p1.d) (with \(D c_1\) being zero modulo \(B^p\) by re-running the (p1.a)--(p1.d) chain on the corresponding K\"ahler derivation of \(c_1\)), we obtain a chain \(b = c_1^p = c_2^{p^2} = \cdots = c_e^{p^e}\) of \(p^e\)-th roots, all lying in \(B\). The chain stabilises in finitely many steps because \(B\) is a finitely-generated \(k\)-algebra with a standard-smooth presentation of relative dimension \(n\) from piece (i); the Frobenius depth \(e\) is uniformly bounded by the chart's structural data. Estimated build: \(\sim 40\)--\(80\) LOC.
    \item[(p1.f)] \emph{Descent to \(\operatorname{range}(\mathtt{algebraMap}\,k\,B)\) via the chart-of-proper-curve integrally-closed-constants helper.} At chart-uniform Frobenius depth \(e\), \(b = c_e^{p^e}\) with \(c_e \in B\) and \(D c_e = 0\) modulo the iterated reduction. The chart-uniform descent forces \(c_e\) to be integral over \(\operatorname{range}(\mathtt{algebraMap}\,k\,B)\) inside \(B\) (via the iter-on-(p1.a)--(p1.d) inductive base: at each Frobenius layer the kernel-of-\(D\) data confines the lifted-\(c_e\) to algebraic dependence on the \(k\)-image inside the chart). By \cref{lem:KaehlerDifferential_constants_in_chart_of_proper_curve} --- invoked on \(V := \Spec B\) viewed as an affine chart of the smooth proper geometrically irreducible ambient scheme supplied by the consumer (see consumer block \cref{lem:chart_algebra_df_zero_factors_through_constant_on_chart} Step 4 for the chart-of-proper-curve hypothesis instantiation; the consumer carries the ambient-scheme data, which this proof step inherits as a standing chart-side premise) --- \(\operatorname{range}(\mathtt{algebraMap}\,k\,B)\) is integrally closed in \(B\). Hence \(c_e \in \operatorname{range}(\mathtt{algebraMap}\,k\,B)\), and Frobenius-stability of \(\operatorname{range}(\mathtt{algebraMap}\,k\,B)\) (the Frobenius of a field of characteristic \(p\) lands inside the image of the prime field, hence inside the image of \(k\)) gives \(b = c_e^{p^e} \in \operatorname{range}(\mathtt{algebraMap}\,k\,B)\). Estimated build: \(\sim 40\)--\(50\) LOC (the bulk concentrated in \cref{lem:KaehlerDifferential_constants_in_chart_of_proper_curve}).
  \end{itemize}
  Composing (p1.a)--(p1.f) closes the characteristic-\(p\) alternative path. Estimated build for (p1) as a whole: \(\sim 140\)--\(230\) LOC, dominated by the (p1.a)--(p1.d) chain (\(\sim 60\)--\(100\) LOC), the Frobenius iteration (p1.e) (\(\sim 40\)--\(80\) LOC), and the chart-of-proper-curve descent (p1.f) (\(\sim 40\)--\(50\) LOC). Mathlib's existing \texttt{Algebra.IsSeparable.iff\_isPurelyInseparable\_extension} delivers only the separable-direction \(\ker D = B^p\) (too restrictive); the (p1.a)--(p1.d) chain above is the standard-smooth-direction version.

  \emph{Closure end state and ordering.} Under the iter-152 alg-closed pivot the signature carries \texttt{[IsAlgClosed k] [CharZero k] [IsDomain B]} (in addition to the finite-type standard-smooth hypotheses), and the live proof is the single-element / perfect-field / Jacobi--Zariski \(H^1\)-cotangent chain (FT.1)--(FT.3) above, a \emph{self-contained assembly of shipped Mathlib lemmas} verified by the iter-154 \texttt{mathlib-analogist} consult to compile end-to-end against snapshot \texttt{b80f227}. It uses neither a separating transcendence basis nor the polynomial-ring helpers (C.a)--(C.c). All three itemisations above and below --- the now-DEAD (C.a)--(C.c) polynomial-ring layer, the (p2) \texttt{Differential.ContainConstants}-via-instance bridge, and the (p1) characteristic-\(p\) Cartier-direction chain --- are recorded for provenance only: (C.a)--(C.c) is superseded by the fraction-field route, while (p1)/(p2) targeted the \emph{bare} \(B\)-only statement (now known false) and carry none of the \texttt{[IsAlgClosed k]} \(+\) \texttt{[IsDomain B]} content; none of the three is on the critical path. The only genuinely new helper in the live route is the transcendental base case \(D_{\operatorname{Frac}(k[X])} X \neq 0\) (FT.3, \({\sim}20\) LOC, verified); the remaining \({\sim}80\)--\(120\) LOC are scalar-tower and instance plumbing around the cited Mathlib lemmas, with no Archon theory gap remaining.

  \emph{Signature-vs-proof reconciliation (iter-146 disposition --- historical, pertains to the demoted (p1)/(p2) records).} The KDM lemma's signature is at its finite-type-\(k\)-algebra (or standard-smooth-of-relative-dimension-\(n\)) generality, but the (p1) char-\(p\) alternative path's descent step (p1.f) relies on \cref{lem:KaehlerDifferential_constants_in_chart_of_proper_curve}, which carries a chart-of-smooth-proper-geometrically-irreducible-ambient-scheme hypothesis not visible in the KDM signature. Neither the (p2) char-\(0\) path nor --- crucially --- the \emph{live} single-element / \(H^1\)-cotangent (FT.1)--(FT.3) route consumes the chart-of-proper-curve helper, so the live route carries no such hidden dependency at all (it is a self-contained Mathlib assembly). This is by design (iter-146 disposition (B.preferred)): the chart-of-proper-curve hypothesis is concentrated at the named helper \cref{lem:KaehlerDifferential_constants_in_chart_of_proper_curve}, which the consumers of KDM (currently a single consumer, \cref{lem:chart_algebra_df_zero_factors_through_constant_on_chart}) invoke in tandem with KDM only when the (p1) char-\(p\) branch fires. A future formalization may either inflate KDM's signature to carry the chart-of-proper-curve hypothesis explicitly, or route the integrally-closed-constants step at the consumer site (so KDM's body proves only the chart-uniform Frobenius-iteration reduction up to integral-closure data); both routes are mathematically equivalent. The blueprint blocks here document the (B.preferred) split.

  \emph{Mathlib status (live route).} Every lemma in the live (FT.1)--(FT.3) route is [verified by compilation, snapshot \texttt{b80f227}] by the iter-154 analogist consult (citation table and eight type-checking \texttt{example} blocks in \texttt{analogies/ftthree-kernel-iter154.md}): \texttt{IsFractionRing.lift} (\texttt{Mathlib.RingTheory.Localization.FractionRing}); \texttt{PerfectField.ofCharZero} (\texttt{Mathlib.FieldTheory.Perfect}); \texttt{Algebra.EssFiniteType.comp}, \texttt{.of\_comp} (\texttt{Mathlib.RingTheory.EssentialFiniteType}); \texttt{Algebra.FormallySmooth.of\_perfectField} (\texttt{Mathlib.RingTheory.Smooth.Field}); \texttt{instSubsingletonH1CotangentOfFormallySmooth} (\texttt{Mathlib.RingTheory.Smooth.Basic}); \texttt{Algebra.H1Cotangent.exact\_\(\delta\)\_mapBaseChange} (\texttt{Mathlib.RingTheory.Kaehler.CotangentComplex}); \texttt{KaehlerDifferential.mapBaseChange\_tmul}, \texttt{map\_D} (\texttt{Mathlib.RingTheory.Kaehler.Basic}); \texttt{Module.FaithfullyFlat.one\_tmul\_eq\_zero\_iff} (\texttt{Mathlib.RingTheory.Flat.FaithfullyFlat.Basic}); \texttt{Algebra.FormallyEtale.of\_isLocalization} (\texttt{Mathlib.RingTheory.Etale.Basic}); \texttt{KaehlerDifferential.isLocalizedModule\_map} (\texttt{Mathlib.RingTheory.Kaehler.Localization}); \texttt{IsLocalizedModule.eq\_zero\_iff} (\texttt{Mathlib.Algebra.Module.LocalizedModule.Basic}); \texttt{KaehlerDifferential.polynomialEquiv}, \texttt{polynomialEquiv\_D} (\texttt{Mathlib.RingTheory.Kaehler.Polynomial}); \texttt{IntermediateField.adjoin.finiteDimensional}, \texttt{Algebra.IsIntegral.of\_finite}, \texttt{IsAlgClosed.algebraMap\_bijective\_of\_isIntegral} (\texttt{Mathlib.FieldTheory.IsAlgClosed.Basic}). The route is therefore a pure Mathlib assembly with no Archon theory gap. The remaining citations in this paragraph pertain to the demoted historical (p1)/(p2) records only.

  \emph{Mathlib status (historical (p1)/(p2) records).} \texttt{KaehlerDifferential.D} and the kernel infrastructure are [verified] in \texttt{Mathlib.RingTheory.Kaehler.Basic}. \texttt{Derivation.map\_algebraMap} is [verified] in \texttt{Mathlib.RingTheory.Derivation.Basic}. \texttt{Differential.ContainConstants} (the primary (p2) closure mechanism) is [verified] in \texttt{Mathlib.RingTheory.Derivation.DifferentialRing} with one Mathlib consumer at \texttt{FieldTheory.Differential.Liouville}; the universal-K\"ahler-to-\texttt{Differential B} bridge of (p2) (basis projection of \(\Omega_{B/k}\) to coefficient derivations \(\partial_i\)) is the substantive (p2) build, leveraging \texttt{Algebra.IsStandardSmooth.free\_kaehlerDifferential} ([verified], \texttt{Mathlib.RingTheory.Smooth.StandardSmoothCotangent}) to fix a basis. \texttt{RingHom.iterateFrobenius\_comm} (the iterated Frobenius commutation that supplies Frobenius's additivity + multiplicativity in (p1.d) and the depth-bounded iteration in (p1.e)) is [verified] in \texttt{Mathlib.Algebra.CharP.Frobenius}. The chart-side \(\ker D = B^p\) Cartier-direction at the heart of (p1.a)--(p1.d) is [gap] in Mathlib snapshot \texttt{b80f227} for the chart-finite finitely-generated case (a standard-smooth chart of relative dimension \(n\)); the project's iter-145+ in-tree build derives it from the explicit chart presentation of \cref{lem:chart_algebra_isPushout_of_affine_product} (or its companion \texttt{Algebra.IsStandardSmoothOfRelativeDimension} structure) by hand. Mathlib's existing \texttt{Algebra.IsSeparable.iff\_isPurelyInseparable\_extension} provides the \(\ker D = B^p\) direction only modulo separability, which is too restrictive here.

  % NOTE iter-151 render-fix: tag `07F4` (Lemma 16.8.3, smoothing ring
  % maps) replaced with `00T7` (Lemma 10.137.6, standard-smooth ⇒
  % Ω free with explicit basis); verbatim source quoted from the bundled
  % `references/stacks-algebra.tex` (the dangling `stacks-07F4.md`
  % pointer cited by the prior render-fix note never existed on disk).
  % SOURCE: [Stacks Project], Tag 00T7 (Lemma 10.137.6, "Standard smooth
  % algebras"), part (2) (read from references/stacks-algebra.tex, L37258).
  % SOURCE QUOTE: "Let $S = R[x_1, \ldots, x_n]/(f_1, \ldots, f_c)
  % = R[x_1, \ldots, x_n]/I$ be a standard smooth algebra. Then
  % \begin{enumerate} \item the ring map \(R \to S\) is smooth, \item the
  % \(S\)-module \(\Omega_{S/R}\) is free on $\text{d}x_{c + 1}, \ldots,
  % \text{d}x_n$, ..."
  \textit{Source: Stacks Tag 00T7 (Lemma 10.137.6, part (2)).}
  \emph{Literature.} Stacks Tag~\textbf{00T7} (Lemma 10.137.6, the standard-smooth presentation \(\Rightarrow\) \(\Omega_{B/k}\) free with basis \(\{\mathrm dx_{c+1}, \ldots, \mathrm dx_n\}\), the structural fact backing the chart-level \(\ker D = B^p\) argument in characteristic \(p\)); Eisenbud, \emph{Commutative Algebra with a View Toward Algebraic Geometry}, \S 16 (``K\"ahler differentials and derivation kernels''); Mathlib's \texttt{Differential.ContainConstants} class (\texttt{Mathlib.RingTheory.Derivation.DifferentialRing} L62--70), the (p2) char-\(0\) closer, with the project-internal (BR.1)--(BR.5) bridge from the universal K\"ahler derivation \(D\) to a family of coordinate \texttt{Differential B} derivations \(\partial_i\) as the substantive (p2) work (the \texttt{Differential.ContainConstants} class is parameterised by an abstract derivation \(B \to B\), not by the universal K\"ahler derivation \(B \to \Omega_{B/k}\), so the bridge is project-internal infrastructure, not Mathlib infrastructure).
\end{proof}

\begin{theorem}\leanok
[(ii-scheme-lift) Scheme-level \(\mathrm df = 0\) implies \(f\) factors through \(\Spec k\)]
  \label{lem:Scheme_Over_ext_of_diff_zero}
  \lean{AlgebraicGeometry.Scheme.Over.ext_of_diff_zero}
  \uses{thm:GrpObj_eq_of_eqOnOpen, lem:chart_algebra_df_zero_factors_through_constant_on_chart}
  
  % NOTE (iter-146 review): the iter-146 Lean closure is a thin
  % renaming of the iter-125 `Scheme.Over.ext_of_eqOnOpen` packaging in
  % `AlgebraicJacobian/Rigidity.lean` — it consumes `eqOnOpen U`
  % directly rather than deriving it from `df = dg` via the chart-algebra
  % (β) chain (Steps 1–2 of the proof below). The `df = dg` hypothesis,
  % the `genus 0` + `smooth proper` hypotheses on `C`, and the
  % `smooth proper geometrically irreducible group scheme` hypothesis
  % on `A` are all DROPPED from the iter-146 Lean signature in favour
  % of the lighter `[IsSeparated A.hom]` + `[IsReduced C.left]` +
  % `[GeometricallyIrreducible C.hom]` triple that the iter-125
  % packaging already takes. Iter-147+ refines the Lean signature
  % to *also* carry `df = dg`, and rewrites the body to derive
  % `eqOnOpen` from it via Steps 1–2 of the recipe (which depends on
  % the deferred (β-core) sub-piece
  % `lem:chart_algebra_df_zero_factors_through_constant_on_chart`).
  % NOTE (iter-154 review): the (β-core) sub-piece
  % `lem:chart_algebra_df_zero_factors_through_constant_on_chart`
  % is a sorry-free delegate to KDM, and KDM
  % (`lem:KaehlerDifferential_mem_range_algebraMap_of_D_eq_zero`) was
  % CLOSED axiom-clean at iter-154 — the forward direction
  % (`ker D ⊆ range algebraMap`) is fully proved, no longer a
  % structured `sorry`. The iter-147 NOTE here previously gated the
  % substantive refinement of this (lift) lemma on KDM body closure;
  % that gate is now lifted. The iter-148+ refinement to carry
  % `df = dg` substantively is unblocked at the KDM level (the
  % disposition above — a thin renaming of the iter-125 packaging —
  % is otherwise unchanged).
  Let \(f, g \colon C \to A\) be two morphisms in \(\Over\,(\Spec k)\) between a smooth proper geometrically irreducible curve \(C\) of genus \(0\) and a smooth proper geometrically irreducible group scheme \(A\) over \(k\), satisfying:
  \begin{itemize}
    \item \(\mathrm df = \mathrm dg\) as morphisms of \(\mathcal O_C\)-modules \(f^{*}\,\Omega_{A/k} \to \Omega_{C/k}\) (equivalently, the difference \(\mathrm d(\mu \circ \langle f, \iota \circ g \rangle) = \mathrm df - \mathrm dg = 0\) vanishes, where \(\mu\) and \(\iota\) are the multiplication and inversion of \(A\) supplied by the \(\GrpObj A\) structure);
    \item \(f\) and \(g\) agree on a non-empty open subset \(U \subseteq C\), i.e.\ \(U.\iota \circ f|_{U} = U.\iota \circ g|_{U}\) as scheme morphisms \(U \to A\).
  \end{itemize}
  Then \(f = g\) as morphisms in \(\Over\,(\Spec k)\).

  In the specialisation \(g := c\) the constant morphism at \(\eta_A \in A(k)\) (with \(\mathrm dc = 0\) automatically because \(c\) factors through \(\Spec k\), and the agreement on \(U\) supplied by promoting the pointing hypothesis \(f(p) = \eta_A\) to a chart around \(p\) via continuity), this delivers the C.2.d conclusion of \cref{thm:rigidity_over_kbar}: \(f\) equals the constant morphism at \(\eta_A\).

  \paragraph{Statement-vs-Lean signature divergence (iter-155).} The \textbf{present} Lean declaration \texttt{AlgebraicGeometry.Scheme.Over.ext\_of\_diff\_zero} is a \emph{thin wrapper} of the iter-125 packaging \texttt{Scheme.Over.ext\_of\_eqOnOpen}: its signature carries only \texttt{[IsSeparated A.hom]} \(+\) \texttt{[IsReduced C.left]} \(+\) \texttt{[GeometricallyIrreducible C.hom]} and consumes the agreement datum \texttt{eqOnOpen}\,\(U\) directly --- it has \textbf{no} \(\mathrm df = \mathrm dg\) hypothesis, no genus-\(0\) / smooth-proper hypotheses on \(C\), and no group-scheme hypotheses on \(A\). The substantive \(\mathrm df = \mathrm dg\) form stated above (with the genus-\(0\) and \(\GrpObj\) hypotheses, closed by Steps 1--3 below) is a \textbf{gated refinement}: it is inert until \(\mathrm df = 0\) is producible, which is itself gated on gaps (i)\(+\)(ii) [route (a)] or on the dual-AV keystone [route (b)] (see the chapter-introduction disposition). Steps 1--3 below are retained as the eventual closing-glue recipe for that refinement, but they are \textbf{not} a present Lean proof obligation: the wrapper is the live decl.
\end{theorem}

\begin{proof}
  
  \leanok
  \emph{Gating note (iter-155).} Steps 1--3 below are the eventual closing-glue recipe for the gated \(\mathrm df = \mathrm dg\) refinement of the statement above; they are inert until \(\mathrm df = 0\) is producible (gaps (i)\(+\)(ii) / route (b)). The \emph{present} Lean decl is the thin \texttt{ext\_of\_eqOnOpen} wrapper and does not run this recipe. The recipe packages the chart-algebra (\(\beta\)) converse chain into the iter-125-refactored \texttt{Scheme.Over.ext\_of\_eqOnOpen} idiom. Three steps:

  \emph{Step 1: reduce \(\mathrm df = \mathrm dg\) to a single difference morphism \(h\) with \(\mathrm dh = 0\).} Using the group structure of \(A\), define \(h := \mu \circ \langle f, \iota \circ g \rangle \colon C \to A\). The K\"ahler-derivation additivity (Mathlib's \texttt{KaehlerDifferential.D\_sub} / \texttt{Derivation.map\_sub} on the algebra side, lifted to the scheme level by the standard functorial transport of differentials) gives \(\mathrm dh = \mathrm df - \mathrm dg = 0\).
  % iter-155: the prior Step-1 prose cited
  % `\cref{lem:GrpObj_mulRight_globalises}`
  % (`AlgebraicGeometry.GrpObj.mulRight_globalises_cotangent`) for the
  % scheme-level lift. That declaration was EXCISED from the Lean tree at
  % iter-145 (it survives only as a design template, not a present decl);
  % the load-bearing dependency on it is removed here. The Kähler-additivity
  % transport used in Step 1 needs no group-cotangent globalisation. The agreement-on-\(U\) hypothesis transports to \(h|_{U} = \eta_A \circ \pi_C|_{U}\) (i.e.\ on \(U\), \(h\) is the constant morphism at \(\eta_A\)), because \(f|_U = g|_U\) implies \(\mu(f, \iota \circ g)|_U = \mu(g, \iota \circ g)|_U = \eta_A \circ \pi_C|_U\) by the group identity \(b \cdot b^{-1} = \eta_A\) (Mathlib's \(\GrpObj\) axiom \texttt{CategoryTheory.GrpObj.mul\_inv\_self} or the equivalent \texttt{lift\_comp\_inv\_*} pattern).

  \emph{Step 2: apply the per-chart helper \cref{lem:chart_algebra_df_zero_factors_through_constant_on_chart} chart-by-chart on \(C\).} For each affine chart \(V = \Spec R \subseteq C.\mathrm{left}\) on which the restriction of \(h\) lands inside a chart \(W = \Spec B\) of \(A.\mathrm{left}\) around the image \(h(V) \subseteq W\) (such chart pairs exist by quasi-compactness and the standard-smooth presentation of \(A.\mathrm{hom}\); chart-uniform existence is automatic from the smoothness of \(A.\mathrm{hom}\) on each fibre), the chart-restriction of \(h\) is a \(k\)-algebra map \(h^{\sharp} \colon B \to R\) with \(\mathrm dB(h^{\sharp}(b)) = 0\) globally (\cref{lem:chart_algebra_df_zero_factors_through_constant_on_chart} consumes the global Step 3 of its proof here). The helper concludes that \(h^{\sharp}(b) \in \operatorname{range}(\mathtt{algebraMap}\,k\,R)\) for every \(b \in B\), i.e.\ \(h^{\sharp}\) factors as \(B \to k \to R\). Assembling chart-by-chart via sheaf compatibility, the scheme-morphism \(h \colon C \to A\) factors through \(\Spec k \to A\) globally: \(h = c' \circ \pi_C\) for a unique constant morphism \(c' \colon \Spec k \to A\) determined by the chart-factorisation's image.

  \emph{Step 3: identify the constant value via Step 1's agreement on \(U\).} From Step 1's agreement \(h|_{U} = \eta_A \circ \pi_C|_{U}\) and Step 2's factorisation \(h = c' \circ \pi_C\), the constant value satisfies \(c' \circ \pi_C|_{U} = \eta_A \circ \pi_C|_{U}\), hence \(c' = \eta_A\) (the structure map \(\pi_C|_{U} \colon U \to \Spec k\) is non-empty since \(U\) is non-empty, so the equation on \(U\) pins \(c'\)). Therefore \(h = \eta_A \circ \pi_C\) as a morphism \(C \to A\), i.e.\ \(h\) is the constant morphism at \(\eta_A\). Inverting Step 1's construction \(h = \mu \circ \langle f, \iota \circ g \rangle\) gives the cancellation \(\mu(f, \iota \circ g) = \eta_A \circ \pi_C \Rightarrow f = g\) via the \(\GrpObj\) cancellation identity (Mathlib's \texttt{CategoryTheory.GrpObj.mul\_inv\_eq\_iff} or the equivalent \texttt{lift\_comp\_inv\_*} chain on \(\GrpObj A\)).

  Finally, the conclusion is packaged via the iter-125-refactored \cref{thm:GrpObj_eq_of_eqOnOpen} (project lemma \texttt{AlgebraicGeometry.Scheme.Over.ext\_of\_eqOnOpen}, \texttt{AlgebraicJacobian/Rigidity.lean}): on a non-empty open \(U \subseteq C\) where \(f|_{U} = g|_{U}\) (supplied by the agreement hypothesis, or by promoting the pointing \(f(p) = g(p)\) to a chart around \(p\) via continuity), since \(C\) is reduced (smooth) and \(A\) is separated (proper), the scheme-morphism equality \(f = g\) globally follows.

  Estimated build: \(\sim 100\)--\(150\) LOC. The chart-cover assembly of Step 2 is the dominant cost (\(\sim 60\)--\(90\) LOC); Steps 1 and 3 are \(\sim 20\)--\(30\) LOC each.

  \emph{Mathlib status.} \texttt{KaehlerDifferential.D\_sub} / \texttt{Derivation.map\_sub} are [verified] in \texttt{Mathlib.RingTheory.Kaehler.Basic} / \texttt{Mathlib.RingTheory.Derivation.Basic}. The \(\GrpObj\) cancellation identity \texttt{CategoryTheory.GrpObj.mul\_inv\_eq\_iff} (and the underlying \texttt{lift\_comp\_inv\_*} axioms of \texttt{CategoryTheory.GrpObj}) are [verified] in \texttt{Mathlib.CategoryTheory.Monoidal.Grp\_}. \cref{thm:GrpObj_eq_of_eqOnOpen} is [verified] as project lemma \texttt{AlgebraicGeometry.Scheme.Over.ext\_of\_eqOnOpen} in \texttt{AlgebraicJacobian/Rigidity.lean}. The chart-by-chart sheaf-assembly of Step 2 reuses the standard pattern already exercised by the project's \texttt{Scheme.Over.ext\_of\_eqOnOpen} body; the chart-factorisation data is supplied directly by \cref{lem:chart_algebra_df_zero_factors_through_constant_on_chart}.

  \emph{Cheap final glue for the gated \(\mathrm df = 0\) production (iter-155).} On the \emph{production} side --- once gaps (i)\(+\)(ii) are built so that \(f^{*}\Omega_{A/k}\) is globally free with a global basis \(e_1, \ldots, e_g\) (gap i) and each \(\mathrm df(e_i) \in H^0(C, \Omega_{C/k}) = 0\) (gap ii) --- the conclusion \(\mathrm df = 0\) is the trivial linear-algebra glue ``a linear map out of a module spanned by a set is zero iff it kills the spanning set''. Mathlib provides this as \texttt{Submodule.linearMap\_eq\_zero\_iff\_of\_span\_eq\_top} ([verified], \texttt{Mathlib.LinearAlgebra.Span.Basic}), applied componentwise to the global-section module (cf.\ \texttt{analogies/df-zero-production-iter155.md}). This glue is \(\sim 5\)--\(20\) LOC and is \emph{not} the bottleneck: it presupposes both hard gaps already closed.
\end{proof}

\textbf{Honest pile cost (pieces (i)+(ii)+(iii), piece (iv) deferred): 1350--2600 LOC over 7--14 iters.} This is the multi-iter shared build whose iter-range is the dominant cost driver of M2 closure per STRATEGY.md \S Sequencing. The iter-126 analogist's honest decomposition supersedes an earlier "800--1500 LOC / 8--12 iter" framing that conflated piece costs and undercounted piece (i).

Per-piece build order: \textbf{piece (i) \(\to\) piece (ii) \(\to\) piece (iii)}. Piece (i) is the dominant cost and the source of the cotangent-trivial input; pieces (ii) and (iii) consume it. \emph{Iter-155 correction:} this LOC accounting and build order describe only the \textbf{converse} ``\(\mathrm df = 0 \Rightarrow\) constant'' assembly; they do \textbf{not} close \cref{thm:rigidity_over_kbar}, because they do not produce \(\mathrm df = 0\). Producing \(\mathrm df = 0\) additionally requires the \textbf{excised} gap-(i) globalisation (\(\sim 800\)--\(1500\) LOC, \texttt{mulRight\_globalises\_cotangent}, removed iter-145) and piece (iv) Serre duality \(H^0(C, \Omega_{C/k}) = 0\) (\(\sim 3000\)--\(8000\) LOC) --- under closure route (a) --- or the dual-AV \(\Pic^0\) keystone under route (b). Piece (iv) is therefore \textbf{on} the keystone's critical path (route (a)), not merely parked; it remains a deferred named gap, not opened until a closure route is committed.

\section{Use in the project}

The named declaration \cref{thm:rigidity_over_kbar} is consumed by:

\begin{itemize}
  \item \textbf{M2.b genus-0 witness} (iter-127 scaffold \cref{def:genusZeroWitness}, then later body iter-145+). The \texttt{isAlbaneseFor} field of \texttt{JacobianWitness C} for a genus-\(0\) curve \(C\) over \(k\) branches on whether \(C(k) \neq \emptyset\). When \(C(k) \neq \emptyset\), the marked point \(P \in C(k)\) supplied by the universal quantifier feeds into the pointing hypothesis. Under the iter-152 alg-closed pivot, \cref{thm:rigidity_over_kbar} is established over the algebraically closed \(\bar k\) (the constancy of \(f \colon C \to A\) is proved after base change to \(\bar k\), where the curve-side \(\mathrm df = 0 \Rightarrow\) constancy chain is valid); the constancy conclusion then \textbf{descends} to a morphism equality over the general base field \(k\) as a short downstream step, using that the faithfully-flat cover \(\Spec \bar k \to \Spec k\) is an epimorphism of schemes (\texttt{AlgebraicGeometry.Flat.epi\_of\_flat\_of\_surjective}), so morphism equality over \(\bar k\) pulls back to morphism equality over \(k\). The detailed descent lives in the genus-\(0\) witness chapter. When \(C(k) = \emptyset\) (Brauer--Severi conics over \(\mathbb Q\) are the standard counterexample), the field is \textbf{vacuously true} at the Lean type level: Lean's \(\forall (P : \mathbf{1} \to C), \mathtt{IsAlbanese}\ldots\) is trivially satisfied when the type \(\mathbf{1} \to C\) is empty.
  \item \textbf{No other in-tree consumer.} The declaration is project-internal; it is not consumed by any \texttt{Cohomology/*.lean}, \texttt{Genus.lean}, \texttt{AbelJacobi.lean}, or \texttt{Differentials.lean} declaration.
\end{itemize}

\section{Mathlib status}

The named declaration is project-internal scaffolding for iter-126. The shared cotangent-vanishing pile (\cref{sec:RigidityKbar_shared_pile}) is upstream Mathlib gap-fill work that the project will execute in-tree over iter-129+ per the iter-126 user hint endorsing the long-LOC Mathlib gap-fill path. Each of the pile's four pieces (i)--(iv) is a meaningful Mathlib contribution candidate; the project will land them as in-tree material first, with upstream PRs as follow-on work after they are stable.
